% !TeX program = xelatex
% 《投射有限群与 Galois 上同调》中文译本主文件
% 原书: Rosalie Chevalley, Profinite groups and Galois cohomology,
%       Universität Wien, Semester Project, Winter 2009-2010.
%
% 原书编号方案(据原书 PDF 逐节核对):
%   * 节号不带章前缀: "1. Definition and examples"; 小节号 "1.1", "2.3" 等;
%   * 定理类环境按"节"连续编号, 且两组各共用一个计数器:
%       甲组: Definition / Notation / Remark / Example
%       乙组: Theorem / Lemma / Proposition / Corollary
%     如第1章第2节: 命题 2.1-2.9(乙), 定义 2.1、记号 2.2(甲);
%   * 公式全文连续编号 (1)-(7), 不按章重置.
\documentclass[
  BookTitleCN={投射有限群与 Galois 上同调},
  BookTitleEN={Profinite Groups and Galois Cohomology},
  OriginalAuthor={Rosalie Chevalley},
  OriginalEdition={学期项目论文(Semester Project)},
  OriginalPublisher={维也纳大学理学院},
  OriginalYear={2009--2010},
  Translator={Fingolfin},
  ModelUsed={Claude},
  TranslationDate={1 Oct 2026}
]{mathtranslation}

\addbibresource{references.bib}

% 有序列表样式控制: 原书定理类环境内的列表一律用 (1) (2) (3) 编号,
% 故各 enumerate 均写作 \begin{enumerate}[label=(\arabic*),leftmargin=*].
\usepackage{enumitem}

\makeatletter
% ---- 与原书一致的定理编号机制 ----
% cls 已按 [section] 定义 definition/lemma/theorem/corollary/proposition/
% remark/example/problem. 现将两组环境并入各自共享的计数器.
\theoremstyle{mtfang}
\newtheorem{notation}[definition]{记号}
\newtheorem{examples}[definition]{例}
\MT@defautorefname{notation}{记号}
\MT@defautorefname{examples}{例}
\AtEndEnvironment{examples}{\envenddiamond}
\crefname{examples}{例}{例}
% 甲组: 记号、例已并入 definition; 再把注、单数例并入 definition
\let\c@remark\c@definition
\let\c@example\c@definition
% 乙组: 引理、命题、推论并入 theorem
\let\c@lemma\c@theorem
\let\c@proposition\c@theorem
\let\c@corollary\c@theorem
\makeatother

% 节与小节编号不带章前缀, 与原书一致
\renewcommand{\thesection}{\arabic{section}}
\renewcommand{\thesubsection}{\arabic{section}.\arabic{subsection}}

% 公式全文连续编号, 显示为 (1), (2), ...
\counterwithout{equation}{chapter}
\renewcommand{\theequation}{\arabic{equation}}

\begin{document}

\maketitle

\frontmatter
% !TeX root = main.tex

\phantomsection
\chapter*{译者说明}
\addcontentsline{toc}{chapter}{译者说明}

本书译自 Rosalie Chevalley 的维也纳大学理学院学期项目论文(Semester Project, Winter 2009--2010)
\emph{Profinite groups and Galois cohomology}，指导教师为 Dietrich Burde，
助教为 Wolfgang Moens. 翻译与 LaTeX 重排以原书 PDF 为最高依据，
MinerU 与 PaddleOCR 的提取文本仅用作工作底稿.

排版遵循以下对应关系：

\begin{enumerate}
    \item 章、节、小节的题号与原题一致：节号不带章前缀（如``1. 定义与例子''、``2.3 性质''）；
    \item 定理类环境按节编号，且与原书一样分为两组各自共用计数器：
          定义、记号、注、例共用一个计数器，定理、引理、命题、推论共用另一个计数器.
          因此同一节中可能出现``命题 2.4 与定义 2.3 相邻''的情形，这是原书的编号方式，并非排印错误；
    \item 行间公式沿用原书全文连续的编号 (1)--(7)；
    \item 原书正文中方括号式引文（如 [Rib99]）全部改为可点击的交叉引用，
          文献细节（定理号、章节号）保留在引文注记中.
\end{enumerate}

原书为学位论文，个别处编号与语法存在笔误，译本一律照原样保留，不作擅自改动；
涉及数学表述的明显讹误则在脚注中说明. 数学术语的译法参照 Ribes 原著中译本
与国内同调代数、域论论著的习惯译名，首次出现时以楷体标出并附英文原词，
全书收录于文末名词索引.

\vspace{0.5cm}
\hfill \Translator

\hfill \TranslationDate

% !TeX root = main.tex

\phantomsection
\chapter*{摘要}
\addcontentsline{toc}{chapter}{摘要}

本项目的目的是研究群的上同调，尤其是 Galois 上同调. 为此，我们将仔细研究有限与无限的
Galois 群. 这些群与投射有限群(profinite group)大致上是等价的，所谓投射有限群即是有限群列的
逆极限并赋予离散拓扑. 因此，本项目有相当大的一部分内容是研究拓扑群与投射有限群. 我们还会
讨论 Galois 理论的一些主题，特别是无限 Galois 扩张. 我们将在适当的时候用同调代数定义群的
上同调，但主要使用的将是另一个等价的定义，它特别适合处理投射有限群. 最后一章将给出关于群
上同调的一些初等结果，并作为进入 Galois 上同调的入门. 特别地，我们将给出 Hilbert 第 90
定理的几个不同版本.

本项目最常用的参考文献是 L.~Ribes 的一本关于 Galois 上同调导论的著作 \cite{rib99}.


\makecontents

\mainmatter
% !TeX root = ../main.tex
% 第 1 章 拓扑群 （原书 CHAPTER 1, Topological groups, 印刷页 5-11）

\chapter{拓扑群}
\label{ch:1}

本章引入拓扑群的概念及其基本性质. 在无限扩张的情形下推广 Galois 理论基本定理时，这些概念是重要的. 读者可在 \cite{pon66} 与 \cite{hus66} 中进一步了解相关内容.

\section{定义与例子}

\begin{definition}\label{c1-def-1-1}
若元素集合 $G$ 满足
\begin{enumerate}[label=(\arabic*),leftmargin=*]
    \item $G$ 是一个群；
    \item $G$ 是一个拓扑空间；
    \item 映射 $G\times G\to G:(a,b)\mapsto ab$ 连续，其中 $G\times G$ 赋予乘积拓扑；
    \item 映射 $g\to G:x\mapsto x^{-1}$ 连续\footnote{原文即如此（作 $g\to G$）；按上下文应理解为 $G\to G$.}.
\end{enumerate}
则称 $G$ 为 \newterm{topological-group}{拓扑群}{topological group}.
\end{definition}

条件 (3) 与 (4) 分别展开叙述如下：
\begin{itemize}
    \item 若 $a,b\in G$，则对元素 $ab$ 的任一邻域 $W$，存在元素 $a$ 与 $b$ 的邻域 $U$ 与 $V$，使得 $UV\subset W$.
    \item 若 $a\in G$，则对元素 $a^{-1}$ 的任一邻域 $V$，存在元素 $a$ 的邻域 $U$，使得 $U^{-1}\subset V$.
\end{itemize}
最后，这两个条件可以合并表述为：
\begin{itemize}
    \item 若 $a,b\in G$，则对元素 $ab^{-1}$ 的任一邻域 $W$，存在元素 $a$ 与 $b$ 的邻域 $U$ 与 $V$，使得 $UV^{-1}\subset W$.
\end{itemize}

\begin{notation}\label{c1-not-1-2}
在本项目中，$G$ 的单位元一律记作 $1$.
\end{notation}

下面给出一些拓扑群的例子.

\begin{examples}\label{c1-ex-1-3}
以下是一些拓扑群的例子.
\begin{itemize}
    \item 若群 $G$ 赋予 \newterm{discrete-topology}{离散拓扑}{discrete topology}(即 $G$ 的每个子集都是开集)，则 $G$ 是拓扑群.
    \item 考虑带加法的实数群 $(\mathbb{R},+)$. 在 $\mathbb{R}$ 上赋予通常拓扑，它构成拓扑群.
    \item 更一般地，任何拓扑向量空间（例如 Banach 空间或 Hilbert 空间）的加法群都是拓扑群.
    \item 全体可逆 $n\times n$ 实矩阵构成的线性群 $GL_{n}(\mathbb{R})$，若赋予把 $GL_{n}(\mathbb{R})$ 视为 $\mathbb{R}^{n\times n}$ 的子集所得到的拓扑，则它是拓扑群.
\end{itemize}
\end{examples}

\begin{definition}\label{c1-def-1-4}
拓扑群的 \newterm{homomorphism-topgrp}{同态}{homomorphism} 是指连续的群同态. 拓扑群的 \newterm{isomorphism-topgrp}{同构}{isomorphism} 是指既是群同构又是同胚的映射.
\end{definition}

\section{初等性质}

拓扑群是均匀的，这意味着后文某些性质只需在单位元附近验证即可.

\begin{proposition}\label{c1-prop-2-1}
设 $G$ 是拓扑群，$a\in G$，映射 $f,g,\varphi:G\to G$ 定义为 $f(x)=xa$，$g(x)=ax$，$\varphi(x)=x^{-1}$，则 $f$、$g$ 与 $\varphi$ 都是 $G$ 的 \newterm{homeomorphism}{同胚}{homeomorphism}.
\end{proposition}

\begin{proof}
这些映射显然是双射；又由拓扑群的公理，它们是连续的.
\end{proof}

\begin{proposition}\label{c1-prop-2-2}
每个拓扑群都是均匀的，即对任意 $p,q\in G$，存在同胚 $f$ 使得 $f(p)=q$.
\end{proposition}

\begin{proof}
设 $p,q\in G$，令 $f:x\mapsto qp^{-1}x$. 由\autoref{c1-prop-2-1}，它是同胚.
\end{proof}

\begin{proposition}\label{c1-prop-2-3}
设 $V$ 是 $a$ 的一个邻域，则存在 $1$ 的邻域 $U$ 使得 $V=aU$.
\end{proposition}

\begin{proof}
定义 $U:=a^{-1}V$. 作为 $V$ 在连续映射 $g:G\to G:x\mapsto ax$ （见\autoref{c1-prop-2-1}）下的原像，$U$ 是 $1$ 的邻域. 而我们有 $V=aU$，结论得证.
\end{proof}

\begin{definition}\label{c1-def-2-1}
设 $G$ 是拓扑群，称邻域 $V$ 是 \newterm{symmetric-neighborhood}{对称的}{symmetric}，若 $V=V^{-1}$.
\end{definition}

\begin{proposition}\label{c1-prop-2-4}
设 $G$ 是拓扑群. 则 $G$ 存在由对称开集构成的拓扑基. 特别地，$G$ 的每个点处都有由对称邻域构成的\newterm{fundamental-system}{邻域基本系}{fundamental system of neighborhoods}.
\end{proposition}

\begin{proof}
设 $U$ 是 $G$ 的开子集，则 $U^{-1}$ 也是开集(它是 $U$ 在某连续映射下的原像)，于是 $V:=U\cap U^{-1}$ 也是开集. 显然有 $V\subseteq U$ 且 $V=V^{-1}$.
\end{proof}

\begin{proposition}\label{c1-prop-2-5}
设 $G$ 是拓扑群，$U$ 是 $x\in G$ 的一个邻域. 则：
\begin{enumerate}[label=(\arabic*),leftmargin=*]
    \item 存在 $x$ 的邻域 $V$，使得 $V^{2}\subseteq U$.
    \item 存在 $x$ 的邻域 $V$，使得 $V^{-1}\subseteq U$.
    \item 存在 $x$ 的邻域 $V$，使得 $VV^{-1}\subseteq U$.
    \item 存在 $x$ 的邻域 $V$，使得 $V^{-1}V\subseteq U$.
\end{enumerate}
\end{proposition}

\begin{proof}
由拓扑群的均匀性，只需对 $1$ 的邻域证明本命题.

上述四条断言都是下列运算连续性的推论：
\begin{enumerate}[label=(\arabic*),leftmargin=*]
    \item 考虑映射 $f_{1}:G\to G:x\mapsto x^{2}$，它连续（因为乘法连续），并令 $V=f_{1}^{-1}(U)$. 由 $1=1^{2}$，$V$ 是 $1$ 的邻域；并且 $V^{2}=f_{1}(V)=f_{1}(f_{1}^{-1}(U))\subseteq U$.
    \item 同理，利用 $f_{2}:G\to G:x\mapsto x^{-1}$ 与 $1=1^{-1}$.
    \item 同理，利用 $f_{3}:G\to G:x\mapsto xx^{-1}$ 与 $1=11^{-1}$.
    \item 同理，利用 $f_{4}:G\to G:x\mapsto x^{-1}x$ 与 $1=1^{-1}1$.
\end{enumerate}
\end{proof}

\begin{notation}\label{c1-not-2-2}
设 $X$ 是拓扑空间，$A\subseteq X$. 以 $\overline{A}$ 记 $A$ 的拓扑闭包.
\end{notation}

\begin{lemma}\label{c1-lem-2-6}
设 $X,Y$ 是两个拓扑空间，$f:X\to Y$ 是连续映射. 则 $f(\overline{A})\subseteq \overline{f(A)}$.
\end{lemma}

\begin{proof}
设 $y\in f(\overline{A})$，$W$ 是 $Y$ 中 $y$ 的一个邻域. 存在 $x\in\overline{A}$ 使得 $f(x)=y$. 由于 $f$ 连续，$U:=f^{-1}(W)$ 是 $x$ 的邻域. 于是 $U\cap A\neq\emptyset$ (因为 $x\in\overline{A}$). 从而
\[
\emptyset\neq f(U\cap A)\subseteq f(U)\cap f(A)\subseteq W\cap \overline{f(A)}.
\]
所以 $y\in\overline{f(A)}$.
\end{proof}

\begin{proposition}\label{c1-prop-2-7}
设 $G$ 是拓扑群，$\mathcal{V}$ 是 $1$ 的开邻域的一个基本系，$D$ 是 $G$ 中的稠密子集. 则 $\mathcal{B}:=\{Ux\mid x\in D,\,U\in\mathcal{V}\}$ 是 $G$ 的拓扑的一个基.
\end{proposition}

\begin{proof}
设 $W$ 是 $G$ 的开子集，$a\in W$. 则 $Wa^{-1}$ 是开集(它是 $W$ 在某连续映射下的原像)，且 $1\in Wa^{-1}$. 于是存在 $U\in\mathcal{V}$ 使得 $UU^{-1}\subseteq Wa^{-1}$ (利用 $\mathcal{V}$ 是基本系以及\autoref{c1-prop-2-5}).

由于 $D$ 稠密，$aD^{-1}$ 也稠密. 事实上，考虑连续（且双射的）映射 $f:G\to G:x\mapsto ax^{-1}$，有 $X=f(X)=f(\overline{D})\subseteq\overline{f(D)}$\footnote{原文即如此作 $X$；按上下文应为 $G$.} （见\autoref{c1-lem-2-6}）. 于是存在 $d\in U\cap aD^{-1}$，从而存在 $m\in D$ 使得 $d=am^{-1}$，这就推出 $m=d^{-1}a$. 从而 $Ud^{-1}a\in\mathcal{B}$.

由于 $d\in U$，我们得到 $Ud^{-1}a\subseteq UU^{-1}a\subseteq Wa^{-1}a=W$；这对一切 $a\in W$ 成立. 于是
\[
\bigcup_{a\in W}Ud^{-1}a\subseteq W.
\]
此外，$1\in Ud^{-1}$ (因为 $d\in U$)，故对一切 $a\in W$ 有 $a\in Ud^{-1}a$. 于是
\[
W\subseteq\bigcup_{a\in W}Ud^{-1}a.
\]
所以 $W=\bigcup_{a\in W}Ud^{-1}a$.
\end{proof}

\begin{proposition}\label{c1-prop-2-8}
设 $G$ 是拓扑群，$\mathcal{V}$ 是 $1$ 的开邻域的一个基本系. 则 $\bigcap_{V\in\mathcal{V}}V=\overline{\{1\}}$.
\end{proposition}

\begin{proof}
证明双向包含：

$\subseteq$：设 $x\in\bigcap_{V\in\mathcal{V}}V$，$W$ 是 $x$ 的一个邻域. 结合\autoref{c1-prop-2-3} 与\autoref{c1-prop-2-4}，知存在对称集 $V\in\mathcal{V}$ 使得 $Vx\subseteq W$. 由 $x$ 的假设，$x\in V=V^{-1}$，从而 $x^{-1}\in V$. 于是 $1=x^{-1}x\in Vx$. 故 $1\in U$，即 $U\cap\overline{\{1\}}\neq\emptyset$\footnote{原文此处即如此作 $U$；按上下文，此邻域应为上文所设的 $W$.}. 从而 $x\in\overline{\{1\}}$.

$\supseteq$：设 $x\in\overline{\{1\}}$，$V\in\mathcal{V}$ 对称(由\autoref{c1-prop-2-4}，可设 $V$ 对称). 由于 $xV$ 是 $x$ 的邻域，$xV\cap 1\neq\emptyset$. 这说明存在 $v\in V$ 使得 $xv=1$，于是 $x=v^{-1}\in V^{-1}=V$. 故 $x\in\bigcap_{V\in\mathcal{V}}V$.
\end{proof}

\begin{proposition}\label{c1-prop-2-9}
设 $G$ 是拓扑群，$A,B\subseteq G$. 则
\begin{enumerate}[label=(\arabic*),leftmargin=*]
    \item $\overline{A}\,\overline{B}=\overline{AB}$.
    \item $\overline{A}^{-1}=\overline{A^{-1}}$.
    \item $x\overline{A}y=\overline{xAy}$ 对一切 $x,y\in G$ 成立.
\end{enumerate}
\end{proposition}

\begin{proof}
(1) 设 $x\in\overline{A}$，$x\in\overline{B}$\footnote{原文即如此（作 $x\in\overline{B}$）；按上下文应为 $y\in\overline{B}$.}，$U$ 是 $1$ 的邻域. 存在 $1$ 的邻域 $V$ 使得 $xVyV\subseteq xyU$ (因为对一切 $b\in G$，映射 $G\to G:a\mapsto ba$ 连续).

由于 $xV$、$yV$ 分别是 $x$、$y$ 的邻域，存在 $a\in A\cap xV$ 与 $b\in B\cap yV$. 于是 $ab\in AB\cap xyU$，即 $xy\in\overline{AB}$.

(2) 证明双向包含.

$\subseteq$：对连续映射 $i:G\to G:z\mapsto z^{-1}$ 用\autoref{c1-lem-2-6}.

$\supseteq$：$A^{-1}\subseteq\overline{A^{-1}}\Rightarrow A\subseteq(\overline{A^{-1}})^{-1}=i(\overline{A^{-1}})$，而后一个集合是闭集，因为 $i$ 是同胚.

于是 $\overline{A}\subseteq(\overline{A^{-1}})^{-1}\Rightarrow(\overline{A})^{-1}\subseteq\overline{A^{-1}}$.

(3) 对同胚 $G\to G:z\mapsto xzy$ 同理.
\end{proof}

\section{拓扑的刻画}

利用拓扑群的均匀性，我们可以用单位元的邻域基来刻画拓扑.

\begin{proposition}\label{c1-prop-3-1}
设 $G$ 是拓扑群，$\mathcal{V}$ 是 $G$ 的单位元 $1$ 的一个邻域基. 则
\begin{enumerate}[label=(\arabic*),leftmargin=*]
    \item 对一切 $V_{1},V_{2}\in\mathcal{V}$，存在 $V\in\mathcal{V}$ 使得 $1\in V\subseteq V_{1}\cap V_{2}$.
    \item 对一切 $V\in\mathcal{V}$，存在 $W\in\mathcal{V}$ 使得 $W^{2}\subseteq V$.
    \item 对一切 $V\in\mathcal{V}$，存在 $W\in\mathcal{V}$ 使得 $W^{-1}\subseteq V$.
    \item 对一切 $V\in\mathcal{V}$ 与一切 $g\in G$，存在 $W\in\mathcal{V}$ 使得 $W\subseteq gVg^{-1}$.
    \item 对一切 $g\in G$，$\{gV\mid V\in\mathcal{V}\}$ 是 $g$ 的一个邻域基.
\end{enumerate}
反之，若 $G$ 是群，$\mathcal{V}$ 是 $G$ 的一个非空子集族且满足 (1)--(4)，则在 $G$ 上存在（唯一的）拓扑使 (5) 成立.
\end{proposition}

\begin{proof}
设 $\mathcal{V}$ 是拓扑群中 $1$ 的邻域基. 则 (1) 是邻域基定义的推论；(2)、(3)、(4) 是\autoref{c1-prop-2-5} 的推论；而 (5) 成立是因为映射 $f:G\to G:x\mapsto gx$ 是同胚.

反之，设 $G$ 是群，$\mathcal{V}$ 是 $G$ 的一个非空子集族且满足 (1)--(4). 注意 (1) 蕴含每个 $V\in\mathcal{V}$ 都含 $1$. 定义
\[
\mathcal{T}=\{U\subseteq G\mid \forall g\in U,\ \exists V\in\mathcal{V}\ \text{使得}\ gV\subseteq U\},
\]
并证明 $\mathcal{T}$ 是 $G$ 上的拓扑. 首先，空集与 $G$ 显然属于 $\mathcal{T}$. 设 $(U_i)_{i\in I}$ 是 $\mathcal{T}$ 的一族元素($I$ 为指标集)，并设 $g\in\bigcup_{i\in I}U_i$. 于是存在 $j\in I$ 使得 $g\in U_j$，由定义存在 $V\in\mathcal{V}$ 使得 $gV\subseteq U_j$，从而 $gV\subseteq\bigcup_{i\in I}U_i$. 再设 $U_{1},U_{2}\in\mathcal{T}$，$g\in U_{1}\cap U_{2}$. 由定义存在 $V_{1},V_{2}\in\mathcal{V}$ 使得 $gV_{1}\subseteq U_{1}$ 且 $gV_{2}\subseteq U_{2}$. 应用 (1) 即得 $V\in\mathcal{V}$ 使得 $V\subseteq V_{1}\cap V_{2}$，这意味着 $gV\subseteq U_{1}\cap U_{2}$. 故 $U_{1}\cap U_{2}\in\mathcal{T}$. 这样我们证明了 $\mathcal{T}$ 给出 $G$ 上的一个拓扑. 容易验证 (5) 成立.

下面证明 $G\to G:g\mapsto g^{-1}$ 连续. 设 $g\in G$，$U\in\mathcal{T}$ （即为开集）且 $g^{-1}\in U$. 我们要找到 $V\in\mathcal{V}$ 使得 $gV\subseteq U^{-1}$. 由 $\mathcal{T}$ 的定义，存在 $V\in\mathcal{V}$ 使得 $g^{-1}V\subseteq U$，从而 $V^{-1}g\subseteq U^{-1}$. 利用 (3) 可得 $V'\in\mathcal{V}$ 使得 $gV'\subseteq U^{-1}$；再利用 (4)，即知存在 $V''\in\mathcal{V}$ 使得 $gV''\subseteq g(g^{-1}Vg)\subseteq U^{-1}$.

再验证乘法 $G\times G\to G:(g_{1},g_{2})\mapsto g_{1}g_{2}$ 连续. 注意对一切 $V_{1},V_{2}\in\mathcal{V}$，集合 $g_{1}V_{1}\times g_{2}V_{2}$ 构成 $G\times G$ 中 $(g_{1},g_{2})$ 的邻域基. 设 $(g_{1},g_{2})\in G\times G$，$U$ 是满足 $g_{1}g_{2}\in U$ 的开集. 我们要找 $V_{1},V_{2}\in\mathcal{V}$ 使得 $g_{1}V_{1}g_{2}V_{2}\subseteq U$. 由于 $U\in\mathcal{T}$，存在 $V\in\mathcal{V}$ 使得 $g_{1}g_{2}V\subseteq U$. 用 (2) 得 $V'\in\mathcal{V}$ 使得 $g_{1}g_{2}V'V'\subseteq U$. 再用 (4)，得 $V''\in\mathcal{V}$ 使得 $V''\subseteq g_{2}V'g_{2}^{-1}$，从而 $g_{1}V''g_{2}V'\subseteq g_{1}(g_{2}V'g_{2}^{-1})g_{2}V'\subseteq U$. 证明完毕.
\end{proof}

\section{拓扑子群}

本节简要讨论拓扑群的一些性质，特别是其子群的性质.

\begin{proposition}\label{c1-prop-4-1}
设 $G$ 是拓扑群，$H$ 是子群. 则 $H$ 是拓扑群.
\end{proposition}

\begin{proof}
只需证明 $H$ 上的乘法与取逆连续，而这显然是对的：它们只是 $G$ 上相应映射到 $H$ 上的限制.
\end{proof}

\begin{proposition}\label{c1-prop-4-2}
设 $G$ 是拓扑群，$H$ 是子群. 则 $\overline{H}$ 是 $G$ 的子群. 此外，若 $H$ 正规，则 $\overline{H}$ 也正规.
\end{proposition}

\begin{proof}
只需证明 $\overline{H}$ 对乘法与取逆封闭\footnote{原文此处用词为 addition（加法）；按上下文应为乘法.}. 设 $a,b\in\overline{H}$，证明 $ab^{-1}\in\overline{H}$.

设 $W$ 是 $ab^{-1}$ 的邻域，则存在 $a$ 与 $b$ 的邻域 $U$ 与 $V$ 使得 $UV^{-1}\subseteq W$ (因为映射 $G\times G\to G:(x,y)\mapsto xy^{-1}$ 连续). 我们有 $U\cap H\neq\emptyset$ 且 $V\cap H\neq\emptyset$ (注意 $a,b\in\overline{H}$). 于是可取 $x\in U\cap H$ 与 $y\in V\cap H$，使得 $xy^{-1}\in UV^{-1}\cap H\subseteq W\cap H$. 故 $W\cap H\neq\emptyset$，从而 $ab^{-1}\in\overline{H}$.

再设 $H$ 正规. 取 $a\in\overline{H}$ 与 $g\in G$，我们证明 $g^{-1}ag\in\overline{H}$，这便蕴含结论.

设 $V$ 是 $g^{-1}ag$ 的邻域，则存在 $a$ 的邻域 $U$ 使得 $g^{-1}Ug\subseteq V$. 由 $a\in\overline{H}$ 知 $U\cap H\neq\emptyset$. 于是存在 $h\in U\cap H$，使得 $g^{-1}hg\in V$. 又因 $H$ 正规，$g^{-1}hg\in H$. 故 $V\cap H\neq\emptyset$.
\end{proof}

\begin{definition}\label{c1-def-4-1}
称拓扑空间是 \newterm{totally-disconnected}{全不连通的}{totally disconnected}，若它的连通子空间只有单点集.
\end{definition}

\begin{examples}\label{c1-ex-4-2}
\
\begin{itemize}
    \item 一切离散空间都是全不连通的.
    \item 有理数集 $\mathbb{Q}$ 与无理数集 $\mathbb{R}\setminus\mathbb{Q}$ 都是全不连通空间.
    \item 实数集 $\mathbb{R}$（赋予通常拓扑）不是全不连通空间.
\end{itemize}
\end{examples}

\begin{lemma}\label{c1-lem-4-3}
设 $X$ 是全不连通拓扑空间. 则对每个 $x\in X$，$\{x\}$ 在 $X$ 中是闭集.
\end{lemma}

\begin{proof}
设 $C$ 是 $\{x\}$ 的拓扑闭包. 若 $A|B$ 是 $C$ 的一个分离且 $x\in A$，则 $A$ 在 $C$ 中闭，从而在 $X$ 中闭. 于是 $A=C$. 这样 $C$ 是连通的，而 $X$ 全不连通，故必有 $C=\{x\}$.
\end{proof}

\begin{proposition}\label{c1-prop-4-4}
设 $G$ 是拓扑群.
\begin{enumerate}[label=(\arabic*),leftmargin=*]
    \item 若 $H$ 是 $G$ 的开（相应地：闭）子群，则每个陪集 $Hg$ 或 $gH$ 都是开（相应地：闭）集.
    \item $G$ 的每个开子群都是闭的. 若 $G$ 紧致，则 $G$ 的每个开子群的指数有限.
    \item $G$ 的每个指数有限的闭子群都是开的.
    \item 若 $H$ 是包含 $G$ 的某个非空开子集 $U$ 的子群，则 $H$ 在 $G$ 中是开集.
    \item $G$ 是 Hausdorff 空间当且仅当 $\{1\}$ 是 $G$ 的闭集. 且若 $K$ 是 $G$ 的正规子群，则 $G/K$ 是 Hausdorff 空间当且仅当 $K$ 在 $G$ 中闭. 若 $G$ 全不连通，则 $G$ 是 Hausdorff 空间.
\end{enumerate}
\end{proposition}

\begin{proof}
(1) 这是\autoref{c1-prop-2-2} 的推论.

(2) 设 $H$ 开，证明 $\overline{H}\subseteq H$. 取 $a\in\overline{H}$. 由于 $H$ 开，$aH$ 是 $a$ 的开邻域，故 $aH\cap H\neq\emptyset$. 于是存在 $h_{1},h_{2}\in H$ 使得 $ah_{1}=h_{2}$，即 $a=h_{2}h_{1}^{-1}\in H$. 至于 $H$ 指数有限：注意各 $gH$ 是开集（由 (1)）、两两不交且并为 $G$. 故若 $G$ 紧致，$H$ 必指数有限.

(3) 设 $H$ 是 $G$ 的指数有限的闭子群. 则对每个 $g\in G$，$gH$ 闭. 而 $G\setminus H=\bigcup_{g\notin H}gH$，$H$ 指数有限蕴含 $G\setminus H$ 是闭集（闭集的有限并），故 $H$ 是开集.

(4) 每个集合 $hU$ 都是开集（由 (1)），且 $H=\bigcup_{h\in H}hU$，故 $H$ 开.

(5) 由一般拓扑知识，Hausdorff 空间中的单点集都是闭集. 反之设 $\{1\}$ 在 $G$ 中闭. 取 $G$ 中不同元素 $a,b$. 由\autoref{c1-prop-2-2}，集合 $\{ab^{-1}\}$ 是闭集(它是 $\{1\}$ 在同胚下的像). 于是存在开集 $U$ 使得 $1\in U$ 且 $U\subseteq G\setminus\{ab^{-1}\}$. 由映射 $G\to G:(x,y)\mapsto xy^{-1}$ （显然）连续\footnote{原文即如此作 $G\to G$；按上下文应为 $G\times G\to G$.}，$U$ 的原像是开集. 于是存在含 $1$ 的开集 $V,W\subseteq G$，使得 $VW^{-1}\subseteq U$. 这样 $a^{-1}b\notin VW^{-1}$，从而 $aV\cap bW=\emptyset$. 由于 $aV$ 与 $bW$ 都是开集，结论得证.

后面的断言由第一条断言、商拓扑的定义与\autoref{c1-lem-4-3} 推出.
\end{proof}

\section{商群}

这里给出关于拓扑群的商的一个命题.

\begin{proposition}\label{c1-prop-5-1}
设 $G$ 是拓扑群，$N$ 是正规子群. 则 $G/N$ 赋予商拓扑后是拓扑群. 此外，典范投影 $\pi:G\to G/N$ 是连续的开同态.
\end{proposition}

\begin{remark}\label{c1-rem-5-1}
商拓扑定义为\footnote{原文即如此作 $U\subseteq X$；按上下文应为 $U\subseteq G$.}
\[
\mathcal{T}=\{U\subseteq X\mid \pi^{-1}(U)\ \text{在}\ G\ \text{中开}\}.
\]
\end{remark}

\begin{proof}
先证 $\pi$ 是开映射. 设 $U$ 是 $G$ 的开子集，则 $\pi^{-1}(\pi(U))=UN$，而 $UN$ 是开集(因为 $U$ 开)，故结论成立. 此外，由 $\mathcal{T}$ 的定义，$\pi$ 的连续性是平凡的.

下面证明 $G/N$ 是拓扑群. $N$ 正规蕴含 $G/N$ 是抽象群. 验证乘法与取逆连续.

设 $W$ 是 $g_{1}g_{2}N\in G/N$ 的邻域，则因 $\pi$ 连续，$\pi^{-1}(W)$ 是 $g_{1}g_{2}$ 的邻域. 由 $G$ 中乘法的连续性，存在 $g_{1},g_{2}$ 的邻域 $U_{1},U_{2}$，使得 $U_{1}U_{2}\subseteq\pi^{-1}(W)$. 于是 $\pi(U_{1})$ 与 $\pi(U_{2})$ 分别是 $g_{1}N$ 与 $g_{2}N$ 的邻域(利用 $\pi$ 是开映射). 又因 $\pi$ 是同态，$\pi(U_{1})\pi(U_{2})=\pi(U_{1}U_{2})\subseteq W$. 故乘法连续.

设 $W$ 是 $gN\in G/N$ 的邻域，则因 $\pi$ 连续，$\pi^{-1}(W)$ 是 $g$ 的邻域. 由 $G$ 中取逆的连续性，存在 $g$ 的邻域 $U$ 使得 $U^{-1}\subseteq\pi^{-1}(W)$. 于是 $\pi(U)$ 是 $gN$ 的邻域，且 $\pi(U)^{-1}=\pi(U^{-1})\subseteq W$. 故取逆连续.
\end{proof}

\section{Lie 群}

我们简要介绍 Lie 群，因为它们提供大量拓扑群的例子. 关于这一主题的更多信息可参见 ??\footnote{原文此处即为 ``??''（未填妥的引文）.}.

\begin{definition}\label{c1-def-6-1}
\newterm{Lie-group}{Lie 群}{Lie group} 是指赋予群结构的有限维光滑流形 $G$，其乘法光滑. 这就是说，存在光滑乘法 $\mu_{G}\times G\to G$\footnote{原文即如此排印；按上下文应理解为 $\mu:G\times G\to G$.}、取逆 $i:G\to G$ 与单位元 $1\in G$，满足群公理.
\end{definition}

光滑流形的定义可见 \cite[Chapter 2]{lang02b}.

\begin{remark}\label{c1-rem-6-2}
注意 Lie 群都是拓扑群. 它们为拓扑群提供了大量例子.
\end{remark}

\begin{examples}\label{c1-ex-6-3}
\
\begin{itemize}
    \item $\mathbb{R}$ 与 $\mathbb{C}$ 对加法构成 Lie 群. 并且，任何有限维实或复向量空间对加法都构成 Lie 群.
    \item $\mathbb{R}^{*}(=\mathbb{R}\setminus\{0\})$ 与 $\mathbb{C}^{*}$ 对乘法构成 Lie 群. $\mathbb{S}=\{z\in\mathbb{C}\mid |z|=1\}$ 对乘法也构成 Lie 群.
    \item 若 $G$ 与 $H$ 是 Lie 群，则乘积 $G\times H$ 是 Lie 群（赋予通常的乘积结构）. 于是 (1) 与 (2) 蕴含环面 $\mathbb{T}_{n}=\mathbb{S}^{n}$ 是 Lie 群.
    \item Lie 群的基本例子是有限维实向量空间 $V$ 上全体可逆线性映射构成的群 $GL(V)$. 我们对 $V=\mathbb{R}^{n}$ 的情形加以验证. 设 $f\in L(V,V)$，取标准基 $\{e_{1},\ldots,e_{n}\}$. 元素 $f(e_{i})\in\mathbb{R}^{n}$ 就是 $f$ 的矩阵的第 $i$ 列. 这给出 $L(V,V)$ 到 $M_{n}(\mathbb{R})=\mathbb{R}^{n^{2}}$ 的映射. 行列式定义光滑函数 $\det:M_{n}(\mathbb{R})\to\mathbb{R}$. 特别地，$GL_{n}(\mathbb{R})=\det^{-1}(\mathbb{R}_{n}^{*})$ 是 $\mathbb{R}^{2}_{n}$\footnote{原文即如此作 $\mathbb{R}^{2}_{n}$；按上下文应为 $\mathbb{R}^{n^{2}}$.} 的开子集，从而是光滑流形. 此外，两矩阵 $A$ 与 $B$ 之积的元素是 $A$ 与 $B$ 的元素的多项式，故乘法定义光滑映射.
    \item 特殊线性群 $SL_{n}(\mathbb{R})$ 定义为行列式映射 $\det:GL_{n}(\mathbb{R})\to\mathbb{R}^{*}$ 的核，它是 $GL_{n}(\mathbb{R})$ 的子群，故为 Lie 群.
    \item 经典群 $O_{n}(\mathbb{R})$ 与 $SO_{n}(\mathbb{R})$ 也是 Lie 群，其中 $O_{n}(\mathbb{R})=\{A\in GL_{n}(\mathbb{R})\mid AA^{T}=A^{T}A=\mathrm{Id}_{n}\}$，$SO_{n}(\mathbb{R})=\{A\in O_{n}(\mathbb{R})\mid \det(A)=1\}$.
\end{itemize}
\end{examples}

下面给出一些不是 Lie 群的拓扑群的例子.

\begin{examples}\label{c1-ex-6-4}
\
\begin{itemize}
    \item 无限维群，例如无限维实向量空间的加法群. 它们不是有限维流形，故不是 Lie 群.
    \item 某些全不连通群，例如域的无限扩张的 Galois 群，或 $p$ 进数域的加法群. 它们的底空间不是实流形，故不是 Lie 群.
\end{itemize}
\end{examples}

% !TeX root = ../main.tex
% 第 2 章 有限与无限的 Galois 理论 （原书 CHAPTER 2, Finite and infinite Galois Theory, 印刷页 13-22 / PDF 页 15-24）

\chapter{有限与无限的 Galois 理论}
\label{ch:2}

% 原书小节号为 "节.小节"，1.2.1 等三级标题不带章前缀
\renewcommand{\thesubsubsection}{\thesection.\arabic{subsection}.\arabic{subsubsection}}
% 有序列表标签照原书样式 (1)(2)...
\renewcommand{\labelenumi}{(\arabic{enumi})}

本章首先给出经典 Galois 理论的概要，特别是有限扩张情形的 Galois 理论基本定理. 然后，我们为 Galois 群赋予一个特定的拓扑，以便陈述并证明无限扩张情形的基本定理.

\section{Galois 理论回顾}

本节中，$\mathbb{K}$、$\mathbb{L}$ 与 $\mathbb{E}$ 恒表示域.

\subsection{一些预备知识}

\begin{definition}\label{c2-def-1-1}
设 $\mathbb{L}$ 是域 $\mathbb{K}$ 的扩张. 若 $\mathbb{L}$ 的每个元素都在 $\mathbb{K}$ 上代数，则称 $\mathbb{L}$ 在 $\mathbb{K}$ 上是\emph{代数的}.
\end{definition}

\begin{example}\label{c2-ex-1-2}
考虑含 $q$ 个元素的有限域 $\mathbb{F}_q$，其中 $q=p^{k}$，$p$ 为素数. 我们回忆：$\mathbb{F}_q$ 的一切有限扩张的阶都是 $q$ 的幂（见 \cite[Thm~5.1]{lang02a} 与 \cite[Corollary~5.2]{lang02a}）.
\end{example}

\begin{definition}\label{c2-def-1-3}
设 $\mathbb{E}$ 与 $\mathbb{L}$ 是 $\mathbb{K}$ 的两个扩张. 从 $\mathbb{L}$ 到 $\mathbb{E}$ 的 \emph{$\mathbb{K}$-同态}是指满足 $\sigma|_{\mathbb{K}}=\mathrm{Id}_{\mathbb{K}}$ 的同态 $\sigma:\mathbb{L}\to\mathbb{E}$.
\end{definition}

\begin{definition}\label{c2-def-1-4}
设 $G$ 是域 $\mathbb{L}$ 的全部自同态所成群的子群. 我们定义
\[
\mathbb{L}^{G}=\{\alpha\in\mathbb{L}\mid \sigma(\alpha)=\alpha,\ \forall\sigma\in G\}.
\]
这是 $\mathbb{L}$ 的子域，称为 $G$ 的\emph{固定域}.
\end{definition}

\subsection{正规扩张与可分扩张}

\subsubsection{正规扩张}

\begin{theorem}\label{c2-thm-1-1}
设 $\mathbb{L}$ 是包含在 $\mathbb{K}$ 的某个代数闭包 $\overline{\mathbb{K}}$ 中的 $\mathbb{K}$ 的代数扩张. 则以下条件等价：
\begin{enumerate}
    \item $\mathbb{L}$ 是 $\mathbb{K}[X]$ 中一族非常数多项式的\newterm{c2-splitting-field}{分裂域}{splitting field}；
    \item 对每个 $\mathbb{K}$-同态 $\sigma:\mathbb{L}\to\overline{\mathbb{K}}$，都有 $\sigma(\mathbb{L})=\mathbb{L}$；
    \item $\mathbb{K}[X]$ 中每个在 $\mathbb{L}$ 中有根的多项式，其全部根都在 $\mathbb{L}$ 中，即它在 $\mathbb{K}$ 中分解为一次因式.\footnote{原文即如此（作 $\mathbb{K}$）；按上下文应理解为在 $\mathbb{L}$（或 $\overline{\mathbb{K}}$）中分解.}
\end{enumerate}
\end{theorem}

\begin{proof}
该定理的证明可参见 \cite[Chapter~5, \S3]{lang02a}.
\end{proof}

\begin{definition}\label{c2-def-1-5}
若 $\mathbb{K}$ 的代数扩张 $\mathbb{L}$ 满足前述定理中的某一条性质，则称 $\mathbb{L}$ 为 $\mathbb{K}$ 的\newterm{c2-normal-extension}{正规扩张}{normal extension}.
\end{definition}

\begin{examples}\label{c2-ex-1-6}
\
\begin{itemize}
    \item 若 $[\mathbb{L}:\mathbb{K}]=2$，则 $\mathbb{L}$ 在 $\mathbb{K}$ 上正规. 事实上，取 $\alpha\in\mathbb{L}\setminus\mathbb{K}$，$\alpha$ 在 $\mathbb{K}$ 上的极小多项式可写成 $f(X)=X^{2}+cX+d$. 于是 $-c+\alpha\in\mathbb{L}$ 就是 $f$ 的另一个根.
    \item 扩张 $\mathbb{Q}\subset\mathbb{Q}[\sqrt{2}]$ 是正规的. 但扩张 $\mathbb{Q}\subset\mathbb{Q}[\sqrt[3]{2}]$ 不是正规的（$X^{3}-2$ 在 $\mathbb{Q}[\sqrt[3]{2}]$ 中不分解）.
\end{itemize}
\end{examples}

\subsubsection{可分扩张}

\begin{definition}\label{c2-def-1-7}
不可约多项式 $f\in\mathbb{K}[X]$ 称为\emph{可分的}，若它没有重根.
\end{definition}

\begin{definition}\label{c2-def-1-8}
设 $\mathbb{L}$ 是 $\mathbb{K}$ 的扩张.

若 $\alpha\in\mathbb{L}$ 在 $\mathbb{K}$ 上代数且 $\min(\alpha,\mathbb{K})$ 是可分的，则称 $\alpha$ 在 $\mathbb{K}$ 上\emph{可分}.

若每个 $\alpha\in\mathbb{E}$\footnote{原文即如此（作 $\mathbb{E}$）；按上下文应为 $\mathbb{L}$.}都在 $\mathbb{K}$ 上可分，则称 $\mathbb{L}$ 为 $\mathbb{K}$ 的\newterm{c2-separable}{可分扩张}{separable extension}.
\end{definition}

\begin{remark}\label{c2-rem-1-9}
下面是可分扩张的一个显式刻画：

若 $\mathbb{K}[X]$ 中每个在 $\mathbb{L}$ 中有根的不可约多项式都是可分的，则 $\mathbb{L}$ 在 $\mathbb{K}$ 上代数且可分.
\end{remark}

\begin{examples}\label{c2-ex-1-10}
\
\begin{itemize}
    \item 若 $\operatorname{char}(\mathbb{K})=0$，则 $\mathbb{K}$ 的每个代数扩张都可分.
    \item 若 $\mathbb{K}$ 有限，则 $\mathbb{K}$ 的每个代数扩张都可分.
    \item 考虑素数域 $\mathbb{F}_p$ 与 $X^{p}-t$ 的一个根 $\alpha$. 则 $\mathbb{F}_p[\alpha]\cong\mathbb{F}_p[X]/(X^{p}-t)$ 不是 $\mathbb{F}_p$ 的可分扩张，因为 $\alpha$ 不可分（$\alpha$ 在 $\mathbb{K}$\footnote{原文即如此（作 $\mathbb{K}$）；按上下文应为 $\mathbb{F}_p$.}上的极小多项式为 $X^{p}-t=(X-\alpha)^{p}$）.
\end{itemize}
\end{examples}

\begin{definition}\label{c2-def-1-11}
设 $\overline{\mathbb{K}}$ 是 $\mathbb{K}$ 的代数闭包. $\mathbb{K}$ 在 $\overline{\mathbb{K}}$ 中的\emph{可分闭包}是指 $\overline{\mathbb{K}}$ 中包含 $\mathbb{K}$ 的一切有限可分扩张的最小子域.

我们将 $\mathbb{K}$ 在其某个代数闭包 $\overline{\mathbb{K}}$ 中的可分闭包记作 $\mathbb{K}_s$.
\end{definition}

\begin{remark}\label{c2-rem-1-12}
换言之，$\mathbb{K}$ 的可分闭包就是包含在 $\overline{\mathbb{K}}$ 中的 $\mathbb{K}$ 的一切有限可分扩张之并.
\end{remark}

下面介绍可分次数与不可分次数，后文将用到它们.

\begin{definition}\label{c2-def-1-13}
设 $\mathbb{L}$ 是 $\mathbb{K}$ 的有限扩张，$\overline{\mathbb{K}}$ 是 $\mathbb{K}$ 的代数闭包. $\mathbb{L}$ 在 $\mathbb{K}$ 上的\emph{可分次数}记作 $[\mathbb{L}:\mathbb{K}]_s$，是指包含同态 $i:\mathbb{K}\to\overline{\mathbb{K}}$ 到 $\mathbb{L}$ 上的扩张的个数.

我们必须验证这一概念是良定义的，即 $[\mathbb{L}:\mathbb{K}]_s$ 与 $\overline{\mathbb{K}}$ 的选取无关. 这一点在 \cite[\S4, Chapter~V]{lang02a} 中完成.
\end{definition}

\begin{proposition}\label{c2-prop-1-2}
设 $\mathbb{E}$ 是 $\mathbb{K}$ 的有限扩张，$\mathbb{L}$ 是 $\mathbb{E}$ 的有限扩张. 则
\[
[\mathbb{L}:\mathbb{K}]_s=[\mathbb{L}:\mathbb{E}]_s\cdot[\mathbb{E}:\mathbb{K}]_s
\]
\end{proposition}

\begin{proof}
参照 \cite[Theorem~4.1, \S1, Chapter~V]{lang02a} 即得证明.
\end{proof}

\begin{proposition}\label{c2-prop-1-3}
设 $\mathbb{L}$ 是 $\mathbb{K}$ 的有限扩张. 则可分次数 $[\mathbb{L}:\mathbb{K}]_s$ 整除次数 $[\mathbb{L}:\mathbb{K}]$. 此外，若特征为 $0$，则商为 $1$；若特征为素数 $p>0$，则商为 $p$ 的幂.
\end{proposition}

\begin{definition}\label{c2-def-1-14}
若 $\mathbb{L}$ 是 $\mathbb{K}$ 的有限扩张，我们称商
\[
\frac{[\mathbb{L}:\mathbb{K}]}{[\mathbb{L}:\mathbb{K}]_s}
\]
为\emph{不可分次数}，记作 $[\mathbb{L}:\mathbb{K}]_i$.
\end{definition}

\subsection{Galois 扩张}

\begin{definition}\label{c2-def-1-15}
若 $\mathbb{L}$ 是 $\mathbb{K}$ 的代数扩张、正规扩张且可分扩张，则称 $\mathbb{L}$ 为 $\mathbb{K}$ 的\newterm{c2-galois-extension}{Galois 扩张}{Galois extension}.
\end{definition}

\begin{remark}\label{c2-rem-1-16}
\
\begin{itemize}
    \item 我们可以更明确地描述 Galois 扩张：

设 $\mathbb{L}$ 是 $\mathbb{K}$ 的 Galois 扩张，$f$ 是 $\mathbb{K}[X]$ 中一个 $m$ 次不可约多项式. 若 $f$ 在 $\mathbb{L}$ 中有根，则 $f$ 在 $\mathbb{L}$ 中有 $m$ 个不同的根. 于是，$\mathbb{L}$ 是 Galois 扩张当且仅当：对每个 $\alpha\in\mathbb{L}$，$\alpha$ 在 $\mathbb{K}$ 上的极小多项式在 $\mathbb{L}$ 中有 $[\mathbb{K}[\alpha]:\mathbb{K}]$ 个不同的根.
    \item 在 Galois 扩张的定义中，我们不要求扩张的次数有限. 后文若无特别说明，Galois 扩张可以是有限的也可以是无限的.
\end{itemize}
\end{remark}

\begin{example}\label{c2-ex-1-17}
考虑素数域 $\mathbb{F}_p$ 的扩张 $\mathbb{F}_{p^{n}}$. 由构造，$\mathbb{F}_{p^{n}}\cong\mathbb{F}_p/f_n(X)$\footnote{原文即如此；应理解为 $\mathbb{F}_p[X]/(f_n(X))$，且断言"$f_n$ 不可约"仅当 $n=1$ 时成立.}，其中 $f_n(X)=X^{p^{n}}-X$ 是 $\mathbb{F}_p$ 上的不可约多项式（见 \cite[Theorem~5.1, \S5, Chapter~V]{lang02a}）. 于是 $\mathbb{F}_{p^{n}}$ 的每个元素都是多项式 $f_n(X)$ 的根，换言之 $\mathbb{F}_{p^{n}}$ 是多项式 $f_n(X)$ 的分裂域. 故 $\mathbb{F}_{p^{n}}$ 是 $\mathbb{F}_p$ 的 Galois 扩张.
\end{example}

下面这条命题在后文中会经常用到.

\begin{proposition}\label{c2-prop-1-4}
若 $\mathbb{L}$ 是 $\mathbb{K}$ 的 Galois 扩张（有限或无限），则 $\mathbb{L}$ 对任何中间域 $\mathbb{E}$（即 $\mathbb{K}\subseteq\mathbb{E}\subseteq\mathbb{L}$）都是 Galois 扩张.
\end{proposition}

\begin{proof}
设 $f$ 是 $\mathbb{E}[X]$ 中在 $\mathbb{L}$ 中有根 $\alpha$ 的不可约多项式. 由于代数性是传递的，可取 $g=\min(\alpha,\mathbb{K})\in\mathbb{K}[X]$，即 $\alpha$ 在 $\mathbb{K}$ 上的极小多项式. $\mathbb{K}$\footnote{原文两处即如此（作 $\mathbb{K}$）；按上下文均应为 $\mathbb{L}$.}的正规性蕴含 $g$ 在 $\mathbb{L}[X]$ 中分解，$\mathbb{K}$ 的可分性蕴含 $g$ 在 $\mathbb{L}$ 中有不同的根（因为 $g$ 在 $\mathbb{L}$ 中有根 $\alpha$）.

由于 $f$ 整除 $g$（在 $\mathbb{E}[X]$ 中），$f$ 也必在 $\mathbb{L}[X]$ 中分解为互异的一次因式.
\end{proof}

\subsection{Galois 群与 Galois 理论基本定理}

本节将特别给出有限情形的 Galois 理论基本定理.

\begin{definition}\label{c2-def-1-18}
设 $\mathbb{L}$ 是 $\mathbb{K}$ 的代数扩张，有限或无限. $\mathbb{L}$ 在 $\mathbb{K}$ 上的\newterm{c2-galois-group}{Galois 群}{Galois group}记作 $\operatorname{Gal}(\mathbb{L},\mathbb{K})$，是 $\mathbb{L}$ 的全部 $\mathbb{K}$-同态构成的集合.
\end{definition}

\begin{definition}\label{c2-def-1-19}
若 Galois 扩张 $\mathbb{L}$（$\mathbb{K}$ 上）的 Galois 群是阿贝尔群（相应地：循环群），则称该扩张为阿贝尔（相应地：循环）扩张.
\end{definition}

下面的例子给出一个 Galois 扩张，并且它为我们理解 Galois 理论主定理提供直觉.

\begin{example}\label{c2-ex-1-20}
考虑域 $\mathbb{Q}$ 及其扩张 $\mathbb{Q}[\sqrt[3]{2},\omega]$，其中 $\omega=e^{\frac{2i\pi}{3}}=-\dfrac{1}{2}+i\dfrac{\sqrt{3}}{2}$. 我们定义
\[
\begin{aligned}
\sigma:\mathbb{Q}[\sqrt[3]{2},\omega]&\longrightarrow\mathbb{Q}[\sqrt[3]{2},\omega]\\
\sqrt[3]{2}&\longmapsto\omega\sqrt[3]{2}\\
\omega&\longmapsto\omega
\end{aligned}
\qquad\text{和}\qquad
\begin{aligned}
\tau:\mathbb{Q}[\sqrt[3]{2},\omega]&\longrightarrow\mathbb{Q}[\sqrt[3]{2},\omega]\\
\sqrt[3]{2}&\longmapsto\sqrt[3]{2}\\
\omega&\longmapsto\omega^{2}=\overline{\omega}.
\end{aligned}
\]
于是可以证明 $\{\mathrm{Id}_{\mathbb{Q}[\sqrt[3]{2},\omega]},\sigma,\sigma^{2},\tau,\tau\sigma,\tau\sigma^{2}\}=\operatorname{Gal}(\mathbb{Q}[\sqrt[3]{2},\omega],\mathbb{Q})$. 从而 $|\operatorname{Gal}(\mathbb{Q}[\sqrt[3]{2},\omega],\mathbb{Q})|=6=[\mathbb{Q}[\sqrt[3]{2},\omega]:\mathbb{Q}]$. （注意 $\mathbb{Q}[\sqrt[3]{2},\omega]$ 是 $\mathbb{Q}$ 的正规扩张.）
\end{example}

\begin{theorem}\label{c2-thm-1-5}
设 $\mathbb{L}$ 是 $\mathbb{K}$ 的有限 Galois 扩张，令 $G=\operatorname{Gal}(\mathbb{L},\mathbb{K})$. 存在双对应
\[
\{H\mid H\leq G\ \text{子群}\}\leftrightarrow\{\mathbb{E}\mid \mathbb{K}\subseteq\mathbb{E}\subseteq\mathbb{L}\ \text{中间域}\}.
\]
更确切地，映射 $H\mapsto\mathbb{L}^{H}$ 与 $\mathbb{E}\mapsto\operatorname{Gal}(\mathbb{L},\mathbb{E})$ 是 $G$ 的子集全体与 $\mathbb{L}$、$\mathbb{K}$ 之间中间域全体之间的互逆双射. 此外，
\begin{enumerate}
    \item 设 $H_{1},H_{2}\leq G$，则 $H_{1}\supseteq H_{2}$ 当且仅当 $\mathbb{L}^{H_{1}}\subseteq\mathbb{L}^{H_{2}}$；
    \item 指数等于次数：$[H_{1}:H_{2}]=[\mathbb{E}^{H_{2}}:\mathbb{E}^{H_{1}}]$；
    \item 若 $H$ 是 $G$ 的子群且 $\sigma\in G$，则 $\mathbb{L}^{\sigma H\sigma^{-1}}=\sigma(\mathbb{L}^{H})$. 反之，若 $\mathbb{E}$ 是中间域，则 $\operatorname{Gal}(\mathbb{L},\sigma\mathbb{E})=\sigma\operatorname{Gal}(\mathbb{L},\mathbb{E})\sigma^{-1}$；
    \item 设 $\mathbb{E}_{1},\mathbb{E}_{2}$ 是中间域，则
    \[
    \operatorname{Gal}(\mathbb{L},\mathbb{E}_{1}\mathbb{E}_{2})=\operatorname{Gal}(\mathbb{L},\mathbb{E}_{1})\cap\operatorname{Gal}(\mathbb{L},\mathbb{E}_{2}).
    \]
    \item $\mathbb{E}$ 是 $\mathbb{K}$ 的 Galois 扩张当且仅当 $\operatorname{Gal}(\mathbb{L},\mathbb{E})$ 是 $G$ 的正规子群，并且此时有
    \[
    \operatorname{Gal}(\mathbb{E},\mathbb{K})\cong G/\operatorname{Gal}(\mathbb{L},\mathbb{E}).
    \]
\end{enumerate}
\end{theorem}

\begin{proof}
该定理的证明可参见 \cite[Chapter~3]{mil08} 或 \cite[\S1, Chapter~IV]{lang02a}.
\end{proof}

\begin{remark}\label{c2-rem-1-21}
对无限 Galois 扩张，该定理一般不成立. 因此我们将为 Galois 群配备一个拓扑，以获得对有限与无限 Galois 扩张都成立的定理.
\end{remark}

\section{Galois 群上的 Krull 拓扑}

如前所述，Galois 理论主定理对无限 Galois 扩张并不成立. 然而，为 Galois 群赋予拓扑后，我们将得到类似的结果. Krull 拓扑是 Galois 群上最自然的（非平凡）拓扑，且具有许多有趣的性质.

注意，本节对域扩张的次数不作任何限制.

\subsection{预备命题}

本节的目的是在 Galois 群上定义一个拓扑，为此我们要用到一些命题和引理.

\begin{lemma}\label{c2-lem-2-1}
设 $\mathbb{L}$ 是 $\mathbb{K}$ 的 Galois 扩张，$\mathbb{M}$ 是中间域，即 $\mathbb{K}\subseteq\mathbb{M}\subseteq\mathbb{L}$. 则从 $\mathbb{M}$ 到 $\mathbb{L}$ 的每个 $\mathbb{K}$-同态都可扩张为从 $\mathbb{L}$ 到 $\mathbb{L}$ 的 $\mathbb{K}$-自同构.
\end{lemma}

\begin{proof}
见 \cite[Chapter~7]{mil08}.
\end{proof}

\begin{lemma}\label{c2-lem-2-2}
设 $\mathbb{L}$ 是 $\mathbb{K}$ 的 Galois 扩张. 对每个满足 $\mathbb{K}\subseteq\mathbb{E}\subseteq\mathbb{L}$ 的 $\mathbb{K}$ 的有限 Galois 扩张 $\mathbb{E}$，映射
\[
\begin{aligned}
\operatorname{Gal}(\mathbb{L},\mathbb{K})&\longrightarrow\operatorname{Gal}(\mathbb{E},\mathbb{K})\\
\sigma&\longmapsto\sigma|_{\mathbb{E}}
\end{aligned}
\]
是满射.
\end{lemma}

\begin{proof}
设 $\sigma\in\operatorname{Gal}(\mathbb{E},\mathbb{K})$，则 $\sigma$ 是从 $\mathbb{L}$ 到 $\mathbb{E}$\footnote{原文即如此；按上下文应为从 $\mathbb{E}$ 到 $\mathbb{L}$ 的 $\mathbb{K}$-同态.}的 $\mathbb{K}$-同态. 于是由\autoref{c2-lem-2-1}，$\sigma$ 可扩张为从 $\mathbb{L}$ 到 $\mathbb{L}$ 的 $\mathbb{K}$-自同构（它就是原像）.
\end{proof}

\begin{notation}\label{c2-not-2-1}
设 $\mathbb{L}$ 是 $\mathbb{K}$ 的 Galois 扩张. 以 $\mathcal{F}_{\mathbb{L}}$ 记满足 $\mathbb{K}\subseteq\mathbb{E}\subseteq\mathbb{L}$ 的 $\mathbb{K}$ 的一切有限 Galois 扩张 $\mathbb{E}$ 构成的族.
\end{notation}

\begin{remark}\label{c2-rem-2-2}
注意，若 $\mathbb{L}$ 是 $\mathbb{K}$ 的 Galois 扩张（即使无限），也有
\[
\mathbb{L}=\bigcup_{\mathbb{E}\in\mathcal{F}_{\mathbb{L}}}\mathbb{E}.
\]
\end{remark}

\begin{lemma}\label{c2-lem-2-3}
设 $\mathbb{L}$ 是 $\mathbb{K}$ 的 Galois 扩张，$\mathbb{M}$ 是中间域. 则
\[
\mathbb{L}^{\operatorname{Gal}(\mathbb{L},\mathbb{M})}=\mathbb{M}.
\]
\end{lemma}

\begin{proof}
由固定域的定义，显然有 $\mathbb{M}\subseteq\mathbb{L}^{\operatorname{Gal}(\mathbb{L},\mathbb{M})}$.

下面证反包含：

设 $\alpha\in\mathbb{L}^{\operatorname{Gal}(\mathbb{L},\mathbb{M})}$. 则 $\operatorname{Gal}(\mathbb{L},\mathbb{M})$ 的定义蕴含对一切 $\sigma\in\operatorname{Gal}(\mathbb{L},\mathbb{M})$ 有 $\sigma(\alpha)=\alpha$. 由于 $\mathbb{L}$ 是有限中间 Galois 扩张之并，存在这样一个扩张 $\mathbb{E}$ 使得 $\alpha\in\mathbb{E}$. 于是由\autoref{c2-lem-2-2}，映射
\[
\operatorname{Gal}(\mathbb{L},\mathbb{M})\to\operatorname{Gal}(\mathbb{E},\mathbb{M}):\sigma\mapsto\sigma|_{\mathbb{E}}
\]
是满射. 从而对一切 $\sigma\in\operatorname{Gal}(\mathbb{E},\mathbb{M})$ 有 $\sigma(\alpha)=\alpha$. 再利用有限扩张情形的定理\autoref{c2-thm-1-5}，即得 $\alpha\in\mathbb{E}^{\operatorname{Gal}(\mathbb{E},\mathbb{M})}=\mathbb{M}$.
\end{proof}

\begin{proposition}\label{c2-prop-2-4}
设 $\mathbb{L}$ 是 $\mathbb{K}$ 的 Galois 扩张. 若 $\mathbb{E}$ 是 $\mathbb{K}$ 的有限中间 Galois 扩张，则 $\operatorname{Gal}(\mathbb{L},\mathbb{E})$ 在 $\operatorname{Gal}(\mathbb{L},\mathbb{K})$ 中正规，且 $\operatorname{Gal}(\mathbb{L},\mathbb{K})/\operatorname{Gal}(\mathbb{L},\mathbb{E})=\operatorname{Gal}(\mathbb{E},\mathbb{K})$.
\end{proposition}

\begin{proof}
取 $\mathbb{E}'\in\mathcal{F}_{\mathbb{L}}$ 使得 $\mathbb{K}\subseteq\mathbb{E}\subseteq\mathbb{E}'\subseteq\mathbb{L}$. 由\autoref{c2-thm-1-5}，因 $\mathbb{E}'$ 是 $\mathbb{K}$ 的有限扩张，有 $\operatorname{Gal}(\mathbb{E}',\mathbb{E})\triangleleft\operatorname{Gal}(\mathbb{E}',\mathbb{K})$. 由\autoref{c2-rem-2-2}，也有 $\operatorname{Gal}(\mathbb{L},\mathbb{E})\triangleleft\operatorname{Gal}(\mathbb{L},\mathbb{K})$.

现在设有 $\xi\in\mathbb{E}$ 与 $\sigma\in\operatorname{Gal}(\mathbb{L},\mathbb{K})$ 使得 $\sigma(\xi)\notin\mathbb{E}$. 由\autoref{c2-lem-2-3}，存在 $\gamma\in\operatorname{Gal}(\mathbb{L},\mathbb{E})$ 使得 $\gamma\sigma(\xi)\neq\sigma(\xi)$，这意味着 $\sigma^{-1}\gamma\sigma(\xi)\neq\xi$. 我们已证 $\operatorname{Gal}(\mathbb{L},\mathbb{E})\triangleleft\operatorname{Gal}(\mathbb{L},\mathbb{K})$，于是存在 $\widetilde{\gamma}\in\operatorname{Gal}(\mathbb{L},\mathbb{E})$ 使得 $\widetilde{\gamma}=\sigma^{-1}\gamma\sigma$. 但这样 $\widetilde{\gamma}(\xi)\neq\xi$，与 $\operatorname{Gal}(\mathbb{L},\mathbb{E})$ 固定 $\mathbb{E}$ 矛盾.

这就证明了 $\sigma(\xi)\in\mathbb{E}$ 对一切 $\xi\in\mathbb{E}$ 与一切 $\sigma\in\operatorname{Gal}(\mathbb{L},\mathbb{K})$ 成立. 于是 $\sigma|_{\mathbb{E}}$ 是 $\mathbb{E}$ 的自同构.

设 $[\tau],[\eta]\in\operatorname{Gal}(\mathbb{L},\mathbb{K})/\operatorname{Gal}(\mathbb{L},\mathbb{E})$，则
\[
\begin{aligned}
\tau\equiv\eta\ &\Leftrightarrow\ \tau\operatorname{Gal}(\mathbb{L},\mathbb{E})=\eta\operatorname{Gal}(\mathbb{L},\mathbb{E})\\
&\Leftrightarrow\ \tau=\eta\gamma\ \text{with }\gamma\in\operatorname{Gal}(\mathbb{L},\mathbb{E})\\
&\Leftrightarrow\ \eta^{-1}\tau=\gamma\ \text{with }\gamma\in\operatorname{Gal}(\mathbb{L},\mathbb{E})\\
&\Leftrightarrow\ \eta^{-1}\tau|_{\mathbb{E}}=\mathrm{id}_{\mathbb{E}}\\
&\Leftrightarrow\ \eta|_{\mathbb{E}}=\tau|_{\mathbb{E}}.
\end{aligned}
\]
这就证明了 $\operatorname{Gal}(\mathbb{L},\mathbb{K})/\operatorname{Gal}(\mathbb{L},\mathbb{E})\cong\{\sigma|_{\mathbb{E}}\mid\forall\sigma\in\operatorname{Gal}(\mathbb{L},\mathbb{K})\}$. 而右端正是 $\mathbb{E}$ 中固定 $\mathbb{K}$ 的自同构之集，即 $\operatorname{Gal}(\mathbb{E},\mathbb{K})$.
\end{proof}

\subsection{定义}

回忆 $\mathcal{F}_{\mathbb{L}}$ 是满足 $\mathbb{K}\subseteq\mathbb{E}\subseteq\mathbb{L}$ 的 $\mathbb{K}$ 的一切有限 Galois 扩张 $\mathbb{E}$ 构成的族.

\begin{definition}\label{c2-def-2-3}
设 $\mathbb{L}$ 是 $\mathbb{K}$ 的 Galois 扩张，$G=\operatorname{Gal}(\mathbb{L},\mathbb{K})$. 考虑有限指数正规子集的全体 $\mathcal{S}=\{\operatorname{Gal}(\mathbb{L},\mathbb{E})\mid\mathbb{E}\in\mathcal{F}_{\mathbb{L}}\}$. 由 $\mathcal{S}$ 作为 $G$ 的单位元邻域基所定义的拓扑，称为 $G$ 的\newterm{c2-krull-topology}{Krull 拓扑}{Krull topology}.
\end{definition}

\begin{proposition}\label{c2-prop-2-5}
Krull 拓扑是良定义的，即按定义取 $\mathcal{S}$，则 $G$ 上存在唯一的拓扑群结构，使 $\mathcal{S}$ 构成单位元的邻域基. 于是 $G:=\operatorname{Gal}(\mathbb{L},\mathbb{K})$ 配备 Krull 拓扑后是拓扑群. 此外，$\mathcal{S}$ 的元素都是开正规子群.
\end{proposition}

\begin{proof}
我们验证族 $\mathcal{S}$ 满足第 $1$ 章\autoref{c1-prop-3-1}的断言 (1)--(4).
\begin{enumerate}
    \item 设 $\operatorname{Gal}(\mathbb{L},\mathbb{E})$ 与 $\operatorname{Gal}(\mathbb{L},\mathbb{E}')$ 是 $\mathcal{S}$ 的元素. 由\autoref{c2-thm-1-5}，有
    \[
    \operatorname{Gal}(\mathbb{L},\mathbb{E}\mathbb{E}')=\operatorname{Gal}(\mathbb{L},\mathbb{E})\cap\operatorname{Gal}(\mathbb{L},\mathbb{E}').
    \]
    而 $\operatorname{Gal}(\mathbb{L},\mathbb{E}\mathbb{E}')$ 是 $\mathcal{S}$ 的元素，因为\autoref{c2-thm-1-5}保证扩张 $\mathbb{E}\mathbb{E}'$ 在 $\mathbb{K}$ 上有限且 Galois.
    \item 设 $\operatorname{Gal}(\mathbb{L},\mathbb{E})\in\mathcal{S}$，则 $\operatorname{Gal}(\mathbb{L},\mathbb{E})\operatorname{Gal}(\mathbb{L},\mathbb{E})\subseteq\operatorname{Gal}(\mathbb{L},\mathbb{E})$，因为它是一个群.
    \item 设 $\operatorname{Gal}(\mathbb{L},\mathbb{E})\in\mathcal{S}$，则 $\operatorname{Gal}(\mathbb{L},\mathbb{E})^{-1}\subseteq\operatorname{Gal}(\mathbb{L},\mathbb{E})$，因为它是一个群.
    \item 设 $\operatorname{Gal}(\mathbb{L},\mathbb{E})\in\mathcal{S}$ 且 $\sigma\in G$. 我们在\autoref{c2-prop-2-4}中已证 $\operatorname{Gal}(\mathbb{L},\mathbb{E})\triangleleft G$，从而 $\operatorname{Gal}(\mathbb{L},\mathbb{E})=\sigma\operatorname{Gal}(\mathbb{L},\mathbb{E})\sigma^{-1}$.
\end{enumerate}
\end{proof}

\begin{remark}\label{c2-rem-2-4}
若 $\mathbb{L}$ 是 $\mathbb{K}$ 的有限 Galois 扩张，则 $\operatorname{Gal}(\mathbb{L},\mathbb{K})$ 的 Krull 拓扑是离散拓扑.
\end{remark}

\begin{remark}\label{c2-rem-2-5}
从今以后，若无特别说明，我们总认为 Galois 群配备 Krull 拓扑.
\end{remark}

\subsection{性质}

\begin{theorem}\label{c2-thm-2-6}
设 $\mathbb{L}$ 是 $\mathbb{K}$ 的 Galois 扩张，令 $G=\operatorname{Gal}(\mathbb{L},\mathbb{K})$. 则配备 Krull 拓扑的群 $G$ 是
\begin{enumerate}
    \item Hausdorff 的，
    \item 紧致的，
    \item 全不连通的
\end{enumerate}
拓扑群.
\end{theorem}

\begin{proof}
\begin{enumerate}
    \item 我们要证：对一切 $\sigma,\tau\in G$ 且 $\sigma\neq\tau$，存在 $\sigma$ 与 $\tau$ 的邻域 $U$ 与 $V$，使得 $\sigma\in U$、$\tau\in V$ 且 $U\cap V=\emptyset$.

    首先注意
    \[
    \bigcap_{U\in\mathcal{S}}U=\bigcap_{\mathbb{E}\in\mathcal{F}_{\mathbb{L}}}\operatorname{Gal}(\mathbb{E},\mathbb{K})=1
    \]\footnote{左端交集跑遍 $\mathcal{S}$，故中间一项原文即如此（作 $\operatorname{Gal}(\mathbb{E},\mathbb{K})$）；按上下文应为 $\operatorname{Gal}(\mathbb{L},\mathbb{E})$.}
    （这可由\autoref{c2-rem-2-2}推出）.

    于是，若 $\sigma,\tau\in G$ 且 $\sigma\neq\tau$，则 $\sigma^{-1}\tau\neq 1$，从而（利用上式）存在 $U_{0}\in\mathcal{S}$ 使得 $\sigma^{-1}\tau\notin U_{0}$. 这样 $\tau\notin\sigma U_{0}$，即 $\tau U_{0}\cap\sigma U_{0}=\emptyset$（$G$ 是群）.
    \item 我们用 Tychonoff 定理证明紧致性. 为此构造同态
    \[
    h:G\to\prod_{\mathbb{E}\in\mathcal{F}_{\mathbb{L}}}\operatorname{Gal}(\mathbb{E},\mathbb{K})=:P
    \]
    定义为
    \[
    h(\sigma)=\prod_{\mathbb{E}\in\mathcal{F}_{\mathbb{L}}}\sigma|_{\mathbb{E}}.
    \]
    注意 $P$ 是紧致的. 事实上，每个 $\operatorname{Gal}(\mathbb{E},\mathbb{K})$ 都是离散有限群，从而紧致，于是由 Tychonoff 定理（见 \cite[Theorem~37.3]{mun75}），$P$ 作为紧致空间之积是紧致的. 如果我们证明 $h$ 是单射、连续且到 $h(G)$ 上是开映射，就得到 $G$ 与 $h(G)$ 之间的同胚（见 \cite[Theorem~26.6]{mun75}）. 进而，若再证明 $h(G)$ 是 $P$ 的闭子集（这蕴含 $h(G)$ 紧致），即得 $G$ 紧致.

    \textbf{单射性：}设 $\sigma\in G$ 满足 $h(\sigma)=1$. 这就是说对每个 $\mathbb{E}\in\mathcal{F}_{\mathbb{L}}$，$\sigma|_{\mathbb{E}}$ 都是恒等映射. 但已知 $\mathbb{L}=\bigcup_{\mathbb{E}\in\mathcal{F}_{\mathbb{L}}}\mathbb{E}$，故 $\sigma=1$.

    \textbf{连续性：}考虑复合
    \[
    G\xrightarrow{\ h\ }P\xrightarrow{\ \pi_{\mathbb{E}}\ }\operatorname{Gal}(\mathbb{E},\mathbb{K}),
    \]
    其中 $\pi_{\mathbb{E}}$ 是典范投影. 要证 $h$ 连续，只需证明对一切 $\mathbb{E}\in\mathcal{F}_{\mathbb{L}}$，$\pi_{\mathbb{E}}\circ h$ 连续. 由于我们处理的是拓扑群，只需在 $1$ 处验证：
    \[
    (\pi_{\mathbb{E}}h)^{-1}(\{1\})=\operatorname{Gal}(\mathbb{L},\mathbb{E})\in\mathcal{S}.
    \]
    注意 $\{1\}$ 在 $P$ 中是开集，因为它在每个 $\operatorname{Gal}(\mathbb{E},\mathbb{K})$ 中都是开集（$\operatorname{Gal}(\mathbb{E},\mathbb{K})$ 有限）.

    \textbf{开映射：}设 $\mathbb{E}\in\mathcal{F}_{\mathbb{L}}$，则
    \[
    h(\operatorname{Gal}(\mathbb{L},\mathbb{E}))=h(G)\cap\big(\prod_{\substack{\mathbb{E}'\neq\mathbb{E}\\ \mathbb{E}'\in\mathcal{F}_{\mathbb{L}}}}\operatorname{Gal}(\mathbb{E}',\mathbb{K})\times\{1\}\big),
    \]
    这是 $h(G)$ 中的开集.

    \textbf{$h(G)$ 闭：}对每对满足 $\mathbb{E}_{2}\subseteq\mathbb{E}_{1}$ 的 $\mathbb{E}_{1},\mathbb{E}_{2}\in\mathcal{F}_{\mathbb{L}}$，定义 $M_{\mathbb{E}_{1}|\mathbb{E}_{2}}:=\{\Sigma\in P\mid \pi_{\mathbb{E}_{1}}(\Sigma)|_{\mathbb{E}_{2}}=\pi_{\mathbb{E}_{2}}(\Sigma)\}$. 我们先证 $M_{\mathbb{E}_{1}|\mathbb{E}_{2}}$ 在 $P$ 中闭. 由于 $\mathbb{E}_{2}$ 是 $\mathbb{K}$ 的有限扩张，可设 $\operatorname{Gal}(\mathbb{E}_{2},\mathbb{K})=\{f_{1},\ldots,f_{r}\}$. 以 $S_{i}$ 记 $f_{i}$ 到 $\mathbb{E}_{1}$ 的扩张之集. 于是
    \[
    M_{\mathbb{E}_{1}|\mathbb{E}_{2}}=\bigcup_{i=1}^{r}\big(\prod_{\substack{\mathbb{E}\neq\mathbb{E}_{1},\mathbb{E}_{2}\\ \mathbb{E}\in\mathcal{F}_{\mathbb{L}}}}\operatorname{Gal}(\mathbb{E},\mathbb{K})\times S_{i}\times\{f_{i}\}\big).
    \]
    有限并中的每个集合都是闭集（作为闭集之积），故 $M_{\mathbb{E}_{1}|\mathbb{E}_{2}}$ 闭.

    我们有
    \[
    h(G)\subseteq\bigcap_{\mathbb{E}_{1}\supseteq\mathbb{E}_{2}}M_{\mathbb{E}_{1}|\mathbb{E}_{2}}
    \]
    若能证明反包含，就得 $h(G)$ 在 $P$ 中闭. 现在，若 $\Sigma\in\bigcap_{\mathbb{E}_{1}\supseteq\mathbb{E}_{2}}M_{\mathbb{E}_{1}|\mathbb{E}_{2}}$，当 $x\in\mathbb{E}$ 时定义 $\sigma(x)=\pi_{\mathbb{E}}(\Sigma)(x)$，这就给出一个自同构 $\sigma:\mathbb{L}\to\mathbb{L}$. 由 $\Sigma\in\cap M_{\mathbb{E}_{1}|\mathbb{E}_{2}}$（并利用\autoref{c2-rem-2-2}），这样的 $\sigma$ 是良定义的. 我们有 $h(\sigma)=\prod_{\mathbb{E}\in\mathcal{F}_{\mathbb{L}}}\pi_{\mathbb{E}}(\Sigma)$，从而
    \[
    h(G)\supseteq\bigcap_{\mathbb{E}_{1}\supseteq\mathbb{E}_{2}}M_{\mathbb{E}_{1}|\mathbb{E}_{2}},
    \]
    故 $h(G)$ 闭.
    \item 由于我们处理的是拓扑群，只需证明 $1$ 的连通分支 $H$ 等于 $\{1\}$.

    对每个 $U\in\mathcal{S}$，定义 $U_{H}:=U\cap H$. 因 $1\in U$，有 $U_{H}\neq\emptyset$；又因 $U$ 开，$U_{H}$ 在 $H$ 中开. 再定义
    \[
    V_{H}=\bigcup_{x\in H\setminus U_{H}}xU_{H}.
    \]
    于是 $V_{H}$ 在 $H$ 中开，因为每个 $xU_{H}$ 都是开集. 我们还有 $U_{H}\cap V_{H}=\emptyset$ 且 $H=U_{H}\cup V_{H}$. 但 $H$ 是连通的，故必有 $V_{H}=\emptyset$；否则 $U_{H}|V_{H}$ 就是 $H$ 的一个分离. 这意味着 $U_{H}=H$，即对一切 $U\in\mathcal{S}$ 有 $U\cap H=H$. 从而
    \[
    H\subseteq\bigcap_{U\in\mathcal{S}}U=\{1\},
    \]
    即 $H=\{1\}$.
\end{enumerate}
\end{proof}

\section{无限 Galois 理论的基本定理}

本节要把有限扩张情形的 Galois 理论主定理推广开来，为此我们还将推广一些命题和引理.

\subsection{一些必要的命题}

第一条命题是\autoref{c2-lem-2-2}的推广，但此时我们是在拓扑群的框架下工作.

\begin{proposition}\label{c2-prop-3-1}
设 $\mathbb{L}$ 是 $\mathbb{K}$ 的 Galois 扩张. 对每个满足 $\mathbb{K}\subseteq\mathbb{E}\subseteq\mathbb{L}$ 的 $\mathbb{K}$ 的有限 Galois 扩张 $\mathbb{E}$，映射
\[
\begin{aligned}
r:\operatorname{Gal}(\mathbb{L},\mathbb{K})&\longrightarrow\operatorname{Gal}(\mathbb{E},\mathbb{K})\\
\sigma&\longmapsto\sigma|_{\mathbb{E}}
\end{aligned}
\]
是连续满射.
\end{proposition}

\begin{proof}
\autoref{c2-lem-2-2}已证该映射是满射.

我们证明单位元 $1_{\operatorname{Gal}(\mathbb{E},\mathbb{K})}$ 的原像在 $\operatorname{Gal}(\mathbb{L},\mathbb{K})$ 中开（拓扑群是均匀的，这就足够了）. 而 $r^{-1}(\{1_{\operatorname{Gal}(\mathbb{E},\mathbb{K})}\})=\operatorname{Gal}(\mathbb{L},\mathbb{E})$，由 Krull 拓扑的定义它是开集.
\end{proof}

\begin{proposition}\label{c2-prop-3-2}
设 $\mathbb{L}$ 是 $\mathbb{K}$ 的 Galois 扩张，记 $G:=\operatorname{Gal}(\mathbb{L},\mathbb{K})$. 则以下断言成立.
\begin{enumerate}
    \item 对每个中间域 $\mathbb{M}$（即 $\mathbb{K}\subseteq\mathbb{M}\subseteq\mathbb{L}$），$\mathbb{L}$ 在 $\mathbb{M}$ 上是 Galois 扩张. 此外，$\operatorname{Gal}(\mathbb{L},\mathbb{M})$ 在 $G$ 中闭，且 $\mathbb{L}^{\operatorname{Gal}(\mathbb{L},\mathbb{M})}=\mathbb{M}$.
    \item 对 $G$ 的每个子群 $H$，$\operatorname{Gal}(\mathbb{L},\mathbb{L}^{H})$ 是 $H$ 的拓扑闭包.
\end{enumerate}
\end{proposition}

\begin{proof}
\begin{enumerate}
    \item 第一条断言蕴含于\autoref{c2-prop-1-4}.

    以 $\mathcal{F}_{\mathbb{M}}$ 记满足 $\mathbb{K}\subseteq\mathbb{E}\subseteq\mathbb{M}$ 的 $\mathbb{K}$ 的一切有限 Galois 扩张 $\mathbb{E}$ 构成的族. 则 $\mathbb{M}$ 是 $\mathcal{F}_{\mathbb{M}}$ 中一切元素之并，这意味着
    \[
    \operatorname{Gal}(\mathbb{L},\mathbb{M})=\cap_{\mathbb{E}\in\mathcal{F}_{\mathbb{M}}}\operatorname{Gal}(\mathbb{L},\mathbb{E}).
    \]
    而各 $\operatorname{Gal}(\mathbb{L},\mathbb{E})$ 都是 $G$ 的开子群，因而也是闭的（见第 $1$ 章\autoref{c1-prop-4-4}(2) 与\autoref{c2-thm-2-6}）. 于是 $\operatorname{Gal}(\mathbb{L},\mathbb{M})$ 作为闭集之交是闭的.

    \autoref{c2-lem-2-3}蕴含最后一条断言.
    \item 我们有 $H\subseteq\operatorname{Gal}(\mathbb{L},\mathbb{L}^{H})$，但 $\operatorname{Gal}(\mathbb{L},\mathbb{L}^{H})$ 是闭的，故 $\overline{H}\subseteq\operatorname{Gal}(\mathbb{L},\mathbb{L}^{H})$.

    现在设 $\sigma\in G\setminus\overline{H}$，则存在 $\mathbb{E}\in\mathcal{F}_{\mathbb{L}}$ 使得 $\sigma\operatorname{Gal}(\mathbb{L},\mathbb{E})\cap H=\emptyset$. 于是 $\sigma\notin H\operatorname{Gal}(\mathbb{L},\mathbb{E})$. 由于 $\sigma$ 不能写成 $\sigma=h\tau$（其中 $h\in H$、$\tau\in\operatorname{Gal}(\mathbb{L},\mathbb{E})$；否则将有 $\sigma(\alpha)=h(\alpha)$ 对一切 $\alpha\in\mathbb{E}$ 成立），故存在 $\alpha\in\mathbb{E}$ 使得 $H$ 固定 $\alpha$（$\Rightarrow\alpha\in\mathbb{L}^{H}$）但 $\sigma(\alpha)\neq\alpha$. 这意味着 $\sigma\notin\operatorname{Gal}(\mathbb{L},\mathbb{L}^{H})$.

    从而 $\operatorname{Gal}(\mathbb{L},\mathbb{L}^{H})\subseteq\overline{H}$.
\end{enumerate}
\end{proof}

\subsection{主定理}

下面证明一个与有限 Galois 扩张的主定理（见\autoref{c2-thm-1-5}）类似的定理.

\begin{theorem}\label{c2-thm-3-3}
设 $\mathbb{L}$ 是 $\mathbb{K}$ 上的 Galois 扩张，$G$ 为其 Galois 群. 存在双射
\[
\{H\mid H\ \text{为}\ G\ \text{的闭子群}\}\leftrightarrow\{\mathbb{M}\ \text{中间域}\mid\mathbb{K}\subseteq\mathbb{M}\subseteq\mathbb{L}\}
\]
它由映射 $H\mapsto\mathbb{L}^{H}$（把闭子群映到它的固定域）与 $\mathbb{M}\mapsto\operatorname{Gal}(\mathbb{L},\mathbb{M})$（把中间域映到它的 Galois 群）给出. 此外，
\begin{enumerate}
    \item 若 $H_{1}$ 与 $H_{2}$ 是 $G$ 的两个闭子群，则 $H_{1}\supseteq H_{2}$ 当且仅当 $\mathbb{L}^{H_{1}}\subseteq\mathbb{L}^{H_{2}}$.
    \item $G$ 的闭子群 $H$ 是开集当且仅当 $\mathbb{L}^{H}$ 在 $\mathbb{K}$ 上的次数有限. 此时 $[G:H]=[\mathbb{L}^{H}:\mathbb{K}]$，即指数与次数一致.
    \item 对每个元素 $\sigma\in G$ 与 $G$ 的每个子群 $H$，有 $\mathbb{L}^{\sigma H\sigma^{-1}}=\sigma(\mathbb{L}^{H})$；若 $\mathbb{M}$ 是中间域，则 $\operatorname{Gal}(\mathbb{L},\sigma\mathbb{M})=\sigma\operatorname{Gal}(\mathbb{L},\mathbb{M})\sigma^{-1}$.
    \item 若 $H$ 是 $G$ 的闭子群，则下述等价成立：

    $H$ 正规当且仅当 $\mathbb{L}^{H}$ 在 $\mathbb{K}$ 上是 Galois 扩张. 此时 $\operatorname{Gal}(\mathbb{L}^{H},\mathbb{K})\cong G/H$.
\end{enumerate}
\end{theorem}

\begin{proof}
要证第一条断言，只需证明映射 $H\mapsto\mathbb{L}^{H}$ 与 $\mathbb{M}\mapsto\operatorname{Gal}(\mathbb{L},\mathbb{M})$ 互为逆映射. 先设 $H$ 是 $G$ 的闭子群，则 $\mathbb{L}$ 在 $\mathbb{L}^{H}$ 上是 Galois 扩张且 $\operatorname{Gal}(\mathbb{L},\mathbb{L}^{H})=H$（见\autoref{c2-prop-3-2}）. 反之，若 $\mathbb{M}$ 是中间域，则 $\operatorname{Gal}(\mathbb{L},\mathbb{M})$ 是 $G$ 的闭子群且 $\mathbb{L}^{\operatorname{Gal}(\mathbb{L},\mathbb{M})}=\mathbb{M}$（见\autoref{c2-prop-3-2}(1)）.
\begin{enumerate}
    \item 设 $H_{1},H_{2}$ 是 $G$ 的闭子群，由定义即有
    \[
    H_{2}\subseteq H_{1}\Rightarrow\mathbb{L}^{H_{1}}\subseteq\mathbb{L}^{H_{2}}\Rightarrow\operatorname{Gal}(\mathbb{L},\mathbb{L}^{H_{1}})\supseteq\operatorname{Gal}(\mathbb{L},\mathbb{L}^{H_{2}})
    \]
    于是 $\operatorname{Gal}(\mathbb{L},\mathbb{L}^{H_{i}})=H_{i}$ 蕴含结论.
    \item 由第 $1$ 章\autoref{c1-prop-4-4}，$H$ 是 $G$ 的闭子群且指数有限蕴含 $H$ 开. 反之用同一条命题，若 $H$ 是 $G$ 的开子群，则 $H$ 闭且指数有限（回忆 $G$ 紧致——见\autoref{c2-thm-2-6}）.

    现在设 $H$ 闭且指数有限. 考虑连续映射
    \[
    \begin{aligned}
    r:\operatorname{Gal}(\mathbb{L},\mathbb{K})&\longrightarrow\operatorname{Hom}_{\mathbb{K}}(\mathbb{L}^{H},\mathbb{K})\\
    \sigma&\longmapsto\sigma|_{\mathbb{L}^{H}}.
    \end{aligned}
    \]
    显然 $\ker(r)=\operatorname{Gal}(\mathbb{L},\mathbb{L}^{H})$ 且 $\operatorname{im}(r)=\operatorname{Hom}_{\mathbb{K}}(\mathbb{L}^{H},\mathbb{K})$（最后一条断言应用\autoref{c2-lem-2-1}）. 于是由群的第一同构定理，
    \[
    \operatorname{Gal}(\mathbb{L},\mathbb{K})/\ker(r)\cong\operatorname{Hom}_{\mathbb{K}}(\mathbb{L}^{H},\mathbb{K}),
    \]
    等价地
    \[
    G/H\cong\operatorname{Hom}_{\mathbb{K}}(\mathbb{L}^{H},\mathbb{K}).
    \]
    \item 设 $\mathbb{M}$ 是中间域，$H=\operatorname{Gal}(\mathbb{L},\mathbb{M})$. 若 $\sigma\in G$，我们要证 $\sigma H\sigma^{-1}\leftrightarrow\sigma\mathbb{M}$.

    为此，设 $\tau\in G$、$\alpha\in\mathbb{L}$. 则 $\tau(\alpha)=\alpha$ 当且仅当 $\sigma\tau\sigma^{-1}(\sigma(\alpha))=\sigma(\alpha)$. 于是 $\operatorname{Gal}(\mathbb{L},\sigma\mathbb{M})=\sigma\operatorname{Gal}(\mathbb{L},\mathbb{M})\sigma^{-1}$. 且 $\sigma(\mathbb{L}^{H})=\sigma\mathbb{M}=\mathbb{L}^{\sigma H\sigma^{-1}}$.
    \item 设 $\mathbb{M}$ 是中间域，$H=\operatorname{Gal}(\mathbb{L},\mathbb{M})$. 由 (3)，$H$ 在 $G$ 中正规当且仅当 $\mathbb{M}$ 在 $G$ 的自然作用下不变. 后一条断言等价于说 $\mathbb{M}$ 是 $\mathbb{K}$ 的在 $G$ 下不变的一族有限扩张之并. 而这利用\autoref{c2-thm-1-5}的第 (3) 条，正相当于要求这些有限扩张都在 $\mathbb{K}$ 上 Galois. 于是 $\mathbb{M}$ 是 Galois 的.

    而同构则由 (2) 推出.
\end{enumerate}
\end{proof}

下面这个例子说明：若去掉拓扑条件，主定理就不成立.

\begin{example}\label{c2-ex-3-1}
考虑素数域 $\mathbb{F}_p$ 及其代数闭包 $\overline{\mathbb{F}}_p$. 注意 $\overline{\mathbb{F}}_p$ 是 $\mathbb{F}_p$ 的 Galois 扩张. 由于 $\overline{\mathbb{F}}_p$ 必须包含一切形如 $X^{p^{n}}-X$（$n\geq 1$）的多项式的根，有 $\mathbb{F}_{p^{n}}\subseteq\overline{\mathbb{F}}_p$. 又 $\overline{\mathbb{F}}_p$ 可写成中间有限扩张之并，故必有 $\overline{\mathbb{F}}_p=\cup_{n\geq 1}\mathbb{F}_{p^{n}}$.

现在考虑 $G=\operatorname{Gal}(\overline{\mathbb{F}}_p,\mathbb{F}_p)$. 以 $\varphi$ 记 Frobenius 自同构，即
\[
\varphi:\overline{\mathbb{F}}_p\to\overline{\mathbb{F}}_p:\alpha\mapsto\alpha^{p}.
\]
由于 $\mathbb{F}_p^{*}$ 是 $p-1$ 阶循环群，对每个 $\alpha\in\mathbb{F}_p$ 有 $\varphi(\alpha)=\alpha$. 故 $\varphi\in G$.

定义 $H=\langle\varphi\rangle$，即由 $\varphi$ 的一切方幂生成的 $G$ 的子群. 下面证明关于 $H$ 的两条引人注目的断言.
\begin{enumerate}
    \item $H\neq G$

    为此要在 $G$ 中找到一个不属于 $H$ 的元素 $\tau$.

    我们先证：存在 $\mathbb{F}_p$ 的无限中间扩张 $\mathbb{M}$，使得 $\mathbb{F}_p\subsetneq\mathbb{M}\subsetneq\overline{\mathbb{F}}_p$. 考虑 $\mathbb{M}=\cup_{n\geq 0}\mathbb{F}_{p^{2n}}$，显然 $\mathbb{M}$ 是无限扩张（因为对一切 $n\geq 0$ 有 $\mathbb{F}_{p^{2n}}\subseteq\mathbb{M}$），且 $\mathbb{M}$ 在 $\mathbb{F}_p$ 上 Galois（因为它是 Galois 扩张之并）.

    下面再证 $\mathbb{M}\neq\overline{\mathbb{F}}_p$. 为此考虑 $\mathbb{F}_{p^{3}}$，显然 $[\mathbb{F}_{p^{3}}:\mathbb{F}_p]=3$ 且 $\mathbb{F}_{p^{3}}\subseteq\overline{\mathbb{F}}_p$，但 $\mathbb{F}_{p^{3}}\nsubseteq\mathbb{M}$. 事实上，若设 $\mathbb{F}_{p^{3}}=\mathbb{F}_p[\alpha]$ 且 $\alpha\in\mathbb{M}$，则对某个 $n\geq 0$ 有 $\alpha\in\mathbb{F}_{p^{2n}}$，从而 $\mathbb{F}_3$\footnote{原文即如此；按上下文应为 $\mathbb{F}_{p^{3}}$.}$\subseteq\mathbb{F}_{p^{2n}}$. 但这样对扩张用次数公式就有：
    \[
    2^{n}=[\mathbb{F}_{p^{2n}}:\mathbb{F}_p]=[\mathbb{F}_{p^{2n}}:\mathbb{F}_{p^{3}}][\mathbb{F}_{p^{3}}:\mathbb{F}_p]=[\mathbb{F}_{p^{2n}}:\mathbb{F}_{p^{3}}]\cdot 3,
    \]
    这是矛盾. 故 $\mathbb{F}_p\subsetneq\mathbb{M}\subsetneq\overline{\mathbb{F}}_p$.

    现在由\autoref{c2-prop-1-4}，$\overline{\mathbb{F}}_p$ 在 $\mathbb{M}$ 上是 Galois 扩张. 取 $\tau\in\operatorname{Gal}(\overline{\mathbb{F}}_p,\mathbb{M})\setminus\{\mathrm{Id}\}\subseteq G$. 若假设 $H=G$，则存在 $n\geq 1$ 使得 $\tau=\varphi^{n}$. 于是 $\varphi^{n}$ 固定 $\mathbb{M}$，故 $\mathbb{M}\subseteq\mathbb{F}_{p^{n}}$. 但 $\mathbb{M}$ 是无限的，矛盾.
    \item $\mathbb{F}_p^{H}=\overline{\mathbb{F}}_p$\footnote{原文即如此（横线在右端）；按上下文，此处及末行的固定域记号应为 $\overline{\mathbb{F}}_p^{\,H}=\mathbb{F}_p$ 与 $\mathbb{F}_p=\overline{\mathbb{F}}_p^{\,G}$.}

    由定义 $\mathbb{F}_p^{H}\subseteq\overline{\mathbb{F}}_p$，故只需证反包含. 设 $\alpha\in\overline{\mathbb{F}}_p$，则 $\alpha\in\mathbb{F}_{p^{n}}$ 对某个 $n\geq 1$ 成立. 这意味着 $\alpha^{p^{n}}-\alpha=0$（见\autoref{c2-ex-1-17}）. 于是 $\alpha^{p^{n}}=\alpha$，从而 $\varphi^{n}(\alpha)=\alpha$. 故 $\alpha\in\mathbb{F}_p^{H}$.
\end{enumerate}
我们得出结论：$H$ 与 $G$ 是两个不同的群，但它们的固定域相同，因为 $\overline{\mathbb{F}}_p=\mathbb{F}_p^{G}$. 这与定理矛盾，除非 $G$ 是 $H$ 的拓扑闭包.
\end{example}

% !TeX root = ../main.tex
% 第 3 章 投射有限群 （原书 CHAPTER 3, Profinite Groups, 印刷页 23-33）

\chapter{投射有限群}
\label{ch:3}

本章定义投射有限群的概念. 投射有限群定义为一类特殊系统的极限. 它们在刻画 Galois 群时十分有用：事实上，每一个 Galois 群都是投射有限群，而每一个投射有限群都可以实现为一个 Galois 群.

\section{定义与例子}

\subsection{逆极限的概念}

\subsubsection{定义}

\begin{definition}\label{c3-def-1-1}
设 $I$ 是一个非空集合，其上配有自反且传递的二元关系 $\leq$，并满足附加性质：任意两个元素都有上界. 则称 $I$ 为 \newterm{c3-directed-set}{有向集}{directed set}.
\end{definition}

\begin{definition}\label{c3-def-1-2}
设 $I$ 是有向集. 以 $I$ 为指标的拓扑空间的 \newterm{c3-inverse-system}{逆系统}{inverse system} $(X_i,f_{ij})_I$ 由一簇拓扑空间 $\{X_i\mid i\in I\}$ 与一簇连续映射 $\{f_{ij}:X_j\to X_i\mid i,j\in I,\,i\leq j\}$ 构成，且满足
\begin{enumerate}[label=(\arabic*),leftmargin=*]
    \item 对一切 $i\in I$，$f_{ii}$ 是 $X_i$ 上的恒等映射；
    \item 对一切 $i\leq j\leq k\in I$，$f_{ij}\circ f_{jk}=f_{ik}$.
\end{enumerate}
\end{definition}

\begin{remark}\label{c3-rem-1-3}
前面的定义不仅适用于拓扑空间范畴，还可以推广到任意范畴.
\end{remark}

\begin{examples}\label{c3-ex-1-4}
\
\begin{enumerate}[label=(\arabic*),leftmargin=*]
    \item 考虑有向集 $(\mathbb{Z},>)$，其中对一切 $m,n\in\mathbb{Z}$，$m>n$ 当且仅当 $n$ 整除 $m$. 定义 $G_m=\mathbb{Z}/m\mathbb{Z}$，并考虑典范投影
    \[
    f_{nm}:\mathbb{Z}/m\mathbb{Z}\to\mathbb{Z}/n\mathbb{Z}:k+m\mathbb{Z}\mapsto k+n\mathbb{Z},
    \]
    当 $n\mid m$ 时它是良定义的. 于是 $(G_m,f_{nm})_{\mathbb{Z}}$ 是一个逆系统.
    \item 设 $p$ 为素数. 现在考虑 $I=\mathbb{N}$ 赋予通常序. 定义 $G_n=\mathbb{Z}/p^n\mathbb{Z}$，并考虑典范投影
    \[
    f_{mn}:\mathbb{Z}/p^n\mathbb{Z}\to\mathbb{Z}/p^m\mathbb{Z}
    \]
    当 $m\leq n$ 时它是良定义的. 于是 $(G_n,f_{mn})_{\mathbb{N}}$ 是一个逆系统.
    \item 设 $G$ 为任意群. 以 $\mathcal{S}$ 记其全体有限指数正规子群之集. 在 $\mathcal{S}$ 上规定序：$U<V$ 当且仅当 $V\subseteq U$，其中 $U,V\in\mathcal{S}$. 定义
    \[
    f_{UV}:G/V\to G/U:gV\mapsto gU,
    \]
    当 $U\leq V$ 时. 于是 $(G/U,f_{UV})_{\mathcal{S}}$ 是一个逆系统.
\end{enumerate}
\end{examples}

\begin{definition}\label{c3-def-1-5}
若 $Y$ 是拓扑空间，称连续映射族 $\{\psi_i:Y\to X_i\}_{i\in I}$ 是\emph{相容的}，若对一切 $i\leq j\in I$ 有 $f_{ij}\psi_j=\psi_i$，即下图交换：
\[
\begin{tikzcd}
& Y \arrow[dl, "\psi_j"'] \arrow[dr, "\psi_i"] & \\
X_j \arrow[rr, "f_{ij}"'] & & X_i.
\end{tikzcd}
\]
\end{definition}

\begin{definition}\label{c3-def-1-6}
拓扑空间逆系统 $(X_i,f_{ij})_I$ 的 \newterm{c3-inverse-limit}{逆极限}{inverse limit} $(X,\varphi_i)$ 是指配有相容族 $\{\varphi_i:X\to X_i\}_{i\in I}$ 的拓扑空间 $X$，它满足如下泛性质：

对每个拓扑空间 $Y$ 与每个相容族 $\{\psi_i:Y\to X_i\}_{i\in I}$，存在唯一的连续映射 $\theta:Y\to X$，使得对一切 $i\leq j\in I$ 下图交换.
\[
\begin{tikzcd}
& Y \arrow[d, dashed, "\exists!\,\theta"'] \arrow[ddl, shift right=3pt, "\psi_j"'] \arrow[ddr, shift left=3pt, "\psi_i"] & \\
& X \arrow[dl, "\varphi_j"'] \arrow[dr, "\varphi_i"'] & \\
X_j \arrow[rr, "f_{ij}"'] & & X_i.
\end{tikzcd}
\]
\end{definition}

\begin{remark}\label{c3-rem-1-7}
内射极限也称为射影极限.\footnote{原文即如此（"An injective limit is also called a projective limit."）；按上下文应为"逆极限亦称射影极限"，injective 疑为 inverse 之误.}
\end{remark}

\subsubsection{逆极限的存在性与唯一性}

\begin{proposition}\label{c3-prop-1-1}
设 $(X_i,f_{ij})_I$ 是拓扑空间的逆系统. 若 $(X,\varphi_i)_I$ 与 $(Y,\psi_i)_I$ 都是 $(X_i,f_{ij})_I$ 的逆极限，则存在同胚 $\Phi:X\to Y$ 使得对一切 $i\in I$ 有 $\psi_i\circ\Phi=\varphi_i$.
\end{proposition}

\begin{proof}
由于 $(X,\varphi_i)_I$ 是逆极限而 $(Y,\psi_i)_I$ 是相容族，存在唯一的 $\Phi:X\to Y$ 使得对一切 $i\in I$ 有 $\psi_i\circ\Phi=\varphi_i$.

由于 $(Y,\psi_i)_I$ 是逆极限而 $(X,\varphi_i)_I$ 是相容族，存在唯一的 $\Psi:Y\to X$ 使得对一切 $i\in I$ 有 $\varphi_i\circ\Psi=\psi_i$.

于是 $\Psi\circ\Phi:X\to X$ 且对一切 $i\in I$ 有 $\varphi_i\Psi\Phi=\psi_i\Phi=\varphi_i$；此外对一切 $i\in I$ 有 $\varphi_i\mathrm{id}_X=\varphi_i$. 由\autoref{c3-def-1-6} 中态射的唯一性，$\Psi\circ\Phi=\mathrm{id}_X$.

同法可证 $\Phi\circ\Psi=\mathrm{id}_Y$.
\end{proof}

\begin{proposition}\label{c3-prop-1-2}
设 $(X_i,f_{ij})_I$ 是拓扑空间的逆系统. 以 $X$ 表示使下图交换的元素 $x\in\prod_{k\in I}X_k$ 之集
\[
\begin{tikzcd}
& \prod_{k\in I}X_k \arrow[dl, "\pi_j"'] \arrow[dr, "\pi_i"] & \\
X_j \arrow[rr, "f_{ij}"'] & & X_i
\end{tikzcd}
\]
这就是说，$X$ 是满足 $f_{ij}\pi_j(x)=\pi_i(x)$（对一切 $i\leq j\in I$）的一切元素 $x\in\prod_{k\in I}X_k$ 之集，其中 $\pi_i$ 是典范投影. 对一切 $i\in I$ 定义 $\varphi_i:=\pi_i|_X$. 则 $(X,\varphi_i)_I$ 是 $(X_i,f_{ij})_I$ 的逆极限.
\end{proposition}

\begin{remark}\label{c3-rem-1-8}
我们给 $\prod_{k\in I}X_k$ 赋予乘积拓扑，给 $X$ 赋予子空间拓扑.
\end{remark}

\begin{proof}
由乘积拓扑的定义，映射 $\varphi_i$ 连续；又由 $X$ 的定义有 $f_{ij}\varphi_j=\varphi_i$. 这就是说 $\{\varphi_i:X\to X_i\mid i\in I\}$ 是相容族.

现在考虑另一个相容族 $\{\psi_i:Y\to X_i\mid i\in I\}$. 我们须证存在唯一的连续映射 $\theta:Y\to X$ 使得对一切 $i\in I$ 有 $\varphi_i\theta=\psi_i$. 为此，定义
\[
\begin{aligned}
\overline{\theta}:Y&\longrightarrow\prod_{k\in I}X_k\\
y&\longmapsto\{\psi_k(y)\}_{k\in I}.
\end{aligned}
\]
我们有 $\pi_i\overline{\theta}=\psi_i$ 对一切 $i\in I$ 连续，这就推出 $\overline{\theta}$ 的连续性. 此外 $\overline{\theta}$ 的像含于 $X$，事实上
\[
f_{ij}\circ\pi_j(\overline{\theta}(y))=f_{ij}\psi_j(y)=\psi_i(y)=\pi_i(\overline{\theta}(y))\quad\forall\, i\leq j\in I,\ \forall\, y\in Y.
\]
这说明我们可以定义映射
\[
\begin{aligned}
\theta:Y&\longrightarrow X\\
y&\longmapsto\overline{\theta}(y)
\end{aligned}
\]
它是连续的. 此映射还满足对一切 $i\in I$ 有 $\varphi_i\theta=\psi_i$. 注意 $\theta$ 的唯一性由其构造即得.
\end{proof}

\begin{notation}\label{c3-not-1-9}
现在可以把逆系统 $(X_i,f_{ij})_I$ 的逆极限记作 $\varprojlim_I X_i$. 并以 $s\varprojlim_I X_i$ 表示上面给出的那个特定构造.
\end{notation}

\subsubsection{性质}

下面给出逆极限的若干性质，后文将会用到.

\begin{proposition}\label{c3-prop-1-3}
设 $(X_i,f_{ij})_I$ 是以 $I$ 为指标的逆系统，记 $X=\varprojlim X_i$.
\begin{enumerate}[label=(\arabic*),leftmargin=*]
    \item 若每个 $X_i$ 都是 Hausdorff 空间，则 $X$ 也是.
    \item 若每个 $X_i$ 都是全不连通空间，则 $X$ 也是.
    \item 若每个 $X_i$ 都是 Hausdorff 空间，则 $s\varprojlim X_i$ 在 $P=\prod_{i\in I}X_i$ 中是闭集.
    \item 若每个 $X_i$ 都是紧致 Hausdorff 空间，则 $X$ 也是.
\end{enumerate}
\end{proposition}

\begin{proof}
由逆极限的唯一性，只需对 $X=s\varprojlim X_i$ 证明结论.
\begin{enumerate}[label=(\arabic*),leftmargin=*]
    \item Hausdorff 空间的乘积与子空间总是 Hausdorff 空间，故结论成立.
    \item 同理.
    \item 由一个拓扑学结果：若 $f,g:Y\to Z$ 是连续映射而 $Z$ 是 Hausdorff 空间，则集合 $\{y\mid f(y)=g(y)\}$ 在 $Y$ 中闭. 由于
    \[
    s\varprojlim X_i=\bigcap_{j\geq i}\{p\in P\mid f_{ij}\pi_j(p)=\pi_i(p)\}
    \]
    （其中 $\pi_i$ 是从 $P$ 到 $X_i$ 的典范投影），若每个 $X_i$ 是 Hausdorff 空间，则 $s\varprojlim X_i$ 是闭集之交，从而 $s\varprojlim X_i$ 是闭集.
    \item 此结论来自两个事实：紧致空间的每个闭子空间是紧致的；紧致空间的乘积是紧致的.
\end{enumerate}
\end{proof}

\subsection{投射有限群}

现在考虑拓扑群的逆极限.

\begin{proposition}\label{c3-prop-1-4}
拓扑群的逆系统 $(G_i,f_{ij})_I$ 的逆极限 $\varprojlim G_i$ 是拓扑群.
\end{proposition}

\begin{proof}
笛卡儿积 $\prod_{k\in I}G_k$ 是拓扑群，这就是说只须证明 $G:=\varprojlim G_i$ 是子群.

显然 $G$ 非空（$1_{\prod_{k\in I}G_k}\in G$）且对乘法与取逆封闭. 事实上 $G=\{g\in\prod_{k\in I}G_k\mid f_{ij}\pi_j(g)=\pi_i(g)\ \forall\, i\leq j\}$，而一切 $f_{ij}$ 与 $\pi_k$（$i,j,k\in I$）都是群同态. 于是 $G$ 是 $\prod_{k\in I}G_k$ 的子群，从而赋予诱导拓扑后是拓扑子群.
\end{proof}

\begin{definition}\label{c3-def-1-10}
\newterm{c3-profinite-group}{投射有限群}{profinite group} $G$ 是指同构于某个有限群逆系统 $(G_i,f_{ij})_I$ 之逆极限的拓扑群，其中各 $G_i$ 均赋予离散拓扑.
\end{definition}

\begin{examples}\label{c3-ex-1-11}
\
\begin{enumerate}[label=(\arabic*),leftmargin=*]
    \item 记 $\widehat{\mathbb{Z}}:=\varprojlim_{m\in\mathbb{Z}}\mathbb{Z}/m\mathbb{Z}$ 为逆系统 $(\mathbb{Z}/m\mathbb{Z},f_{nm})_{\mathbb{Z}}$ 的逆极限（见例~\ref{c3-ex-1-4} (2)\footnote{原文即如此；按上下文应为 例 1.4 (1).}）. 称 $\widehat{\mathbb{Z}}$ 为 \emph{Prüfer 群}. 注意 $\widehat{\mathbb{Z}}$ 是满足
    \[
    a_m\equiv a_n\bmod n,\quad\forall\, n\mid m.
    \]
    的序列 $(a_m)$ 的全体等价类之集.
    \item 定义 $\mathbb{Z}_p:=\varprojlim_{n\in\mathbb{N}}\mathbb{Z}/p^n\mathbb{Z}$ 为逆系统 $(\mathbb{Z}/p^n\mathbb{Z},f_{mn})_{\mathbb{N}}$ 的逆极限（见例~\ref{c3-ex-1-4} (2)）. $\mathbb{Z}_p$ 是 $p$-进整数的加法群.
    \item 每个有限群 $G$ 特别地都是投射有限群. 事实上它是平凡逆系统 $(G,\mathrm{Id}_G)_{\{1\}}$ 的极限.
    \item 整数群 $\mathbb{Z}$ 不是投射有限群.
\end{enumerate}
\end{examples}

\begin{definition}\label{c3-def-1-12}
设 $G$ 为任意群，以 $\mathcal{S}$ 记其全体有限指数正规子群之集. 如例~\ref{c3-ex-1-4} 所述，$(G/U,f_{UV})_{\mathcal{S}}$ 是逆系统，我们定义
\[
\widehat{G}:={\varprojlim_{U\in\mathcal{S}}}G/U
\]
为 $G$ 的射完备.
\end{definition}

\section{投射有限群的一般性质}

本节的目标是给出投射有限群的一个特殊刻画. 事实上，我们将证明：投射有限群恰是紧致、全不连通、Hausdorff 群.

\subsection{若干引理与命题}

\begin{proposition}\label{c3-prop-2-1}
设 $(G_i,f_{ij})_I$ 是有限群（赋予离散拓扑）的逆系统. 我们构造 $G:=s\varprojlim_{i\in I}G_i$，并以 $\varphi_i:G\to G_i$ 记投影同态. 则 $\{\ker(\varphi_i)\mid i\in I\}$ 是 $G$ 的单位元 $1$ 的开邻域基本系.
\end{proposition}

\begin{proof}
考虑 $\prod_{i\in I}G_i$ 中形如
\[
\bigl(\prod_{i\neq i_1,\ldots,i_t}G_i\bigr)\times\{1\}_{i_1}\times\ldots\times\{1\}_{i_t}
\]
的 $1$ 的邻域族，其中 $i_1,\ldots,i_t\in I$ 为任意有限多个指标，而 $\{1\}_i$ 表示 $G_i$ 中只含其单位元的子集. 由于诸 $G_i$ 是离散的，上述邻域族构成 $\prod_{i\in I}G_i$ 的单位元的邻域基本系（见乘积拓扑的定义——\cite[\S19]{mun75}）. 由于 $I$ 是有向集且所涉指标集有限，存在 $i_0\in I$ 使得 $i_0\geq i_r$ 对一切 $1\leq r\leq t$ 成立. 于是有
\[
G\cap\bigl[(\prod_{i\neq i_0}G_i)\times\{1\}_{i_0}\bigr]=G\cap\bigl[(\prod_{i\neq i_1,\ldots,i_t}G_i)\times\{1\}_{i_1}\times\ldots\times\{1\}_{i_t}\bigr]
\]
因为 $f_{i_ri_0}(1)=1$ 对一切 $1\leq r\leq t$ 成立.

这样，$G$ 中 $1$ 的开邻域基本系由全体形如 $G\cap[(\prod_{i\neq i_0}G_i)\times\{1\}_{i_0}]$ 的集合给出. 而我们看到 $\ker(\varphi_{i_0})=G\cap[(\prod_{i\neq i_0}G_i)\times\{1\}_{i_0}]$.
\end{proof}

\begin{lemma}\label{c3-lem-2-2}
设 $X$ 是紧致 Hausdorff 拓扑空间. 设 $x\in X$，而 $\{U_i\mid i\in I\}$ 是包含 $x$ 的全体紧开集之族. 则 $A=\cap_{i\in I}U_i$ 是连通的.
\end{lemma}

\begin{proof}
注意 $X$ 是 Hausdorff 空间且 $U_i$ 紧致，蕴含 $U_i$ 对一切 $i\in I$ 是闭集（见 \cite[Theorem~26.3]{mun75}）.

假设 $A$ 存在分离，即假设存在闭集 $U,V$ 使得 $A=U\cup V$ 且 $U\cap V=\emptyset$. 由于 $X$ 是正规空间（$X$ 紧致 Hausdorff——见 \cite[Theorem~32.3]{mun75}），存在开集 $U',V'$ 使得 $U\subseteq U'$，$V\subseteq V'$ 且 $U'\cap V'=\emptyset$. 于是有：
\[
U\cup V\subseteq U'\cup V'\Rightarrow X\setminus(U'\cap V')\subseteq X\setminus(U\cap V),
\]
即 $[X\setminus(U'\cap V')]\cap A=\emptyset$. 但 $X\setminus(U'\cap V')$ 是闭集，故也紧致. 由 $X\setminus(U'\cap V')\subseteq X\setminus A=\cup_{i\in I}X\setminus U_i$（诸 $X\setminus U_i$ 是开集），存在有限集 $J\subseteq I$ 使得
\[
[X\setminus(U'\cap V')]\cap(\bigcap_{i\in J}U_i)=\emptyset.
\]
考虑 $B:=\cap_{i\in J}U_i$. $B$ 是开集，因为 $|J|<\infty$ 且 $U_i$ 对一切 $i\in I$ 是开集；$B$ 是紧致的，因为它是 $X$ 的闭子集（$B$ 也是闭集之交）. 我们有 $x\in B$ 且 $B=(B\cap U')\cup(B\cap V')$. 假设 $x\in B\cap U'$. 由于 $B\cap U'$ 既开又紧，对某个 $i_0\in I$ 有 $B\cap U'=U_{i_0}$. 这意味着 $A\subseteq B\cap U'\subseteq U'$，于是
\[
A\cap V\subseteq A\cap V'\subseteq U'\cap V'=\emptyset\Rightarrow V=\emptyset,
\]
故结论得证.
\end{proof}

\begin{lemma}\label{c3-lem-2-3}
设 $G$ 是紧致、Hausdorff、全不连通的拓扑群. 则 $1$ 的每个邻域都含有一个开正规子群. 并且该子群在 $G$ 中指数有限.
\end{lemma}

\begin{proof}
考虑 $\{U_i\mid i\in I\}$：包含 $1$ 的全体紧开集之族. 如同前面的引理，一个拓扑学结果蕴含 $U_i$ 对每个 $i\in I$ 是闭集.

由\autoref{c3-lem-2-2} 以及 $X$ 全不连通\footnote{原文即如此作 $X$；按上下文此处应为 $G$.}，我们有 $\{1\}=\cap_{i\in I}U_i$.

设 $U$ 是 $\{1\}$ 的一个开邻域. 则 $G\setminus U$ 是闭集，故紧致. 我们有
\[
(G\setminus U)\cap(\bigcap_{i\in I}U_i)=\emptyset,
\]
这说明存在有限集 $J\subseteq I$ 使得
\[
(G\setminus U)\cap(\bigcap_{i\in J}U_i)=\emptyset.
\]
令 $A:=\cap_{i\in J}U_i$，则 $A$ 作为有限个开集之交是开集，且 $A$ 紧致，因为它是 $G$ 的闭子空间（注意 $A$ 是闭集之交）. 这意味着 $A$ 是 $1$ 的邻域. 此外 $A\subseteq U$.

定义 $F:=(G\setminus A)\cap A^2$. 注意 $A^2$ 紧致，因为它是 $A\times A$（紧致——见 \cite[Theorem~26.7]{mun75}）在连续映射 $(x,y)\mapsto xy$ 下的像. 于是 $A^2$ 是闭集，从而 $F$ 也是闭集.

设 $V$ 是 $1$ 的对称开邻域，使得 $AV\cap F=\emptyset$ 且 $V\subseteq A$（这样的邻域存在，由第 1 章的\autoref{c1-prop-2-4}）. 于是有 $AV\subseteq A^2\Rightarrow AV\cap(G\setminus A)\subseteq A^2\cap(G\setminus A)=F$. 由 $AV\cap F=\emptyset$ 可知 $AV\cap(G\setminus A)=\emptyset$. 从而 $AV\subseteq A\Rightarrow AV^n\subseteq A$ 对一切 $n\geq 1$ 成立，这就推出 $V^n\subseteq A$（因为 $1\in A$）.

定义 $K:=\cup_{n\geq 1}V^n\subseteq A$，它是 $G$ 的开子集. 由于 $G$ 紧致，由\autoref{c1-prop-4-4}，$K$ 指数有限. 于是可写 $G:=\cup_{m=1}^r x_mK$. $K$ 指数有限蕴含 $K$ 的共轭类只有有限多个（$xK=yK\Rightarrow$ 存在 $k\in K$ 使 $x=yk\Rightarrow xKx^{-1}=ykKk^{-1}y^{-1}=yKy^{-1}$）. 故
\[
H:=\bigcap_{x\in G}xKx^{-1}=\bigcap_{m=1}^r x_mKx_m^{-1}.
\]
由于 $x_mKx_m^{-1}$ 对一切 $1\leq m\leq r$ 是开集（见\autoref{c1-prop-4-4}），$H$ 是 $G$ 的开正规子群. 这也意味着 $H$ 指数有限（\autoref{c1-prop-4-4}）.
\end{proof}

\begin{lemma}\label{c3-lem-2-4}
设 $(X_i,f_{ij})_I$ 是拓扑空间的逆系统. 设 $(X,\varphi_i)_{i\in I}$ 是相容族，且一切 $h_i$ 都是满射.\footnote{原文即如此作 $h_i$；按上下文应为 $\varphi_i$.} 则或者 $\varprojlim X_i$ 为空，或者诱导映射 $\theta:X\to\varprojlim X_i$ 把 $X$ 映满 $\varprojlim X_i$ 的一个稠密子集.
\end{lemma}

\begin{proof}
我们将证明：对 $\varprojlim X_i$ 的每个基本开集 $V$，都可找到 $y\in X$ 使得 $\theta(y)\in V$.

考虑
\[
V=(\varprojlim X_i)\cap(\prod_{i\in I\setminus\{i_1,\ldots,i_n\}}X_i\times U_1\times\ldots U_n),
\]
其中对每个 $1\leq j\leq n$，$U_j$ 是 $X_{i_j}$ 的非空开子集.

由于 $I$ 是有向集，可找到 $i_0\in I$ 使得 $i_0>i_j$ 对一切 $1\leq j\leq n$ 成立. 若取 $x=\prod_{i\in I}X_i\in V$，则 $f_{i_ji_0}(x_{i_0})=x_{i_j}$ 对一切 $1\leq j\leq n$ 成立.

由于 $\varphi_{i_0}$ 是满射，存在 $y\in X$ 使得 $\varphi_{i_0}(y)=x_{i_0}$. 这意味着 $\pi_{i_0}(\theta(y))=x_{i_0}$，于是 $\pi_{x_j}(\theta(y))=f_{i_ji_0}(\pi_{x_0}(\theta(y)))=x_{i_j}$\footnote{原文即如此作 $\pi_{x_j}$ 与 $\pi_{x_0}$；按上下文应为 $\pi_{i_j}$ 与 $\pi_{i_0}$.}，从而 $\theta(y)\in V$.
\end{proof}

\subsection{投射有限群的其他刻画}

\begin{theorem}\label{c3-thm-2-5}
投射有限群恰是紧致、全不连通的 Hausdorff 群.
\end{theorem}

\begin{proof}
假设 $G$ 是投射有限群.

由投射有限群的定义，存在有限群 $(G_i)_{i\in I}$ 与映射 $f_{ij}:X_j\to X_i$\footnote{原文即如此；按上下文应为 $f_{ij}:G_j\to G_i$.}（对一切 $i,j\in I$，$i\leq j$）使得 $G=\varprojlim_{i\in I}G_i$. 由于一切 $G_i$ 都是有限离散拓扑群，它们都是 Hausdorff、紧致且全不连通的. 于是由\autoref{c3-prop-1-3}，$G$ 是 Hausdorff 紧致全不连通群.

反之，假设 $G$ 是 Hausdorff、紧致且全不连通的拓扑群.

当 $U$ 是 $G$ 的开正规子群时，记 $U<_O G$.

考虑族 $\mathcal{S}=\{U\mid U<_O G\}$. 由于 $G$ 紧致，$U$ 为开集蕴含 $G/U$ 有限（第 1 章的\autoref{c1-prop-4-4}）. 由\autoref{c3-lem-2-3}，$\mathcal{S}$ 是 $1\in G$ 的开邻域基.

对每一对满足 $U\subseteq V$ 的 $U,V\in\mathcal{S}$，考虑自然映射：
\[
f_{U,V}:G/U\to G/V:gU\mapsto gV.
\]
于是显然 $(G/U,f_{U,V})_{\mathcal{S}}$ 是群的逆系统.

考虑典范同态构成的相容族 $\varphi_U:G\to G/U$，$U\in\mathcal{S}$，我们得到一个映射
\[
\varphi:G\to s\varprojlim_{\mathcal{S}}G/U\subseteq\prod_{U\in\mathcal{S}}G/U.
\]
我们将证明 $\varphi$ 是单射、连续且满射. 然后由 $G$ 紧致即得 $\varphi$ 是拓扑同构（见 \cite[Theorem~26.6]{mun75}）.
\begin{itemize}
    \item \textbf{$\varphi$ 是单射：}设 $\sigma\in G$ 满足 $\varphi(\sigma)=1$. 则 $\sigma\in U$ 对一切 $U\in\mathcal{S}$ 成立，从而 $\sigma\in\cap_{U\in\mathcal{S}}U=\{1\}$.
    \item \textbf{$\varphi$ 是连续的：}我们证明复合
    \[
    G\xrightarrow{\varphi}s\varprojlim_{\mathcal{S}}G/U\xrightarrow{\pi_U}G/U
    \]
    对每个 $U\in\mathcal{S}$ 连续. 事实上这就够了，因为 $s\varprojlim_{\mathcal{S}}G/U\subseteq\prod_{U\in\mathcal{S}}G/U$（见 \cite[Theorem~19.6]{mun75}）. 而
    \[
    (\pi_U\varphi)^{-1}\{1\}=\varphi^{-1}(\pi_U^{-1}(\{1\}))=U
    \]
    在 $G$ 中是开集.

    （注意对一切 $U\in\mathcal{S}$，$\{1\}$ 是 $G/U$ 中 $1$ 的开邻域基——因为 $\{1\}$ 在 $G$ 中本来就是开邻域基.）
    \item \textbf{$\varphi$ 是满射：}由\autoref{c3-lem-2-4}，$\varphi(G)$ 在 $s\varprojlim_{\mathcal{S}}G/U$ 中稠密. 我们证明 $\varphi(G)$ 是闭集，从而 $\varphi(G)=s\varprojlim_{\mathcal{S}}G/U$. $\varphi(G)$ 是闭集，因为 $\varphi(G)=\pi_U^{-1}(\varphi_U(G))=\pi_U^{-1}(G/U)$（这里用到 $\varphi_U$ 是满射）而 $\pi_U$ 连续.
\end{itemize}
\end{proof}

\begin{remark}\label{c3-rem-2-1}
一些作者依据\autoref{c3-thm-2-5}，把投射有限群定义为 Hausdorff、紧致、全不连通的拓扑群.
\end{remark}

\begin{corollary}\label{c3-cor-2-6}
设 $G$ 是投射有限群. 令 $\mathcal{S}:=\{G\ \text{的开正规子群}\}$，则有
\[
G\cong{\varprojlim_{U\in\mathcal{S}}}G/U.
\]
\end{corollary}

\subsection{定理的推论}

\begin{proposition}\label{c3-prop-2-7}
设 $G$ 是投射有限群.
\begin{enumerate}[label=(\arabic*),leftmargin=*]
    \item 设 $\{H_i\mid i\in I\}$ 是 $G$ 的一族闭子群，且 $\cap_{i\in I}H_i\leq U\leq G$，其中 $U$ 是 $G$ 的开子群. 则存在 $I$ 的有限子集 $J$ 使得 $\cap_{j\in J}H_j\leq U$.
    \item 设 $\{U_i\mid i\in I\}$ 是 $G$ 的一族开子群，且 $\cap_{i\in I}U_i=1$. 令 $\mathcal{V}=\{\cap_{j\in J}U_j\mid J\subseteq I\ \text{有限}\}$. 则 $\mathcal{V}$ 是 $G$ 中 $1$ 的邻域基本系.
\end{enumerate}
\end{proposition}

\begin{proof}\
\begin{enumerate}[label=(\arabic*),leftmargin=*]
    \item 考虑紧空间 $G\setminus U$ 的开覆盖 $\{G\setminus H_i\mid i\in I\}$（注意 $G\setminus U$ 是紧空间 $G$ 的闭子集）. 于是存在有限子覆盖 $\{G\setminus H_j\mid j\in J\}$，其中 $J$ 是 $I$ 的有限子集. 这意味着
    \[
    G\setminus U\subseteq\bigcup_{j\in J}G\setminus H_j,
    \]
    从而
    \[
    \bigcap_{j\in J}G\setminus H_j\subseteq U.
    \]
    \item 由 (1) 即得.
\end{enumerate}
\end{proof}

\begin{proposition}\label{c3-prop-2-8}
投射有限群 $G$ 的闭子群 $H$ 是投射有限的. 确切地说，若 $G=\varprojlim_{\mathcal{S}}G/U$，其中 $\mathcal{S}=\{G\ \text{的开正规子群}\}$，则
\[
H\cong{\varprojlim_{\mathcal{S}}}HU/U\cong{\varprojlim_{\mathcal{S}}}H/H\cap U.
\]
\end{proposition}

\begin{proof}
我们使用\autoref{c3-thm-2-5} 的刻画. 显然 $H$ 是 Hausdorff 且全不连通的，因为 $G$ 如此. 定理 \cite[Theorem~26.2, Chapter~3]{mun75} 蕴含 $H$ 紧致（因为 $H$ 是闭集）.

第二个断言由\autoref{c3-lem-2-4} 得出. 事实上，典范映射
\[
\begin{aligned}
\varphi_U&:H\to HU/U\\
\pi_U&:H\to H/H\cap U
\end{aligned}
\]
对 $\mathcal{S}$ 中一切 $U$ 都是满射.
\end{proof}

\begin{proposition}\label{c3-prop-2-9}
投射有限群 $G$ 模闭子群 $H$ 的商群 $G/H$ 是投射有限群. 事实上，若以 $\mathcal{S}$ 记 $G$ 的一切开正规子群之集，则
\[
G/H\cong{\varprojlim_{\mathcal{S}}}G/HU.
\]
\end{proposition}

\begin{proof}
考虑典范投影
\[
\pi:G\to G/HU,
\]
它是满射. 由于 $\pi(H)=0$\footnote{原文即如此；按上下文应为 $\pi(H)=1$.}，我们可以定义
\[
\widetilde{\pi}:G/H\to G/HU
\]
并再次使用\autoref{c3-lem-2-4}.
\end{proof}

\begin{proposition}\label{c3-prop-2-10}
设 $I$ 是指标集，$(G_i)_{i\in I}$ 是一族投射有限群. 则 $G:=\prod_{i\in I}G_i$ 是投射有限群.
\end{proposition}

\begin{proof}
显然：Hausdorff、紧致、全不连通群的乘积仍是 Hausdorff、紧致、全不连通群.
\end{proof}

\begin{corollary}\label{c3-cor-2-11}
设 $(G_i,f_{ij})_I$ 是投射有限群的逆系统. 则 $G=\varprojlim_I G_i$ 是投射有限的.
\end{corollary}

\begin{proof}
用\autoref{c3-prop-1-3} 可见 $G$ 是 $\prod_{i\in I}G_i$ 的闭子群，再用\autoref{c3-cor-2-11}\footnote{原文即如此（此处自指推论 2.11）；按上下文应为 命题 2.8.}.
\end{proof}

\begin{example}\label{c3-ex-2-2}
设 $\mathbb{P}$ 为全体素数之集. 我们将证明
\[
\widehat{\mathbb{Z}}\cong\prod_{p\in\mathbb{P}}\mathbb{Z}_p.
\]
我们构造自然投影
\[
\alpha^m_p:\mathbb{Z}_p\to\mathbb{Z}/p^{m_p}\mathbb{Z}
\]
其中 $m=\prod_{q\in\mathbb{P}}q^{m_q}$. 于是有映射
\[
\alpha^m:\prod_{p\in\mathbb{P}}\alpha^m_p:\prod_{p\in\mathbb{P}}\mathbb{Z}_p\to\prod_{p\in\mathbb{P}}\mathbb{Z}/p^{m_p}\mathbb{Z},
\]
它是满射（可以容易地构造出一个原像）.

注意，利用中国剩余定理（见 \cite[Corollary~2.2, \S2, Chapter~2]{lang02a}），我们有
\[
\prod_{p\in\mathbb{P}}\mathbb{Z}/p^{m_p}\mathbb{Z}\cong\mathbb{Z}/m\mathbb{Z}.
\]
利用\autoref{c3-lem-2-4}，映射 $\alpha^m$（$m\in\mathbb{Z}$）诱导出一个连续满射
\[
\alpha:\prod_{p\in\mathbb{P}}\mathbb{Z}_p\to\widehat{\mathbb{Z}}={\varprojlim_{m}}\mathbb{Z}/m\mathbb{Z}.
\]
另一方面，设 $(x_p)_{p\in\mathbb{P}}$ 是 $\prod_{p\in\mathbb{P}}\mathbb{Z}_p$ 中满足 $\alpha((x_p)_\mathbb{P})=0$ 的元素. 这意味着 $\alpha^m((x_p)_\mathbb{P})=0$ 对每个 $m\in\mathbb{Z}$ 成立（由 $\widehat{\mathbb{Z}}$ 的定义）. 从而 $\alpha^m_q((x_p)_\mathbb{P})=0$ 对每个 $m\in\mathbb{Z}$ 与每个 $q\in\mathbb{P}$ 成立. 于是对一切 $p\in\mathbb{P}$ 有 $x_p=0$，即 $(x_p)_\mathbb{P}=0$. 这说明 $\alpha$ 是单射.

由于 $\prod_{p\in\mathbb{P}}\mathbb{Z}_p$ 紧致，我们得到 $\alpha$ 是拓扑同构（见 \cite[Theorem~26.6, Chapter~3]{mun75}）.
\end{example}

\begin{proposition}\label{c3-prop-2-12}
$\varprojlim$ 是从投射有限群的逆系统范畴（指标集 $I$ 为有向集）到投射有限群范畴的正合函子.
\end{proposition}

\begin{proof}
设
\[
1\longrightarrow(H_i,\varphi_{ij})\xrightarrow{(f_i)_I}(G_i,\psi_{ij})\xrightarrow{(g_i)_I}(K_i,\eta_{ij})\longrightarrow 1
\]
是投射有限群逆系统（在 $I$ 上）的一个正合列. 我们须证
\[
1\longrightarrow{\varprojlim_I}\,H_i\xrightarrow{f}{\varprojlim_I}\,G_i\xrightarrow{g}{\varprojlim_I}\,K_i\longrightarrow 1
\]
正合.
\begin{itemize}
    \item 先证 $f$ 是单射：设 $(x_i)_I\in{\varprojlim_I}\,H_i$ 满足 $f((x_i)_I)=0$. 则对一切 $i\in I$ 有 $f_i(x_i)=0$. 而 $f_i$ 是单射，故对一切 $i\in I$ 有 $x_i=0$. 于是 $(x_i)_I=0$.
    \item 其次，证明 $\ker(g)=\mathrm{im}(f)$：设 $(x_i)_I\in{\varprojlim_I}\,H_i$，
    \[
    (g\circ f)((x_i)_I)=g((f_i(x_i))_I)=((g_i\circ f_i)(x_i))_I=0
    \]
    因为逆系统序列是正合的. 这说明 $\mathrm{im}(f)\subseteq\ker(g)$.

    设 $(y_i)_I\in\ker(g)$. 这等价于 $g((y_i)_I)=0$，于是 $(g_i(y_i))_I=0$. 从而对一切 $i\in I$ 有 $y_i\in\ker(g_i)$. 而逆系统序列是正合的，故对每个 $i\in I$，存在 $x_i\in H_i$ 使得 $f_i(x_i)=y_i$. 于是 $\ker(g)\subseteq\mathrm{im}(f)$.
    \item 最后须证 $g$ 是满射：设 $(z_i)_I\in{\varprojlim_I}\,K_i$. 考虑集合 $Y_i=g_i^{-1}(x_i)$\footnote{原文即如此作 $x_i$；前文所设为 $(z_i)_I$，按上下文应为 $g_i^{-1}(z_i)$，下文诸 $x_i$ 同此.}，它是 $G_i$ 中的紧集，因为它是闭集的原像. 注意 $(g_i)_I$ 是逆系统范畴中的态射，这意味着当 $i\leq j$ 时下图交换\footnote{交换图上方箭头原文即如此标 $\varphi_{ij}$；按上下文应为 $\psi_{ij}$.}：
    \[
    \begin{tikzcd}
    G_j \arrow[r, "\varphi_{ij}"] \arrow[d, "g_j"'] & G_i \arrow[d, "g_i"] \\
    K_j \arrow[r, "\eta_{ij}"'] & K_i
    \end{tikzcd}
    \]
    于是当 $i\leq j$ 时，$\psi_{ij}(Y_j)\subseteq Y_i$，事实上对每个 $y_j\in Y_j$ 有
    \[
    \begin{aligned}
    g_i(\psi_{ij}(y_j))\,&=(g_i\circ\psi_{ij})(y_j)\\
    &=(\eta_{ij}\circ g_j)(y_j)\\
    &=\eta_{ij}(x_j)=x_i,
    \end{aligned}
    \]
    从而 $(Y_i,\psi_{ij})$ 是一族非空紧集的逆系统. 故 ${\varprojlim_I}\,Y_i\neq\emptyset$（见 \cite[Proposition~8, \S9, Chapter~1]{bou71}）. 而对每个 $(y_i)_I\in{\varprojlim_I}\,Y_i$，有 $g((y_i)_I)=(x_i)_I$，故结论得证.
\end{itemize}
\end{proof}

\section{Galois 群与投射有限群}

我们已在第 2 章（\autoref{c2-thm-2-6}）看到，Galois 群是 Hausdorff、紧致且全不连通的. 而由\autoref{c3-thm-2-5}，投射有限群亦然. 本节将把这一点说得更清楚.

\begin{theorem}\label{c3-thm-3-1}
假设 $\mathbb{L}$ 是域 $\mathbb{K}$ 的 Galois 扩张. 固定 $\mathcal{F}=\mathcal{F}=\{\mathbb{E}\mid \mathbb{K}\subseteq\mathbb{E}\subseteq\mathbb{L}$，$\mathbb{E}$ 为 $\mathbb{K}$ 的有限 Galois 扩张$\}$\footnote{原文即如此（$\mathcal{F}=\mathcal{F}$ 重复，疑为衍文）.}. 则 $\mathrm{Gal}(\mathbb{L},\mathbb{K})$ 是有限群 $\mathrm{Gal}(\mathbb{E},\mathbb{K})$（$\mathbb{E}\in\mathcal{F}$）的逆极限；特别地，$\mathrm{Gal}(\mathbb{L},\mathbb{K})$ 是投射有限群.
\end{theorem}

\begin{proof}
由 Galois 理论，$\mathrm{Gal}(\mathbb{E},\mathbb{K})$ 是有限群. 若 $\mathbb{E}_1,\mathbb{E}_2\in\mathcal{F}$ 且 $\mathbb{E}_1\subseteq\mathbb{E}_2$，我们定义
\[
\begin{aligned}
\varphi_{\mathbb{E}_1\mathbb{E}_2}:\mathrm{Gal}(\mathbb{E}_2,\mathbb{K})&\longrightarrow\mathrm{Gal}(\mathbb{E}_1,\mathbb{K})\\
\sigma&\longmapsto\sigma|_{\mathbb{E}_1}.
\end{aligned}
\]
于是 $(\mathrm{Gal}(\mathbb{E},\mathbb{K}),\varphi_{\mathbb{E}_1\mathbb{E}_2})_{\mathcal{F}}$ 是以 $\mathcal{F}$ 为指标的逆系统.

考虑限制映射 $\mathrm{Gal}(\mathbb{L},\mathbb{K})\to\mathrm{Gal}(\mathbb{E},\mathbb{K})$，我们得到群同态 $\theta:\mathrm{Gal}(\mathbb{L},\mathbb{K})\to\prod_{\mathbb{E}\in\mathcal{F}}\mathrm{Gal}(\mathbb{E},\mathbb{K})$. $\theta$ 的像含于 $\varprojlim_{\mathbb{E}\in\mathcal{F}}\mathrm{Gal}(\mathbb{E},\mathbb{K})$，因为
\[
\varphi_{\mathbb{E}_1,\mathbb{E}_2}(\pi_{\mathbb{E}_2}(\theta(\sigma)))=\varphi_{\mathbb{E}_1,\mathbb{E}_2}(\sigma|_{\mathbb{E}_2})=\sigma|_{\mathbb{E}_1}=\pi_{\mathbb{E}_1}(\theta(\sigma)).
\]
现在给定 $(\sigma_\mathbb{E})_{\mathbb{E}\in\mathcal{F}}\in s\varprojlim_{\mathbb{E}\in\mathcal{F}}\mathrm{Gal}(\mathbb{E},\mathbb{K})$ 与 $x\in\mathbb{K}$，我们构造
\[
\psi((\sigma_\mathbb{E}))(x)=\sigma_\mathbb{M}(x)
\]
其中 $\mathbb{M}\in\mathcal{F}$ 是满足 $x\in\mathbb{M}$ 的一个中间扩张（这样的 $\mathbb{M}$ 总是存在）. $\psi$ 的这一构造与 $\mathbb{M}$ 的选取无关，因为 $\sigma$ 属于 $\varprojlim_{\mathbb{E}\in\mathcal{F}}\mathrm{Gal}(\mathbb{E},\mathbb{K})$. 容易看出 $\psi((\sigma_\mathbb{E}))\in\mathrm{Gal}(\mathbb{L},\mathbb{K})$ 且 $\psi$ 是 $\theta$ 的逆映射. 于是 $\theta$ 是同构. 从而作为抽象群，$\mathrm{Gal}(\mathbb{L},\mathbb{K})$ 同构于 $\varprojlim_{\mathbb{E}\in\mathcal{F}}\mathrm{Gal}(\mathbb{E},\mathbb{K})$.

下面我们将证明 $1$ 处的拓扑基之间存在同构. 回忆 $\varprojlim_{\mathbb{E}\in\mathcal{F}}\mathrm{Gal}(\mathbb{E},\mathbb{K})$ 中 $1$ 的开邻域基由
\[
\{\ker(\pi_\mathbb{E})\mid \mathbb{E}\in\mathcal{F}\}
\]
给出（见\autoref{c2-thm-3-3}）\footnote{原文作 "proposition 3.3"；本章并无命题 3.3，按上下文疑指第 2 章定理 3.3（Galois 群的 Krull 拓扑）.}. 但对每个 $\mathbb{F}\in\mathcal{F}$，我们有
\[
\ker(\pi_\mathbb{F})=\{(\sigma_\mathbb{E})\in s\varprojlim_{\mathbb{E}\in\mathcal{F}}\mathrm{Gal}(\mathbb{E},\mathbb{K})\mid \pi_\mathbb{F}((\sigma_\mathbb{E}))=\mathrm{Id}_\mathbb{F}\}
\]
注意
\[
\theta(\mathrm{Gal}(\mathbb{L},\mathbb{F}))=\prod_{\mathbb{E}\in\mathcal{F}\setminus\{\mathbb{F}\}}\mathrm{Gal}(\mathbb{E},\mathbb{K})\times\{\mathrm{Id}_\mathbb{F}\}
\]
且由 $\mathrm{im}(\theta)\subseteq s\varprojlim_{\mathbb{E}\in\mathcal{F}}\mathrm{Gal}(\mathbb{E},\mathbb{K})$，我们得到
\[
\ker(\pi_\mathbb{F})=\theta(\mathrm{Gal}(\mathbb{L},\mathbb{F}))
\]
对一切 $\mathbb{F}\in\mathcal{F}$ 成立. 这就证明了同胚，于是
\[
\mathrm{Gal}(\mathbb{L},\mathbb{K})\cong s\varprojlim_{\mathbb{E}\in\mathcal{F}}\mathrm{Gal}(\mathbb{E},\mathbb{K}).
\]
\end{proof}

\begin{proof}
你可以在 \cite[Chapter~3.3]{wil98} 中找到该引理的证明.\footnote{原文即如此：此处并无对应的引理陈述，此证明块疑为排版残留.}
\end{proof}

\begin{theorem}\label{c3-thm-3-2}
每个投射有限群 $G$ 都同构（作为拓扑群）于一个 Galois 群.
\end{theorem}

\begin{proof}
回忆 $G$ 是 Hausdorff、紧致、全不连通的拓扑群（见\autoref{c3-thm-2-5}）.

设 $\mathbb{F}$ 是任意域. 我们考虑 $S$：集合族 $G/N$ 的不交并，其中 $N$ 遍历 $G$ 的全体开正规子群（我们记 $N<_O G$）. 把 $S$ 的元素看作未定元，并构造域 $\mathbb{L}=\mathbb{F}(S)$（以 $\mathbb{F}$ 中元素为系数、以 $S$ 中元素为未定元的有理函数域）.

群 $G$ 以自然方式作用于 $S$：若 $g\in G$ 且 $g'N\in G/N$，则 $g(g'N)=gg'N$（对任何 $N\in S$）. 于是这诱导出 $G$ 在 $\mathbb{L}$ 上的一个作用.

定义 $\mathbb{K}=\mathbb{L}^G$ 为由 $G$ 固定的 $\mathbb{K}$ 的元素构成的 $\mathbb{L}$ 的子域.\footnote{原文即如此（"the elements of $\mathbb{K}$ fixed by $G$"）；按上下文应为 $\mathbb{L}$ 中由 $G$ 固定的元素.}

我们现在的目标是证明 $\mathbb{L}$ 是 $\mathbb{K}$ 的 Galois 扩张且 Galois 群为 $G$.

若 $\alpha\in\mathbb{L}$，考虑 $G$ 的子群
\[
G_\alpha=\{\sigma\in G\mid \sigma(\alpha)=\alpha\}.
\]
若出现在 $\alpha$ 的有理表达式中的未定元为 $\{t_i\in G/N_i\mid i=1,\ldots,n\}$，则
\[
\{1\}\subseteq\bigcap_{i=1}^n N_i\subseteq G_\alpha.
\]
事实上，若 $n\in N_i$ 且 $gN_i\in G/N_i$，我们有
\[
n(gN_i)=ngN_i=gg^{-1}ngN_i=gN_i,
\]
因为 $g^{-1}ng\in N_i$（$N_i$ 是 $G$ 的正规子群）. 这意味着 $N_i$ 在 $G/N_i$ 上的作用是平凡的. 由于一切 $N_i$ 都是 $G$ 的开子群，$\cap_{i=1}^n N_i$ 是开集. 于是利用第 1 章的\autoref{c1-prop-4-4} 第 (4) 部分可知 $G_\alpha$ 是开集. 从而，利用同一命题（第 (2) 部分），$G_\alpha$ 指数有限.

由于元素的稳定子群的指数等于其轨道的基数，$\alpha$ 在 $G$ 作用下的轨道是有限集. 设该轨道为 $O_\alpha=\{\alpha_1,\ldots,\alpha_r\}$（诸 $\alpha_i$ 两两不同），并考虑多项式\footnote{原文即如此；按上下文应为 $\prod_{i=1}^r(X-\alpha_i)$.}
\[
f(X)=\prod_{i=1}^r(X\alpha_i).
\]
由于 $G$ 把 $O_\alpha$ 变为自身（由轨道的定义），对一切 $1\leq i\leq r$，$\alpha_i$ 属于 $\mathbb{K}$，这意味着 $f(X)\in\mathbb{K}[X]$. 而 $\alpha$ 作为 $f(X)$ 的根，在 $\mathbb{K}$ 上是代数元. 此外，注意 $f(X)$ 的一切根两两不同，这意味着 $\alpha$ 在 $\mathbb{K}$ 上是可分的. 于是 $\mathbb{L}$ 是 $\mathbb{K}$ 的可分扩张.

注意 $\mathbb{K}[\alpha_1,\ldots,\alpha_r]$ 是 $\mathbb{K}$ 的正规扩张（作为非常数多项式 $f(X)\in\mathbb{K}[X]$ 的分裂域）. 于是 $\mathbb{L}$ 是 $\mathbb{K}$ 的一族正规扩张之并，故 $\mathbb{L}$ 在 $\mathbb{K}$ 上正规.

这样我们就证明了 $\mathbb{L}$ 是 $\mathbb{K}$ 的 Galois 扩张.

我们记 $H:=\mathrm{Gal}(\mathbb{L},\mathbb{K})$. 显然，由于 $G$ 作用于 $\mathbb{L}$ 且固定 $\mathbb{K}=\mathbb{L}^G$，有 $G\subseteq H$. 下面证明等号成立. 考虑包含映射 $i:G\hookrightarrow H$. 下一步将证明此映射连续.

设 $U$ 是 $H$ 的正规开子群，并考虑 $\mathbb{L}^U$. 由于 $U$ 指数有限（见第 1 章的\autoref{c1-prop-4-4}），$\mathbb{L}^U$ 是 $\mathbb{K}$ 的有限 Galois 扩张（用第 2 章的\autoref{c2-thm-3-3}）. 设 $\mathbb{L}^U=\mathbb{K}(\alpha'_1,\ldots,\alpha'_s)$，其中 $\alpha'_i\in\mathbb{L}$. 于是
\[
G\cap U\supseteq\bigcap_{i=1}^s G_{\alpha'_i}.
\]
这说明 $G\cap U$ 是开集（因为它含有一个开子群——见第 1 章的\autoref{c1-prop-4-4}）. 于是 $G$ 是开集，从而 $G$ 是 $H$ 的闭子群.

最后，利用关于无限 Galois 扩张的定理（第 2 章的\autoref{c2-thm-3-3}）：由 $\mathbb{L}^G=\mathbb{L}^H$ 且 $G$ 是闭集，得 $H=G$.
\end{proof}

% !TeX root = ../main.tex
% 第 4 章 群的上同调 （原书 CHAPTER 4, Group Cohomology, 印刷页 35-41）

\chapter{群的上同调}
\label{ch:4}

下一章（按原书行文如此\footnote{原文为 ``In the following chapter''，按上下文即指本章.}）我们定义\newterm{c4-cohomology}{上同调群}{cohomology group}.这一概念源于拓扑学，后来被用于群论中以提供不变量.\\
本章许多证明留给读者.

\section{概说}

\subsection{（余）同调模}

\begin{definition}\label{c4-def-1-1}
一个\emph{左 $R$-模}（其中 $R$ 是环）是指配有标量乘法 $R\times M\to M:(r,m)\mapsto r\cdot m$ 的阿贝尔群 $(M,+)$，使得对一切 $m,m'\in M$ 与一切 $r,r'\in R$，都有
\begin{enumerate}[label=(\arabic*),leftmargin=*]
    \item $r\cdot(m+m')=r\cdot m+r\cdot m'$，
    \item $(r+r')\cdot m=r\cdot m+r'\cdot m$，
    \item $(rr')\cdot m=r\cdot(r'\cdot m)$，
    \item $1\cdot m=m$.
\end{enumerate}
\end{definition}

\begin{definition}\label{c4-def-1-2}
设 $M$ 与 $N$ 是两个 $R$-模，一个 $R$-\emph{映射}（或 $R$-\emph{同态}）是指群同态 $f:M\to N$，使得对一切 $m\in M$ 与 $r\in R$，
\[
r\cdot f(m)=f(r\cdot m).
\]
\end{definition}

可以类似地定义右 $R$-模，但当 $R$ 交换时两者一致.因此，为简化起见，从此刻起我们假定 $R$ 是一个交换环.

\begin{notation}\label{c4-not-1-3}
\begin{enumerate}[label=(\arabic*),leftmargin=*]
    \item 表示 $R$ 在 $R$-模 $M$ 上的作用时，我们常写 $rm$ 而不写 $r\circ m$.
    \item $R$-模与 $R$-映射构成的类是一个阿贝尔范畴，记作 $\mathscr{M}_R$.
\end{enumerate}
\end{notation}

这里我们介绍（余）同调的最一般概念.

\begin{definition}\label{c4-def-1-4}
$R$ 上的一个\newterm{c4-chain-complex}{链复形}{chain complex} $C$ 是一列 $R$-模 $(C_n)_{n\in\mathbb Z}$ 与 $R$-映射 $(d_n:C_n\to C_{n-1})_{n\in\mathbb Z}$，称为\newterm{c4-coboundary}{（余）边缘算子}{(co)boundary}\footnote{原文此处用词为 differentiations（微分）.}，
\[
C:\ldots\to C_{n+1}\xrightarrow{d_{n+1}}C_n\xrightarrow{d_n}C_{n-1}\to\ldots
\]
使得对一切 $n\in\mathbb Z$ 有 $d_n\circ d_{n+1}=0$.
\end{definition}

\begin{definition}\label{c4-def-1-5}
链复形的一个\emph{态射} $f:C\to C'$ 是一族 $R$-模的态射 $(f_n:C_n\to C_n')_{n\in\mathbb Z}$，使得对一切 $n$，下图交换
\[
\begin{tikzcd}
C_n \arrow[d,"f_n"'] \arrow[r,"d_n"] & C_{n-1} \arrow[d,"f_{n-1}"] \\
C_n' \arrow[r,"d_n'"'] & C_{n-1}'.
\end{tikzcd}
\]
\end{definition}

\begin{definition}\label{c4-def-1-6}
设 $C$ 是一个链复形，我们定义其第 $n$ 个\newterm{c4-homology}{同调群}{homology}\footnote{原文作 homology module（同调模）.}为
\[
H_n(C):=\ker d_n/\operatorname{im} d_{n+1}.
\]
\end{definition}

若 $f:C\to C'$ 是链复形的态射，我们定义
\[
H_n(f):H_n(C)\to H_n(C'):[z]\mapsto[f_n(z)].
\]
下面是上同调的定义.

\begin{definition}\label{c4-def-1-7}
$R$ 上的一个\emph{上链复形} $C$ 是一列 $R$-模 $(D_n)_{n\in\mathbb Z}$ 与 $R$-映射 $(\delta_n:C_n\to C_{n-1})_{n\in\mathbb Z}$\footnote{原文此处记号即如此不自洽：模的序列写作 $(D_n)$，映射写作 $\delta_n:C_n\to C_{n-1}$，而下面的图中模却写作 $D^{n-1}$、$D^n$、$D^{n+1}$，箭头标签亦上、下标混用（$\delta^{n-1}$ 与 $\delta_n$），此处一律照录.}，
\[
D:\ldots\to D^{n-1}\xrightarrow{\delta^{n-1}}D^n\xrightarrow{\delta_n}D^{n+1}\to\ldots
\]
使得对一切 $n\in\mathbb Z$ 有 $\delta^{n+1}\circ\delta^n=0$.
\end{definition}

\begin{definition}\label{c4-def-1-8}
设 $D$ 是一个上链复形，我们定义其第 $n$ 个上同调模为
\[
H^n(D):=\ker\delta^n/\operatorname{im}\delta^{n-1}.
\]
\end{definition}

\subsection{Hom 函子}

\begin{definition}\label{c4-def-1-9}
设 $\mathscr C$ 与 $\mathscr D$ 是范畴.从 $\mathscr C$ 到 $\mathscr D$ 的一个\emph{协变函子} $F$ 是指如下对应：它给每个对象 $C\in\mathscr C$ 配备一个对象 $F(C)\in\mathscr D$，并给 $\mathscr C$ 中每个态射 $f:C_1\to C_2$ 配备一个 $\mathscr D$ 中的态射 $F(f):F(C_1)\to F(C_2)$，使得
\begin{itemize}
    \item $F(\mathrm{id}_C)=\mathrm{id}_{F(C)}$，对每个对象 $C\in\mathscr C$；
    \item $F(g\circ f)=F(g)\circ F(f)$，对 $\mathscr C$ 中一切态射 $f:C_1\to C_2$ 与 $g:C_2\to C_3$.
\end{itemize}
从 $\mathscr C$ 到 $\mathscr D$ 的一个\emph{逆变函子} $F$ 是指如下对应：它给每个对象 $C\in\mathscr C$ 配备一个对象 $F(C)\in\mathscr D$，并给 $\mathscr C$ 中每个态射 $f:C_1\to C_2$ 配备一个 $\mathscr D$ 中的态射 $F(f):F(C_2)\to F(C_1)$，使得
\begin{itemize}
    \item $F(\mathrm{id}_C)=\mathrm{id}_{F(C)}$，对每个对象 $C\in\mathscr C$；
    \item $F(g\circ f)=F(f)\circ F(g)$，对 $\mathscr C$ 中一切态射 $f:C_1\to C_2$ 与 $g:C_2\to C_3$.
\end{itemize}
\end{definition}

\begin{definition}\label{c4-def-1-10}
设 $M$ 是一个固定的 $R$-模.\\
我们定义函子 $\operatorname{Hom}(\square, M)$：它从范畴 $\mathscr{M}_R$ 到范畴 $\mathscr A$，把每个 $R$-模 $N$ 送到从 $N$ 到 $M$ 的一切 $R$-映射的集合 $\operatorname{Hom}_R(N, M)$，并把每个 $R$-映射 $\varphi:N_1\to N_2$ 送到映射\footnote{原文即如此；按上下文应写作 $\operatorname{Hom}(\varphi, M)$ 与 $g\mapsto g\circ\varphi$.}
\[
\operatorname{Hom}(f, M):\operatorname{Hom}(N_2, M)\to\operatorname{Hom}(N_1, M):f\mapsto f\circ\varphi.
\]
其次，我们定义函子 $\operatorname{Hom}(M, \square)$：它从范畴 $\mathscr{M}_R$ 到范畴 $\mathscr A$，把每个 $R$-模 $N$ 送到从 $M$ 到 $N$ 的一切 $R$-映射的集合 $\operatorname{Hom}_R(M, N)$，并把每个 $R$-映射 $\varphi:N_1\to N_2$ 送到映射
\[
\operatorname{Hom}(f, M):\operatorname{Hom}(M, N_1)\to\operatorname{Hom}(M, N_2):f\mapsto\varphi\circ f.
\]
\end{definition}

\begin{proposition}\label{c4-prop-1-1}
\mbox{}
\begin{itemize}
    \item $\operatorname{Hom}_R(\square, M)$ 是从 $\mathscr{M}_R$ 到 $\mathscr A$ 的一个定义良好的逆变函子.此外，该函子是左正合的.
    \item $\operatorname{Hom}_R(M, \square)$ 是从 $\mathscr{M}_R$ 到 $\mathscr A$ 的一个定义良好的协变函子.此外，该函子是左正合的.
\end{itemize}
\end{proposition}

\begin{proof}
证明可参见 \cite[\S2.1]{rot09}.
\end{proof}

\begin{remark}\label{c4-rem-1-11}
从链复形 $C$ 出发得到上链复形有许多方式.最自然的是定义 $C^n:=C_{-n}$ 且 $\delta_m:=d_{-n}$\footnote{原文即如此；$\delta_m$ 的下标与右端 $d_{-n}$ 不相符，按上下文应为 $\delta_n$.}.\\
但另一种方式更有意思：我们只需施加函子 $\operatorname{Hom}(\square, G)$. 确切地说，$C^n:=\operatorname{Hom}_R(C_n, R)$ 且 $\delta^n:=d^*_{n+1}=\operatorname{Hom}_R(d_{n+1}, R)$. 我们把这一上链复形记作
\[
\operatorname{Hom}_R(C, R):\ldots\to\operatorname{Hom}_R(C_{n-1}, R)\xrightarrow{\delta^{n-1}}\operatorname{Hom}_R(C_n, R)\xrightarrow{\delta_n}\operatorname{Hom}_R(C_{n+1}, R)\to\ldots
\]
\end{remark}

下面的两个定理都是经典结果.

\begin{theorem}\label{c4-thm-1-2}
设 $0\to C'\xrightarrow{i}C\xrightarrow{p}C''\to 0$ 是链复形的一个\newterm{c4-exact-sequence}{正合列}{exact sequence}，则存在模的长正合列
\[
\ldots\to H_n(C')\xrightarrow{H_n(i)}H_n(C)\xrightarrow{H_n(p)}H_n(C'')\xrightarrow{\delta_n}H_{n-1}(C')\to\ldots.
\]
\end{theorem}

\begin{proof}
该定理的证明可参见 \cite[\S6.1]{rot09}.
\end{proof}

\begin{theorem}\label{c4-thm-1-3}
设 $0\to C'\xrightarrow{i}C\xrightarrow{p}C''\to 0$ 是上链复形的一个正合列，则存在模的长正合列
\[
\ldots\to H^n(C')\xrightarrow{H^n(i)}H^n(C)\xrightarrow{H^n(p)}H^n(C'')\xrightarrow{\delta^n}H^{n+1}(C')\to\ldots.
\]
\end{theorem}

\section{投射模与内射模}

本节我们定义投射模与内射模的概念.

\begin{definition}\label{c4-def-2-1}
称 $R$-模 $P$ 是\newterm{c4-projective-module}{投射模}{projective module}，若对 $R$-模的任一满同态 $\beta:M\to N$ 与任一态射 $\alpha:P\to N$，存在态射 $\gamma:P\to M$ 使得 $\beta\circ\gamma=\alpha$，即下图交换
\[
\begin{tikzcd}
& P \arrow[d,"\alpha"] \arrow[ld,dotted,"\gamma"'] \\
M \arrow[r,"\beta"'] & N \arrow[r] & 0.
\end{tikzcd}
\]
\end{definition}

下面给出的定理将为我们提供投射模的另一个刻画.

\begin{theorem}\label{c4-thm-2-1}
设 $P$ 是一个 $R$-模，则下面各条等价：
\begin{enumerate}[label=(\arabic*),leftmargin=*]
    \item $P$ 是投射的；
    \item $\operatorname{Hom}(P, \square)$ 是正合的；
    \item 每个 $R$-模的正合列 $0\to M'\to M\to P\to 0$ 都分裂；
    \item $P$ 是某个自由模的直和项，即存在 $R$-模 $M$ 使得 $P\oplus M$ 是自由的.
\end{enumerate}
\end{theorem}

\begin{proof}
证明可参见 \cite[Proposition~3.2, Proposition~3.3 \& Theorem~3.5, Chapter~3]{rot09}.
\end{proof}

我们现在精确定义什么是 $R$-模\footnote{原文即如此；按上下文此处所指应为自由 $R$-模.}.

\begin{definition}\label{c4-def-2-2}
称 $R$-模 $M$ 是\emph{自由的}，若存在集合 $X=\{m_i\mid i\in I\}$（$I$ 可为任意指标集），称为 $M$ 的一组基，使得每个 $m\in M$ 都有唯一表达式
\[
m=\sum_{i\in I}r_im_i
\]
其中 $r_i\in R$，且 $r_i\neq 0$ 的情形只有有限多个.
\end{definition}

\begin{proposition}\label{c4-prop-2-2}
设 $F$ 是一个自由 $R$-模，$M$、$N$ 为任意 $R$-模. 若 $\beta:M\to N$ 是满同态，则对每个 $\alpha:F\to N$，存在 $R$-同态 $\gamma:F\to M$ 使下图交换
\[
\begin{tikzcd}
& F \arrow[d,"\alpha"] \arrow[ld,dotted,"\gamma"'] \\
M \arrow[r,"\beta"'] & N \arrow[r] & 0.
\end{tikzcd}
\]
\end{proposition}

\begin{proof}
该命题的证明可参见 \cite[Theorem~3.1, Chapter~3]{rot09}.
\end{proof}

\begin{corollary}\label{c4-cor-2-3}
每个自由模都是投射模.
\end{corollary}

我们现在定义投射模的对偶概念，即内射模.

\begin{definition}\label{c4-def-2-3}
称 $R$-模 $I$ 是\newterm{c4-injective-module}{内射模}{injective module}，若对任一单同态 $\beta:M\to N$ 与任一 $R$-模的态射 $\alpha:M\to I$，存在态射 $\gamma:N\to I$ 使得 $\gamma\circ\beta=\alpha$，即下图交换
\[
\begin{tikzcd}
& I \\
0 \arrow[r] & M \arrow[u,"\alpha"'] \arrow[r,"\beta"'] & N. \arrow[ul,dotted,"\gamma"']
\end{tikzcd}
\]
\end{definition}

\begin{theorem}\label{c4-thm-2-4}
设 $I$ 是一个 $R$-模，则下面各条等价：
\begin{enumerate}[label=(\arabic*),leftmargin=*]
    \item $I$ 是内射的；
    \item $\operatorname{Hom}_R(\square, I)$ 是正合函子；
    \item 每个 $R$-模的正合列 $0\to I\to M\to M'\to 0$ 都分裂.
\end{enumerate}
\end{theorem}

\begin{proof}
证明可参见 \cite[Proposition~3.25 \& 3.26, \S3.2]{rot09}.
\end{proof}

\section{\texorpdfstring{分解与 $R$-模的上同调}{分解与 R-模的上同调}}

\subsection{分解}

下面定义模的分解，它将提供用来定义上同调的链复形.

\begin{definition}\label{c4-def-3-1}
设 $M$ 是一个 $R$-模. $M$ 在 $R$ 上的一个\newterm{c4-projective-resolution}{投射分解}{projective resolution}（相应地：\emph{自由分解}）是指 $R$-模的一个正合列
\[
\ldots\to P_n\xrightarrow{d_n}P_{n-1}\to\ldots\to P_1\xrightarrow{d_1}P_0\xrightarrow{\varepsilon}M\to 0,
\]
其中每个 $P_n$ 都是投射模（相应地：自由模）.
\end{definition}

\begin{notation}\label{c4-not-3-2}
这样的投射分解有时记作 $P_\bullet\to M$.
\end{notation}

这些分解不仅有用，我们还知道：对任何 $R$-模，这样的分解总可以找到.

\begin{theorem}\label{c4-thm-3-1}
每个 $R$-模都有自由分解.
\end{theorem}

\begin{proof}
参见 \cite[Proposition~6.2, \S6.1]{rot09}.
\end{proof}

\begin{corollary}\label{c4-cor-3-2}
每个 $R$-模都有投射分解.
\end{corollary}

\begin{definition}\label{c4-def-3-3}
设 $M$ 是一个 $R$-模. $M$ 在 $R$ 上的一个\emph{内射分解}是指 $R$-模的一个正合列
\[
0\to M\xrightarrow{\varepsilon}I^0\xrightarrow{\delta^0}I^1\to\ldots\to I^n\xrightarrow{\delta^n}I^{n+1}\to\ldots,
\]
其中每个 $I^n$ 都是内射模.
\end{definition}

\begin{notation}\label{c4-not-3-4}
这样的内射分解有时记作 $M\to I^\bullet$.
\end{notation}

\begin{theorem}\label{c4-thm-3-3}
每个 $R$-模都有内射分解.
\end{theorem}

\begin{proof}
参见 \cite[Proposition~6.4, \S6.1]{rot09}.
\end{proof}

现在定义删失（余）复形，因为它们处理起来更方便.

\begin{definition}\label{c4-def-3-5}
设 $M$ 是一个 $R$-模. 若
\[
X:\ldots\to X_n\xrightarrow{d_n}X_{n-1}\to\ldots\to X_1\to X_0\to M\to 0
\]
是一个链复形，我们定义 $X$ 的\emph{删失复形}为链复形
\[
X_M:\ldots\to X_n\xrightarrow{d_n}X_{n-1}\to\ldots\to X_1\to X_0\to 0.
\]
类似地，若
\[
Y:0\to M\to Y^0\xrightarrow{\delta^0}Y^1\to\ldots\to Y^n\xrightarrow{\delta^n}Y^{n+1}\to\ldots
\]
是一个上链复形，我们定义 $Y$ 的\emph{删失余复形}为上链复形
\[
Y_M:0\to Y^0\xrightarrow{\delta^0}Y^1\to\ldots\to Y^n\xrightarrow{\delta^n}Y^{n+1}\to\ldots
\]
\end{definition}

\subsection{\texorpdfstring{$R$}{R}-模的上同调}

本节我们定义函子 Ext，它是上同调概念的推广\footnote{换言之，$\operatorname{Ext}$ 是函子 $\operatorname{Hom}$ 的\newterm{c4-derived-functor}{导出函子}{derived functor}（原文未使用这一术语）.}.

\begin{definition}\label{c4-def-3-6}
设 $M$、$N$ 是 $R$-模. 令 $P_M$ 为 $M$ 的某个投射分解的删失复形. 构成复形 $\operatorname{Hom}_R(P_M, N)$，即
\[
\operatorname{Hom}_R(P_M, N):0\to\operatorname{Hom}_R(P_0, N)\xrightarrow{d^*_1}\ldots\to\operatorname{Hom}_R(P_n, N)\xrightarrow{d^*_{n+1}}\operatorname{Hom}_R(P_{n+1}, N)\to\ldots
\]
我们把 $M$ 以 $N$ 为系数的第 $n$ 个上同调 $\operatorname{Ext}^n_R(M, N)$ 定义为上述上链复形的第 $n$ 个上同调，即
\[
\operatorname{Ext}^n_R(M, N):=H^n(\operatorname{Hom}(P_M, N)).
\]
\end{definition}

\begin{theorem}\label{c4-thm-3-4}
设 $M$、$N$ 是 $R$-模，$P_M$ 与 $P_M'$ 是 $M$ 的两个投射分解. 设 $\operatorname{Ext}^n_R(M, N)$ 与 $\operatorname{Ext}^{'n}_R(M, N)$ 分别是与 $P_M$、$P_M'$ 相联的上同调模，则
\[
\operatorname{Ext}^n_R(M, N)\cong\operatorname{Ext}^{'n}_R(M, N).
\]
\end{theorem}

\begin{proof}
证明可参见 \cite[Proposition~6.40, \S6.2]{rot09}.
\end{proof}

\begin{remark}\label{c4-rem-3-7}
有些作者利用内射分解来定义 $M$ 以 $N$ 为系数的上同调.\\
令 $I_N$ 为 $N$ 的某个内射分解的删失复形. 构成复形 $\operatorname{Hom}_R(M, I_N)$，即
\[
\operatorname{Hom}_R(M, I_N):0\to\operatorname{Hom}_R(M, I^0)\to\ldots\to\operatorname{Hom}_R(M, I^n)\xrightarrow{\delta^n_*}\operatorname{Hom}_R(M, E^{n+1})\to\ldots
\]
则 $M$ 以 $N$ 为系数的第 $n$ 个上同调为
\[
H^n(M, N)=H^n(\operatorname{Hom}_R(M, I_N)).
\]
这与上面的定义并无歧义，因为 \cite[Theorem~6.67, \S6.2]{rot09} 证明了\footnote{原文上一行写作 $E^{n+1}$，按上下文应为 $I^{n+1}$.}
\[
H^n(\operatorname{Hom}(M, I_N))=H^n(\operatorname{Hom}(P_M, N)).
\]
\end{remark}

\section{群的上同调}

现在我们来精确定义\newterm{c4-group-cohomology}{群的上同调}{group cohomology}. 下文 $R$ 恒为带单位元 $1$ 的交换环.

\begin{definition}\label{c4-def-4-1}
设 $G$ 是一个群，群环 $R[G]$ 是以 $G$ 为基的自由阿贝尔群，即 $G$ 中元素的一切以 $R$ 中元素为系数的形式线性组合构成的群（并赋予由 $G$ 的乘法诱导的乘法）：
\[
R[G]=\Bigl\{\sum_{\substack{g\in G\\ \text{finite}}}r_gg\ \Big|\ r_g\in R\Bigr\}.
\]
整数群\footnote{原文作 integral group，按上下文应为 integral group ring（整数群环）.}是指对整数环 $\mathbb Z$ 取群环所得的 $\mathbb Z[G]$，该环也写作 $\mathbb ZG$.
\end{definition}

\begin{definition}\label{c4-def-4-2}
设 $(G,\cdot)$ 是一个群. 一个 $G$-\emph{模}是指配有映射 $G\times M\to M:(g,m)\mapsto g\cdot m$ 的阿贝尔群 $(M,+)$，且对 $m,n\in M$ 与 $g_1,g_2,g\in G$ 满足：
\begin{enumerate}[label=(\arabic*),leftmargin=*]
    \item $(g_1g_2)\cdot(m)=g_1\cdot(g_2\cdot m)$；
    \item $g\cdot(m+n)=g\cdot m+g\cdot n$；
    \item $1_G\cdot m=m$.
\end{enumerate}
\end{definition}

\begin{remark}\label{c4-rem-4-3}
若 $M$ 是一个 $G$-模，则按
\[
\Bigl(\sum_{\substack{g\in G\\ \text{finite}}}n_gg\Bigr)\cdot m=\sum_{\substack{g\in G\\ \text{finite}}}n_g(g\cdot m)
\]
定义作用后，$M$ 成为一个 $\mathbb ZG$-模. 反之，若 $M$ 是一个 $\mathbb ZG$-模，则按 $g\cdot m=(1g)\cdot m$ 定义作用后，$M$ 成为一个 $G$-模.
\end{remark}

\begin{definition}\label{c4-def-4-4}
设 $G$ 是一个群. 称 $G$-模 $M$ 是\emph{平凡的}，若 $G$ 的每个元素在 $M$ 上的作用都是恒等映射，即 $g\cdot m=m$.
\end{definition}

\begin{definition}\label{c4-def-4-5}
设 $G$ 是一个群，$M$ 是一个 $G$-模. 把整数 $\mathbb Z$ 视为平凡 $G$-模，并定义 $G$ 以 $M$ 为系数的第 $n$ 个上同调群为\footnote{原文公式中写作 $H^n(G, N):=\operatorname{Ext}^n_{\mathbb ZG}(\mathbb Z, N)$，按上下文 $N$ 应为 $M$.}
\[
H^n(G, N):=\operatorname{Ext}^n_{\mathbb ZG}(\mathbb Z, N).
\]
\end{definition}

\section{Shapiro 引理与诱导模}

本节我们介绍 Shapiro 引理与\newterm{c4-induced-module}{诱导模}{induced module}. 它们在后面研究 Galois 群的上同调时将派上用场.\\
设 $G$ 是一个群.

\begin{definition}\label{c4-def-5-1}
称 $\mathbb ZG$-模 $A$ 是 $G$-\emph{无循环的}（$G$-acyclic），若对一切 $n\geq 1$ 有 $H^n(G, A)=\{0\}$.
\end{definition}

\begin{definition}\label{c4-def-5-2}
群同态 $f:G'\to G$ 诱导出环同态 $\mathbb ZG'\to\mathbb ZG$，仍记作 $f$，确切地说，$f:\sum n_{g'}g'\mapsto\sum n_{g'}f(g')$. 设 $A$ 是一个 $\mathbb ZG$-模，它可以被视为 $\mathbb ZG'$-模：若 $g'\in G'$ 且 $a\in A$，则 $g'a=f(g')a$. 把视为 $G'$-模的 $G$-模 $A$ 记作 $UA$，并称
\[
U_f:\mathscr{M}_{\mathbb ZG}\to\mathscr{M}_{\mathbb ZG'}
\]
为一个\emph{换群函子}.
\end{definition}

\begin{lemma}\label{c4-lem-5-1}
设 $f:G'\to G$ 是群同态，$U_f:\mathscr{M}_{\mathbb ZG}\to\mathscr{M}_{\mathbb ZG'}$ 是相应的换群函子.
\begin{enumerate}[label=(\arabic*),leftmargin=*]
    \item 若 $\mathbf P$ 是一个 $G$-循环复形，则 $\mathbf{UP}$ 是一个 $G'$-循环复形.
    \item 设 $H\subseteq G$，$f:H\to G$ 为包含同态. 若 $\mathbf P$ 是投射 $G$-模，则 $UP$ 是投射 $H$-模. 此外，若 $\mathbf P$ 是 $\mathbb ZG$-模 $A$ 在 $\mathbb ZG$ 上的一个投射分解，则 $U\mathbf P$ 是 $UA$ 在 $\mathbb ZS$-上的一个投射分解\footnote{原文此处即止于 ``$\mathbb ZS$-''，按上下文应为 $\mathbb ZH$.}.
\end{enumerate}
\end{lemma}

\begin{proof}
\begin{enumerate}[label=(\arabic*),leftmargin=*]
    \item 由 \cite[Proposition~8.3, \S8.2]{rot09}，我们知道 $U:\mathscr{M}_{\mathbb ZG}\to\mathscr{M}_{\mathbb ZG'}$ 是一个正合可加函子. 于是 $\mathbf{UP}$ 是一个复形，此外 $H^n(G', A)=\{0\}$.
    \item 取 $R$ 为这样一个集合：它含有 $G/H$ 的每个陪集中的（且仅一个）元素. 于是 $G$ 是不交并 $\bigcup_{r\in R}rH$. 这意味着对每个 $g\in G$，存在唯一的 $r\in R$ 与 $h\in H$ 使得 $g=rh$. 从而我们可以写成 $\mathbb ZG=\bigoplus_{r\in R}r(\mathbb ZH)$（作为 $\mathbb ZS$-模的直和），故 $\mathbb ZG$ 是一个自由 $\mathbb ZH$-模.

设 $P$ 是一个投射 $\mathbb ZG$-模. 这等价于存在 $\mathbb ZG$-模 $M$ 使得 $P\oplus M$ 是自由的. 于是 $UP$ 是一个投射 $\mathbb ZH$-模.

第二条断言由 (1) 推出.
\end{enumerate}
\end{proof}

\begin{remark}\label{c4-rem-5-3}
设 $H$ 是 $G$ 的子群，$A$ 是一个 $\mathbb ZH$-模. 则 $\operatorname{Hom}_{\mathbb ZH}(\mathbb ZG, A)$ 赋予如下作用后是一个 $\mathbb ZG$-模：
\[
gf:\mathbb ZG\to A:x\mapsto gf(g^{-1}x),
\]
对一切 $g\in G$ 与一切 $f:\mathbb ZG\to A$ 成立. 我们特别考虑 $H=\{1\}$ 的情形（此时 $\mathbb ZH=\mathbb Z$）.
\end{remark}

其次，我们介绍 Eckman-Shapiro 引理（即通常所称的 Shapiro 引理）.

\begin{theorem}\label{c4-thm-5-2}
设 $G$ 是一个群，$H$ 是子群，$A$ 是一个 $\mathbb ZH$-模. 则
\[
H^n(H, A)\cong H^n\bigl(G, \operatorname{Hom}_{\mathbb ZH}(\mathbb ZG, A)\bigr)
\]
对一切 $n\geq 0$ 成立.
\end{theorem}

在本证明中，我们要用到张量积的概念，记作 $\otimes$. 关于它的更多信息，参见 \cite[\S2.2]{rot09}.

\begin{proof}
记 $U:\mathscr{M}_{\mathbb ZG}\to\mathscr{M}_{\mathbb ZH}$ 为换群函子. 设 $\mathbf P=\ldots\to P_1\to P_0\to\mathbb Z\to 0$ 是 $\mathbb Z$ 在 $\mathbb ZG$ 上的一个自由分解，则
\[
H^n\bigl(G, \operatorname{Hom}_{\mathbb ZH}(\mathbb ZG, A)\bigr)=H^n\bigl(\operatorname{Hom}_{\mathbb ZG}(\mathbf P, \operatorname{Hom}_{\mathbb ZG}(\mathbb ZG, A))\bigr),
\]
因为这正是上同调群的定义. 利用张量积与函子 Hom 之间的伴随关系（确切地说，参见 \cite[Theorem~2.75, \S2.2]{rot09}），我们有
\[
\operatorname{Hom}_{\mathbb ZG}\bigl(P_i, \operatorname{Hom}_{\mathbb ZH}(\mathbb ZG, A)\bigr)\cong\operatorname{Hom}_{\mathbb ZH}\bigl(P_i\bigotimes_{\mathbb ZG}\mathbb ZG, A\bigr).
\]
其次，利用张量积的一条性质（见 \cite[Proposition~2.58, \S2.2]{rot09}），有 $P_i\otimes_{\mathbb ZG}\cong P_i$\footnote{原文即如此；按上下文应为 $P_i\otimes_{\mathbb ZG}\mathbb ZG\cong P_i$.}，并用事实 $UP_i\cong P_i$. 于是对一切 $i$，我们得到
\[
\operatorname{Hom}_{\mathbb ZG}\bigl(P_i, \operatorname{Hom}_{\mathbb ZH}(\mathbb ZG, A)\bigr)\cong\operatorname{Hom}_{\mathbb ZH}(UP_i, A).
\]
由 \autoref{c4-lem-5-1}(2)，$U\mathbf P$ 是 $\mathbb Z$ 在 $\mathbb ZH$ 上的一个投射分解，从而
\[
H^n(H, A)\cong H^n(\operatorname{Hom}_{\mathbb ZH}(U\mathbf P, A),
\]
参见 \autoref{c4-thm-3-4}. 此外，存在复形的同构
\begin{equation}\label{c4-eq-1}
\operatorname{Hom}_{\mathbb ZH}(U\mathbf P, A)\cong\operatorname{Hom}_{\mathbb ZG}\bigl(\mathbf P, \operatorname{Hom}_{\mathbb ZH}(\mathbb ZG, A)\bigr).
\end{equation}
总之，由 \eqref{c4-eq-1} 我们得到同构
\[
H^n(H, A)\cong H^n\bigl(G, \operatorname{Hom}_{\mathbb ZH}(\mathbb ZG, A)\bigr).
\]
\end{proof}

\begin{definition}\label{c4-def-5-4}
设 $G$ 是一个群，$A$ 是一个阿贝尔群. 我们定义\emph{余诱导模}为 $\mathbb ZG$-模 $\operatorname{Hom}_{\mathbb Z}(\mathbb ZG, A)$.
\end{definition}

\begin{remark}\label{c4-rem-5-5}
上述构造与 \autoref{c4-rem-5-3} 中取 $H=\{1\}$ 的构造相同，这里注意到每个阿贝尔群都可以视为带平凡作用的 $\mathbb Z$-模.
\end{remark}

\begin{proposition}\label{c4-prop-5-3}
每个余诱导 $\mathbb ZG$-模 $A$ 都是 $G$-循环的：$H^n(G, A)=\{0\}$ 对一切 $n\geq 1$ 成立.
\end{proposition}

\begin{proof}
考虑子群 $H=\{1\}$ 并应用 \autoref{c4-thm-5-2}，则
\[
H^n\bigl(G, \operatorname{Hom}_{\mathbb Z}(\mathbb ZG, A)\bigr)\cong H^n(H, A).
\]
但该上同调群是平凡的（见 \cite[Corollary~9.28, \S9.2]{rot09}）.
\end{proof}

% !TeX root = ../main.tex
% 第 5 章 Galois 上同调 （原书 CHAPTER 5, Galois Cohomology, 印刷页 45-56）

\chapter{Galois 上同调}
\label{ch:5}

本章为投射有限群定义上同调群，注意我们将采用与前一章略有不同的定义. 最后，我们陈述并证明 Hilbert 第 $90$ 定理的各种形式.

\section{上同调群}

处理投射有限群意味着我们必须考虑拓扑. 因此我们定义离散模与特定的上同调群.

\subsection{离散 \texorpdfstring{$G$}{G}-模}
\label{c5-sub-1-1}

\begin{definition}\label{c5-def-1-1}
设 $(G,\cdot)$ 是一个投射有限群. 一个\newterm{c5-discrete-gmodule}{离散 $G$-模}{discrete $G$-module}（或简称 $G$-模，若不会引起混淆）是指赋予离散拓扑且 $G$ 在其上连续作用的阿贝尔群 $(A,+)$. 这就是说，$G$-模是一个阿贝尔群 $A$，配有连续映射 $G\times A\to A:(g,a)\mapsto g\cdot a$（其中 $G\times A$ 赋予乘积拓扑），满足：
\begin{enumerate}[label=(\arabic*),leftmargin=*]
    \item $(g_1g_2)\cdot(a)=g_1\cdot(g_2\cdot a)$；
    \item $g\cdot(a+b)=g\cdot a+g\cdot b$；
    \item $1_G\cdot a=a$；
\end{enumerate}
对一切 $a,b\in A$ 与 $g_1,g_2,g\in G$ 成立.
\end{definition}

\begin{remark}\label{c5-rem-1-2}
若 $G$ 是有限群，则一切 $G$-模都是离散 $G$-模.
\end{remark}

\begin{proposition}\label{c5-prop-1-1}
设 $G$ 是投射有限群，$A$ 是阿贝尔群. 设 $G\times A\to A$ 是 $G$ 在 $A$ 上的一个作用，满足\autoref{c5-def-1-1} 的条件 (1)、(2) 与 (3). 则下列各条等价：
\begin{enumerate}[label=(\arabic*),leftmargin=*]
    \item $G\times A\to A$ 连续；
    \item 对每个 $a\in A$，$a$ 的稳定子 $U_a=\{g\in G\mid g\cdot a=a\}$ 是 $G$ 的开子群；
    \item $A=\bigcup_U A^U$，其中 $U$ 取遍 $G$ 的全体开子群，且
    \[
    A^U=\{a\in A\mid g\cdot a=a\ \forall\, g\in U\}.
    \]
\end{enumerate}
\end{proposition}

\begin{proof}
记 $f:G\times A\to A$ 为 $G$ 在 $A$ 上的作用.

(1) $\Rightarrow$ (2)：考虑限制 $f_a:G\times\{a\}\to A$，因 $f$ 连续，它是连续映射. 于是 $U_a=f_a^{-1}(\{a\})$ 作为开集的原像是开集（注意 $A$ 赋予离散拓扑）.

(2) $\Rightarrow$ (3)：显然 $\bigcup_U A^U\subseteq A$. 设 $a\in A$，则 $a$ 属于 $A^{U_a}$. 而 $U_a$ 是 $G$ 的开子群，故得结论.

(3) $\Rightarrow$ (1)：我们证明对每个 $a\in A$，$f^{-1}(\{a\})$ 是 $G$ 的开子群. 于是若 $a\in A$，则 $a$ 属于某个 $A^U$，其中 $U$ 是 $G$ 的开子群. 这意味着对一切 $g\in U$ 有 $g\cdot a=a$. 从而
\[
1_G\times\{a\}\in U\times\{a\}\subseteq f^{-1}(\{a\}).
\]
因此 $f^{-1}(\{a\})$ 是开集，因为 $G\times A$ 是拓扑群且 $U\times\{a\}$ 是开集（见\autoref{c1-prop-4-4}）.
\end{proof}

下面是离散 $G$-模的一些例子.

\begin{examples}\label{c5-ex-1-3}
\
\begin{itemize}
    \item 对任意投射有限群 $G$ 与任意阿贝尔群 $A$，我们可定义作用
    \[
    G\times A\to A:(g,a)\mapsto a.
    \]
    此作用称为 $A$ 上的\emph{平凡作用}. 而 $A$ 是一个 $G$-模，称为\emph{平凡 $G$-模}.
    \item 设 $\mathbb{L}$ 与 $\mathbb{K}$ 是域，且 $\mathbb{L}$ 是 $\mathbb{K}$ 的 Galois 扩张. 回忆 $\mathrm{Gal}(\mathbb{L},\mathbb{K})$ 是投射有限群. 考虑阿贝尔群 $(\mathbb{L},+)$ 与作用
    \[
    G\times\mathbb{L}\to\mathbb{L}:(\sigma,\alpha)\mapsto\sigma(\alpha).
    \]
    则 $(\mathbb{L},+)$（有时记作 $\mathbb{L}^+$）是离散 $G$-模. 赋予同样的作用，$(\mathbb{L}^*,\cdot)$ 以及 $\mathbb{L}$ 中全体单位根构成的群（对乘法）也是离散 $G$-模.
\end{itemize}
\end{examples}

\begin{definition}\label{c5-def-1-4}
设 $G$ 是投射有限群，$A$、$B$ 是两个离散 $G$-模. 一个 $G$-同态或 $G$-映射 $\varphi:A\to B$ 是指满足下式的群同态
\[
\varphi(g\cdot a)=g\cdot\varphi(a),
\]
对一切 $g\in G$ 与一切 $a\in A$ 成立.
\end{definition}

\begin{remark}\label{c5-rem-1-5}
离散 $G$-模与 $G$-映射构成的类是一个阿贝尔范畴，记作 $\mathscr{M}_G$.
\end{remark}

\subsection{上同调群的定义}
\label{c5-sub-1-2}

设 $G$ 是投射有限群. 对每个 $n\in\mathbb{N}$，以 $G^n$ 记 $G$ 的 $n$ 个副本的笛卡尔积.

考虑一个离散 $G$-模 $A$. 我们令
\[
C^n_c(G,A):=\{f:G^n\to A\mid f\ \text{连续}\}.
\]
注意 $C^n_c(G,A)$ 对加法构成阿贝尔群（因为 $A$ 是阿贝尔群）. 它称为 $n$-\emph{上链}群. 此外 $C^0_c(G,A)$ 同构于 $A$（作为从 $0$ 到 $A$ 的映射群）.

对每个 $n\geq 1$ 定义群同态
\[
d_{n+1}:C^n_c(G,A)\to C^{n+1}_c(G,A),
\]
使得
\[
\begin{aligned}
(d_{n+1}f)(g_1,\ldots,g_{n+1})=\ &g_1\cdot f(g_2,\ldots,g_{n+1})\\
+\ &\sum_{i=1}^{n}(-1)^i f(g_1,\ldots,g_ig_{i+1},\ldots,g_{n+1})\\
+\ &(-1)^{n+1}f(g_1,\ldots,g_n),
\end{aligned}
\]
对一切 $f\in C^n_c(G,A)$ 与一切 $(g_1,\ldots,g_{n+1})\in G^{n+1}$ 成立. 对 $n=0$，我们定义同态
\[
d_0:C^0_c(G,A)\to C^1_c(G,A),
\]
使得
\[
(d_0a)(g)=g\cdot a-a,
\]
对一切 $a\in A$ 与一切 $g\in G$ 成立. 映射 $d_n$ 称为\emph{上边缘算子}.

\begin{lemma}\label{c5-lem-1-2}
对每个 $n\geq 0$，$d_n$ 是群同态，并且还有 $d_{n+1}\circ d_n=0$.
\end{lemma}

\begin{proof}
证明可参见 \cite[Chapter~2, \S2.2]{bur04}.
\end{proof}

\begin{remark}\label{c5-rem-1-6}
此引理蕴含映射 $d_n$ 是某个上链复形的微分
\[
\overline{C}_c(G,A):=(C^n_c(G,A),d_n).
\]
\end{remark}

\begin{definition}\label{c5-def-1-7}
设 $G$ 是投射有限群，$A$ 是离散 $G$-模.
\begin{itemize}
    \item 称 $G$ 的 $n$-\newterm{c5-cocycle}{上闭链}{cocycle}群为 $Z^n_c:=\ker(d_n)$，称 $G$ 的 $n$-上边缘群为 $B^n_c(G,A):=\operatorname{im}d_{n-1}$. \autoref{c5-lem-1-2} 蕴含 $B^n_c(G,A)$ 是 $Z^n_c(G,A)$ 的子群.
    \item 我们定义 $G$ 关于系数 $A$ 的第 $n$ 个\newterm{c5-continuous-cohomology}{（连续）上同调群}{(continuous) cohomology group}为群
    \[
    H^n_c(G,A):=Z^n_c(G,A)/B^n_c(G,A).
    \]
\end{itemize}
\end{definition}

\begin{remark}\label{c5-rem-1-8}
当 $A$ 仅仅是一个 $G$-模而不要求上链连续时，我们也可构造上同调. 构造方式相同. 此时 $G$ 关于系数 $A$ 的上同调群记作 $H^n_c(G,A)$\footnote{原文即如此；此处应为 $H^n(G,A)$——$H^n_c$ 已被用于连续上同调，否则下文所谓记号上不存在歧义即不成立.}. 可以证明这一上同调群的定义与第 $4$ 章的一致，即记号上不存在歧义. 当 $G$ 是离散群时，连续性假设自动满足. 这意味着 $H^n_c(G,A)=H^n(G,A)$. 于是我们也可使用第 $4$ 章中给出的性质.
\end{remark}

\section{低维上同调群的解释}

设 $G$ 是投射有限群，$A$ 是离散 $G$-模. 我们研究小 $n\in\mathbb{N}$ 时的上同调群 $H^n_c(G,A)$.

\subsection{平凡情形}
\label{c5-sub-2-1}

对 $n=0$，我们有
\[
H^0_c(G,A)\cong Z^0_c(G,A)=\{a\in A\mid d_0(a)=0\}=\{a\in A\mid g\cdot a-a=0\ \forall\, g\in G\}.
\]
于是 $H^0_c(G,A)$ 同构于 $A^G$，即 $A$ 中在 $G$ 作用下保持不变的元素之集.

\subsection{第一上同调群}
\label{c5-sub-2-2}

现在我们考虑 $n=1$. 如前所定义 $H^1_c(G,A)=Z^1_c(G,A)/B^1_c(G,A)$. $1$-上闭链群可明确描述为
\[
\begin{aligned}
Z^1_c(G,A)=\ &\{f:G\to A\ \text{连续}\mid d_1(f)=0\}\\
=\ &\{f:G\to A\ \text{连续}\mid f(g_1g_2)=g_1\cdot f(g_2)+f(g_1),\ \forall\, g_1,g_2\in G\}.
\end{aligned}
\]
而 $1$-上边缘群为
\[
\begin{aligned}
B^1_c(G,A)=\ &\{f:G\to A\ \text{连续}\mid \exists\, a\in A\ \text{使得}\ d_0(a)=f\}\\
=\ &\{f:G\to A\ \text{连续}\mid \exists\, a\in A\ \text{使得}\ f(g)=g\cdot a-a,\ \forall\, g\in G\}.
\end{aligned}
\]
$Z^1_c(G,A)$ 与 $B^1_c(G,A)$ 中的元素分别称为连续\newterm{c5-crossed-homomorphism}{交叉同态}{crossed homomorphism}与\newterm{c5-principal-crossed}{主交叉同态}{principal crossed homomorphism}.

\begin{example}\label{c5-ex-2-1}
若 $G$ 在 $A$ 上平凡作用，则 $H^1_c(G,A)$ 是从 $G$ 到 $A$ 的全体连续群同态构成的群.
\end{example}

\subsection{第二上同调群}
\label{c5-sub-2-3}

对 $n=2$，$H^2_c(G,A)=Z^2_c(G,A)/B^2_c(G,A)$，其中
\[
\begin{aligned}
Z^2_c(G,A)=\ &\{f:G\times G\to A\ \text{连续}\mid d_2(f)=0\}\\
=\ &\{f:G\times G\to A\ \text{连续}\mid g_1f(g_2,g_3)+f(g_1,g_2g_3)=f(g_1g_2,g_3)+f(g_1,g_2),\\
&\hspace{4em}\forall\, g_1,g_2,g_3\in G\},
\end{aligned}
\]
且
\[
\begin{aligned}
B^2_c(G,A)=\ &\{f:G\times G\to A\ \text{连续}\mid \exists\,\varphi\in C^1_c(G,A)\ \text{使得}\ f=d_1(\varphi)\}\\
=\ &\{f:G\times G\to A\ \text{连续}\mid \exists\,\varphi\in C^1_c(G,A)\ \text{使得}\\
&\hspace{2em}f(g_1,g_2)=g_1\varphi(g_2)-\varphi(g_1g_2)+\varphi(g_1),\ \forall\, g_1,g_2\in G\}.
\end{aligned}
\]
$Z^2_c(G,A)$ 中的元素称为连续因子集.

\section{上同调群的函子性质}

\subsection{相容映射}
\label{c5-sub-3-1}

\subsubsection{定义}

\begin{definition}\label{c5-def-3-1}
设 $G$ 与 $G'$ 是两个投射有限群. 设 $A\in\mathscr{M}_G$ 且 $A'\in\mathscr{M}_{G'}$. 考虑投射有限群的连续同态 $\varphi:G\to G'$ 与群同态 $\psi:A'\to A$. 若
\[
\psi(\varphi(g)\cdot a')=g\cdot f(a'),\footnote{原文即如此；此处 $f$ 应为 $\psi$.}
\]
对一切 $g\in G$ 与一切 $a'\in A'$ 成立，则称 $\psi$ 与 $\varphi$ 是\emph{相容映射}.

这等价于要求：当 $A'$ 视为带有作用 $g\cdot a'=\varphi(g)\cdot a'$（对每个 $a'\in A'$ 与 $g\in G$）的 $G$-模时，$\psi$ 是一个 $G$-映射.
\end{definition}

\begin{example}\label{c5-ex-3-2}
设 $\mathbb{L}$ 与 $\mathbb{E}$ 是 $\mathbb{K}$ 的 Galois 扩张，且 $\mathbb{K}\subseteq\mathbb{E}\subseteq\mathbb{L}$. 我们考虑限制
\[
\pi:\mathrm{Gal}(\mathbb{L},\mathbb{K})\to\mathrm{Gal}(\mathbb{E},\mathbb{K}),
\]
与单射
\[
i:\mathbb{E}^*\hookrightarrow\mathbb{L}^*.
\]
则这两个映射相容.
\end{example}

\begin{lemma}\label{c5-lem-3-1}
设 $\varphi$ 与 $\psi$ 是如\autoref{c5-def-3-1} 中的相容映射. 它们在 $n$-上链群（对 $n\geq 0$）上诱导出同态
\[
(\varphi,\psi):C^n_c(G',A')\to C^n_c(G,A)
\]
定义为
\[
[(\varphi,\psi)(f)](g_1,\ldots,g_n)=\psi[f(\varphi(g_1),\ldots,\varphi(g_n))].
\]
特别地，$(\varphi,\\psi)$ 是上链复形之间的映射，即下图对 $n\geq 0$ 交换：
\[
\begin{tikzcd}
\cdots \arrow[r, "d_n"] & C^n_c(G',A') \arrow[r, "d_{n+1}"] \arrow[d, "{(\varphi,\psi)}"'] & C^{n+1}_c(G',A') \arrow[d, "{(\varphi,\psi)}"'] \arrow[r] & \cdots \\
\cdots \arrow[r, "d_n"] & C^n_c(G,A) \arrow[r, "d_{n+1}"] & C^{n+1}_c(G,A) \arrow[r] & \cdots
\end{tikzcd}
\]
\end{lemma}

\begin{proof}
通过直接计算即可轻易验证此引理.
\end{proof}

\begin{corollary}\label{c5-cor-3-2}
若 $\varphi$、$\psi$ 是如\autoref{c5-def-3-1} 中的相容映射，则 $(\varphi,\psi)$ 在上同调群上诱导出同态，记作
\[
(\varphi,\psi):H^n_c(G',A')\to H^n_c(G,A).
\]
\end{corollary}

\begin{proof}
利用上述图的交换性.
\end{proof}

\subsubsection{相容映射的函子性质}

设 $G_i$（$i=1,2,3$）是投射有限群，$A_i\in\mathscr{M}_{G_i}$. 考虑连续群同态
\[
G_1\xrightarrow{\ \varphi_1\ }G_2\xrightarrow{\ \varphi_2\ }G_3
\]
与群同态
\[
A_1\xleftarrow{\ \psi_1\ }A_2\xleftarrow{\ \psi_2\ }A_3.
\]
其中 $\varphi_1$、$\psi_1$ 是相容映射，$\varphi_2$、$\psi_2$ 亦然. 它们诱导出两个群同态
\[
\begin{aligned}
(\varphi_1,\psi_1):H^n_c(G_2,A_2)\ &\to\ H^n_c(G_1,A_1)\\
(\varphi_2,\psi_2):H^n_c(G_3,A_3)\ &\to\ H^n_c(G_2,A_2).
\end{aligned}
\]
于是 $\varphi_1\circ\varphi_2$\footnote{原文即如此；按各映射的定义域，$\varphi_1\circ\varphi_2$ 与 $\psi_2\circ\psi_1$ 无定义，应为 $\varphi_2\circ\varphi_1$ 与 $\psi_1\circ\psi_2$. 下式右端原文写作 $(\varphi_2\circ\varphi_1,\psi_1\circ\psi_2)$，无误.} 与 $\psi_2\circ\psi_1$ 相容，且
\[
(\varphi_1,\psi_1)\circ(\varphi_2,\psi_2)=(\varphi_2\circ\varphi_1,\psi_1\circ\psi_2).
\]
此外，若 $\varphi$ 是 $G$ 的恒等映射，$\psi$ 是 $A$ 的恒等映射，则 $(\varphi,\psi)$ 是 $H^n_c(G,A)$ 的恒等映射.

于是，对每个投射有限群 $G$ 与每个 $n\geq 0$，$H^n_c(G,-)$ 是从 $\mathscr{M}_G$ 到 $\mathscr{A}$ 的协变函子.

\subsection{正向极限}
\label{c5-sub-3-2}

我们定义正向极限，以便考虑 $G$-模的正向极限. 这将有助于描述上同调群的极限.

\begin{definition}\label{c5-def-3-3}
设 $I$ 是有向集，$\mathscr{C}$ 是一个范畴. 由 $I$ 指标化的 $\mathscr{C}$ 中对象的一个\emph{正向系} $(X_i,f_{ij})_I$ 由一族对象 $\{X_i\mid i\in I\}$ 与一族 $\mathscr{C}$ 中的态射 $\{f_{ij}:X_i\to X_j\mid i,j\in I,\,i\leq j\}$ 构成，且满足
\begin{enumerate}[label=(\arabic*),leftmargin=*]
    \item 对一切 $i\in I$，$f_{ii}$ 是 $X_i$ 上的恒等映射；
    \item 对一切 $i\leq j\leq k\in I$，$f_{jk}\circ f_{ij}=f_{ik}$.
\end{enumerate}
\end{definition}

\begin{definition}\label{c5-def-3-4}
若 $Y$ 是 $\mathscr{C}$ 的一个对象，我们称 $\mathscr{C}$ 的一族态射 $\{\psi_i:X_i\to Y\}_{i\in I}$ 是\emph{相容的}，若对每个 $i\leq j\in I$ 有 $\psi_j\circ f_{ij}=\psi_i$，即下图交换：
\[
\begin{tikzcd}
& Y & \\
X_i \arrow[ur,"\psi_i"'] \arrow[rr,"f_{ij}"'] & & X_j. \arrow[ul,"\psi_j"']
\end{tikzcd}
\]
\end{definition}

\begin{definition}\label{c5-def-3-5}
$(X,\varphi_i)$ 称为 $\mathscr{C}$ 中对象的一个正向系 $(X_i,f_{ij})_I$ 的\emph{余极限}，若 $X$ 配有相容族 $\{\varphi_i:X_i\to X\}_{i\in I}$，且满足下列泛性质：对每个对象 $Y$ 与每个相容族 $\{\psi_i:X_i\to Y\}_{i\in I}$，存在唯一的 $\mathscr{C}$ 中态射 $\theta:Y\to X$\footnote{原文即如此；按图示与余极限的泛性质，此处应为 $\theta:X\to Y$.}，使得下图对一切 $i\leq j\in I$ 交换：
\[
\begin{tikzcd}
& Y & \\
X_i \arrow[ur,"\psi_i"'] \arrow[dr,"\varphi_i"'] \arrow[rr,"f_{ij}"'] & & X_j \arrow[ul,"\psi_j"'] \arrow[dl,"\varphi_j"'] \\
& X \arrow[uu,dashed,"{\exists!\,\theta}"'] &
\end{tikzcd}
\]
\end{definition}

\begin{proposition}\label{c5-prop-3-3}
正向极限在同构意义下唯一.
\end{proposition}

\begin{proof}
与\autoref{c3-prop-1-1} 的证明类似.
\end{proof}

\begin{notation}\label{c5-not-3-6}
由于不再有歧义，我们以 $\varinjlim_I X_i$ 记 $(X_i,f_{ij})_I$ 的正向极限.
\end{notation}

我们在第 $3$ 章中已经看到，当处于拓扑空间范畴时，极限（或逆极限）有一个特殊的刻画. 现在我们给出群范畴（及群同态）中正向极限的这样一个刻画.

\begin{proposition}\label{c5-prop-3-4}
设 $(X_i,f_{ij})_I$ 是一个群的正向系. 考虑集合 $Z=\coprod_I X_i$（$X_i$ 的不交并）并定义等价关系 $\sim$，使得 $(x_i,i)\sim(x_j,j)$ 当且仅当 $i,j\in I$ 且 $x_i\in X_i$、$x_j\in X_j$，并且存在 $k\in I$ 满足 $i,j\leq k$ 与 $f_{ik}(x_i)=f_{jk}(x_j)$. 则 $\varinjlim_I X_i\cong Z/\!\sim$.
\end{proposition}

\begin{proof}
容易看出 $\sim$ 是等价关系. 我们定义态射
\[
\varphi_i:X_i\to Z/\!\sim,
\]
它由包含映射 $\iota:X_i\to Z:x_i\mapsto(x_i,i)$ 诱导. 按构造，$\{\varphi_i:X_i\to Z/\!\sim\}$ 是相容系.

考虑另一个群 $Y$ 与相容族 $\{\psi_i:X_i\to Y\}_{i\in I}$. 我们定义
\[
\theta:Z\to Y:(x_i,i)\mapsto\psi_i(x_i).
\]
设 $(x_i,i),(x_j,j)\in Z$ 且 $(x_i,i)\sim(x_j,j)$. 我们要证明 $\theta((x_i,i))=\theta((x_j,j))$. 但 $(x_i,i)\sim(x_j,j)$ 蕴含存在 $k\in I$ 满足 $i,j\leq k$ 与 $f_{ik}(x_i)=f_{jk}(x_j)$. 于是
\[
\begin{aligned}
\theta((x_i,i))&=\psi_i(x_i)\\
&=\psi_k(f_{ik}(x_i))\\
&=\psi_k(f_{jk}(x_j))\\
&=\psi_j(x_j)=\theta((x_j,j)).
\end{aligned}
\]
因此可定义 $\widetilde{\theta}:Z/\!\sim\ \to Y$. 此外显然有 $\widetilde{\theta}\circ\varphi_i=\psi_i$ 对每个 $i\in I$ 成立，且 $\widetilde{\theta}$ 唯一. 于是 $Z/\!\sim$ 是 $(X_i,f_{ij})_I$ 的正向极限.
\end{proof}

\subsection{极限的上同调群的性质}
\label{c5-sub-3-3}

现在我们考虑一个有向指标集 $(I,\leq)$ 与两个系：投射有限群的逆系 $(G_i,\varphi_{ij})_I$ 与阿贝尔群的正向系 $(A_i,\lambda_{ij})$（定义见\autoref{c5-sub-3-2}），使得对每个 $i\in I$ 有 $A_i\in\mathscr{M}_{G_i}$，且对 $i\leq j$，$\varphi_{ij}$ 与 $\lambda_{ij}$ 是相容映射. 这诱导出一个正向系 $(H^n_c(G_i,A_i),(\varphi_{ij},\lambda_{ij}))_I$.

我们记 $G:=\varprojlim_I G_i$ 与 $A:=\varinjlim_I G_i$\footnote{原文即如此；此处应为 $A:=\varinjlim_I A_i$.}，并有典范态射
\[
\begin{aligned}
\pi_i:G\ &\to\ G_i\\
\lambda_i:A_i\ &\to\ A.
\end{aligned}
\]

\begin{lemma}\label{c5-lem-3-5}
$A$ 可通过下列作用视为 $G$-模：对 $a\in A$ 与 $g\in G$，取 $i\in I$ 与 $a_i\in A_i$ 使得 $\lambda_i(a_i)=a$. 然后定义
\[
g\cdot a=\lambda_i((\pi_i(g))\cdot a_i).
\]
此作用定义良好且连续.
\end{lemma}

\begin{proof}
我们先证明此作用与 $i\in I$ 的选取无关. 回忆
\[
A\cong\coprod_I A_i/\!\sim,
\]
其中 $\alpha_i\sim\alpha_j\Leftrightarrow$ $\alpha_i\in A_i$、$\alpha_j\in A_j$ 且存在 $k\in I$ 满足 $i,j\leq k$ 与 $\lambda_{ik}(\alpha_i)=\lambda_{jk}(\alpha_j)$. 因此设 $i,j\in I$，$a_i,a_j$ 分别属于 $A_i,A_j$，使得 $\lambda_j(a_j)=a$ 且 $\lambda_i(a_i)=a$. 因 $I$ 是有向集，存在 $k\in I$ 满足 $i,j\leq k$. 于是，利用逆极限的性质以及 $\pi_{ik}$ 与 $\lambda_{ik}$ 的相容性，
\[
\begin{aligned}
\lambda_{ik}[(\pi_i(g))\cdot a_i]&=\lambda_{ik}[(\pi_{ik}(\pi_k(g)))\cdot a_i]\\
&=\pi_k(g)\cdot\lambda_{ik}(a_i).
\end{aligned}
\]
同理 $\lambda_{jk}[(\pi_j(g))\cdot a_j]=\pi_k(g)\cdot\lambda_{jk}(a_j)$. 而由正向极限的定义，
\[
\lambda_k(\lambda_{ik}(a_i))=\lambda_i(a_i)=a=\lambda_j(a_j)=\lambda_k(\lambda_{jk}(a_j)).
\]
因 $\lambda_k$ 单射，$\lambda_{ik}(a_i)=\lambda_{jk}(a_j)$. 这意味着
\[
\lambda_{ik}[(\pi_i(g))\cdot a_i]=\lambda_{jk}[(\pi_j(g))\cdot a_j].
\]
从而
\[
\lambda_i[(\pi_i(g))\cdot a_i]=\lambda_j[(\pi_j(g))\cdot a_j].
\]
现在容易看出此作用作为连续映射的复合是连续的.
\end{proof}

\begin{theorem}\label{c5-thm-3-6}
若 $G_i$ 与 $A_i$ 如上，则对每个 $n\geq 0$，
\[
H^n_c(G,A)\cong\varinjlim_I H^n_c(G_i,A_i).
\]
\end{theorem}

\begin{proof}
可以证明 $\varinjlim$ 在阿贝尔群范畴 $\mathscr{A}$ 中是正合函子（如\autoref{c3-prop-2-12}）. 这意味着
\[
\varinjlim_I H^n_c(G_i,A_i)\cong H^n_c(\varinjlim_I\overline{C}(G_i,A_i))
\]
其中上链复形 $\overline{C}(G_i,A_i)$ 由诱导映射
\[
\overline{\lambda}_{ij}=(\pi_{ij},\lambda_{ij}):\overline{C}(G_i,A_i)\to\overline{C}(G_j,A_j)
\]
\begin{sloppypar}
对每个 $i\leq j$（见\autoref{c5-sub-3-1}）构成一个正向系. 于是，要证明本命题，我们只需找一个从 $\varinjlim_I C^n_c(G_i,A_i)$ 到 $C^n_c(G,A)$ 且与 $d_n$ 交换的同构. 定义
\end{sloppypar}
\[
\begin{aligned}
\varphi_i:C^n_c(G_i,A_i)\ &\to\ C^n_c(G,A)\\
f\ &\mapsto\ \varphi_i(f):=\lambda_i\circ f\circ\pi_i.
\end{aligned}
\]
注意对每个满足 $i\leq j$ 的 $i,j\in I$，有 $\varphi_j\circ\overline{\lambda}_{ij}=\varphi_i$，事实上
\[
\begin{aligned}
\varphi_j\circ\overline{\lambda}_{ij}(f)&=\varphi_j(\lambda_{ij}\circ f\circ\pi_{ij})\\
&=\lambda_j\circ\lambda_{ij}\circ f\circ\pi_{ij}\circ\pi_j\\
&=\lambda_i\circ f\circ\pi_i\\
&=\varphi_i(f),
\end{aligned}
\]
对每个 $f\in C^n_c(G_i,A_i)$. 于是 $\varphi_i$ 诱导出一个态射
\[
\varphi:\varinjlim_I C^n_c(G_i,A_i)\to C^n_c(G,A),
\]
因为 $\{\varphi_i:C^n_c(G_i,A_i)\to C^n_c(G,A)\}_I$ 是相容族（见\autoref{c5-sub-3-2} 中正向极限的定义）.

容易看出 $\varphi$ 与上边缘算子 $d_k$（$k\geq 0$）交换.

\emph{$\varphi$ 是单射：}设 $f\in\varinjlim_I C^n_c(G_i,A_i)$ 使得 $\varphi(f)=0$. 取 $i_0\in I$ 与 $f_{i_0}\in C^n_c(G_i,A_i)$\footnote{原文即如此；应为 $C^n_c(G_{i_0},A_{i_0})$.} 使得 $\overline{\lambda}_{i_0}(f_{i_0})=f$. 现在记
\begin{equation}\label{c5-eq-2}
f_j:=\overline{\lambda}_{i_0 j}(f_{i_0}),
\end{equation}
对所有 $j\leq i_0$，这特别蕴含
\begin{equation}\label{c5-eq-3}
\overline{\lambda}_j(f_j)=f,
\end{equation}
而可推出 $0=\varphi(f)=\lambda_j\circ f_j\circ\pi_j$ 对所有 $j\leq i_0$ 成立\footnote{原文即如此；此处两处 $j\leq i_0$ 应为 $j\geq i_0$——$\overline{\lambda}_{i_0 j}$ 仅对 $i_0\leq j$ 有定义，且下文即取 $j\geq i_0$.}.\ 我们定义
\[
X_j:=\{g_j=(g_{j1},\ldots,g_{jn})\in G_j^n\mid f_j(g_j)\neq 0\}.
\]
然而，我们只需证明存在某个 $j\geq i_0$ 使得 $X_j=\emptyset$（$\Rightarrow f_j=0$），这将推出 $f=0$（因为 \eqref{c5-eq-2} 与 \eqref{c5-eq-3} —— 见\autoref{c5-sub-3-2}）.

显然有
\[
G_j^n\subseteq\bigcup_{a\in A_j}f_j^{-1}(\{a\}),
\]
但由于 $f_j$ 是连续映射，$A_j$ 是离散 $G$-模且 $G_j^n$ 是紧致群（作为紧致的乘积），$f_j$ 只取有限多个值. 于是 $X_j$ 是闭集（作为有限个闭空间的并）从而紧致（作为 $G_j^n$ 的子群）.

考虑 $i_0\leq i\leq j$ 并设 $g_j\in X_j$，则
\[
\begin{aligned}
0\neq f_j(g_j)&=(\overline{\lambda}_{ij}(f_i))(g_j)\\
&=\lambda_{ij}\circ f_i\circ\pi_{ij}(g_j),
\end{aligned}
\]
而我们有 $f_i\circ\pi_{ij}(g_j)\neq 0$. 这意味着 $\pi_{ij}(g_j)\in X_i$. 于是对每个 $i_0\leq i\leq j$ 有 $\pi_{ij}(X_j)\subseteq X_i$，此外 $\{X_i,\pi_{ij}\}_{i\leq i_0}$ 是紧致空间的逆系. 显然对 $g=(g_1\ldots,g_n)\in\varprojlim_{i\leq i_0}X_i\subseteq G^n$，有 $\varphi(f)(g)\neq 0$. 故 $\varprojlim_{i\leq i_0}X_i=\emptyset$. 现在利用 \cite[Proposition~8, \S9, Chapter~1]{bou71}，我们得到存在某个 $i\geq i_0$ 使得 $X_i=\emptyset$.

\emph{$\varphi$ 是满射：}设 $f\in C^n_c(G,A)$. 我们要找一个连续映射 $f_i:G_i^n$，使得对某个 $i\in I$ 有 $f=\varphi_i(f_i)=\lambda_i\circ f_i\circ\pi_i$. 如前所述，$f$ 连续、$A$ 离散且 $G^n$ 紧致蕴含 $f$ 只取有限多个值. 设 $f(G^n)=\{a_1,\ldots,a_q\}\subseteq A$. 于是存在 $i_0\in I$ 使得 $\lambda_{i_0}A_{i_0}\supseteq f(G^n)$.

考虑 $U_1$ 为 $G$ 的一个正规子群，使得 $f$ 在 $G^n$ 中 $U_1^n$ 的陪集上为常数. 由于 $\{\pi_l^{-1}(U_l)\mid U_l$ 是 $G$ 的正规子群$\}$ 是 $G$ 中 $1$ 的一个开邻域基（见第 3 章\autoref{c3-thm-2-5} 与\autoref{c3-cor-2-6}\footnote{原文作 ``see proposition and corollary 2.6 in chapter 3''；按第 3 章的实际编号，此处应指定理 2.5 与推论 2.6.}），存在 $G_i$ 的一个正规子群 $U_i$ 使得 $U_1\supseteq U:=\pi_i^{-1}(U_i)$ 对某个 $i\geq i_0$ 成立. 注意 $i\geq i_0$ 蕴含
\begin{equation}\label{c5-eq-4}
f(G^n)\subseteq\lambda_{i_0}(A_{i_0})=\lambda_i\circ\lambda_{i_0 i}(A_{i_0})\subseteq\lambda_i(A_i).
\end{equation}
于是
\[
f=\overline{f}\circ p,
\]
其中 $p:G^n\to G^n/U^n$ 是典范投影，而 $\overline{f}:G^n/U^n\to A$ 定义为 $f(gU^n)=f(g)$（由 $U$ 的构造知定义良好）. 注意
\begin{equation}\label{c5-eq-5}
\operatorname{im}(\overline{f})\subseteq\operatorname{im}(f).
\end{equation}
反之 $\pi_i$ 诱导出单射 $\pi_i':G^n/U^n\to G_i^n/U_i^n$（单射性来自 $\pi_i(U^n)\subseteq U_i^n$）. 确切地说，对每个 $\tilde{g}\in G^n$，$\pi_i'(\tilde{g}U^n)=p_i\circ\pi_i(\tilde{g})$，这意味着 $\pi_i'\circ p=p_i\circ\pi_i$.

设 $\overline{f}_i:G_i^n/U_i^n\to A_i$ 是任意满足 $\lambda_i\circ\overline{f}_i\circ\pi_i'=\overline{f}$ 的映射，即使得下图交换的映射：
\[
\begin{tikzcd}
G^n/U^n \arrow[r,"\overline{f}"] \arrow[d,"\pi_i'"'] & A \\
G_i^n/U_i^n \arrow[r,dashed,"\overline{f}_i"'] & A_i \arrow[u,"\lambda_i"']
\end{tikzcd}
\]
由 \eqref{c5-eq-4}、\eqref{c5-eq-5} 及逆极限的定义，这样的映射总是存在. 此外 $\overline{f}_i$ 连续，因为 $A$ 离散、$\overline{f}$ 连续且 $\pi_i'$ 是开映射.

现在定义 $f_i=\overline{f}_i\circ p_i$，其中 $p_i:G_i^n\to G_i^n/U_i^n$ 是典范投影. $f_i$ 作为连续映射的复合显然连续，而且
\[
\begin{aligned}
\lambda_i\circ f_i\circ\pi_i&=\lambda_i\circ\overline{f}_i\circ p_i\circ\pi_i\\
&=\lambda_i\circ\overline{f}_i\circ\pi_i'\circ p\\
&=\overline{f}\circ p=f.
\end{aligned}
\]
\end{proof}

\begin{example}\label{c5-ex-3-7}
设 $\mathbb{L}$ 是 $\mathbb{K}$ 的 Galois 扩张. 回忆 $\mathcal{F}$ 是所有满足 $\mathbb{K}\subseteq\mathbb{E}\subseteq\mathbb{L}$ 且 $\mathbb{E}$ 是 $\mathbb{K}$ 的有限 Galois 扩张的中间域 $\mathbb{E}$ 之集. 回忆
\[
\mathrm{Gal}(\mathbb{L},\mathbb{K})=\varprojlim_{\mathcal{F}}\mathrm{Gal}(\mathbb{E},\mathbb{K}).
\]
我们通过在 $\mathcal{F}$ 上赋予包含偏序并考虑平凡包含 $f_{ij}:\mathbb{E}_i^+\hookrightarrow\mathbb{E}_j^+$（当 $\mathbb{E}_i^+\subseteq\mathbb{E}_j^+$ 时）来定义正向系 $(\mathbb{E},f)_{\mathcal{F}}$. 于是我们要证明 $\mathbb{L}^+\cong\varinjlim_{\mathcal{F}}\mathbb{E}^+\cong\coprod_{\mathbb{E}\in\mathcal{F}}\mathbb{E}^+/\!\sim$（见\autoref{c5-sub-3-2}）. 我们定义
\[
\begin{aligned}
\varphi:\coprod_{\mathbb{E}\in\mathcal{F}}\mathbb{E}^+\ &\to\ \mathbb{L}^+\\
\overline{(\alpha,\mathbb{E})}\ &\mapsto\ \alpha,
\end{aligned}
\]
其中 $\overline{(\alpha,\mathbb{E})}$ 表示 $(\alpha,\mathbb{E})$ 关于关系 $\sim$ 的等价类. 验证此映射定义良好、单且满是容易的（利用 $\mathbb{L}=\cup_{\mathcal{F}}\mathbb{E}$）. 于是，利用前面的\autoref{c5-thm-3-6}，我们有
\[
H^n_c(\mathrm{Gal}(\mathbb{L},\mathbb{K}),\mathbb{L}^+)\cong\varinjlim_{\mathcal{F}}H^n_c(\mathrm{Gal}(\mathbb{E},\mathbb{K}),\mathbb{E}^+).
\]
对乘法群 $\mathbb{L}^*$ 而非 $\mathbb{L}^+$，可证明同样的结论.
\end{example}

\section{\texorpdfstring{Hilbert 第 $90$ 定理}{Hilbert 第 90 定理}}

\newterm{c5-hilbert90}{Hilbert 第 $90$ 定理}{Hilbert's Theorem 90}原本是关于数域的循环扩张的定理，但它有许多推广. 该定理得名于它是 David Hilbert 一部著名著作《Zahlbericht》（发表于 $1897$ 年）中的第 $90$ 个定理，尽管该定理有时被归于 Kummer. 后来它于 $1919$ 年由 Andreas Speiser 推广. 不过这一结果也已知出自 Emmy Noether. 我们将讨论这一定理的各种形式. 我们的参考文献是 \cite{lor98}.

在本节中，我们记 $H^n(G,A)$ 以代替 $H^n_c(G,A)$，从而简化记号.

\subsection{原始的 Hilbert 第 \texorpdfstring{$90$}{90} 定理}
\label{c5-sub-4-1}

为陈述 Hilbert 定理，我们必须定义域扩张的范. 相关内容在 \cite[\S5, Chapter~VI]{lang02a} 中有详细讨论.

\begin{definition}\label{c5-def-4-1}
设 $\mathbb{L}$ 是 $\mathbb{K}$ 的有限扩张. 设 $[\mathbb{L}:\mathbb{K}]_s=r$，且当特征为素数 $p>0$ 时 $[\mathbb{L}:\mathbb{K}]_i=p^\mu$，否则为 $1$.

设 $\sigma_1,\ldots\sigma_r$ 是 $\mathbb{L}$ 在 $\mathbb{K}$ 的一个代数闭包 $\overline{\mathbb{K}}$ 中的不同嵌入. 若 $\alpha$ 是 $\mathbb{L}$ 的一个元素，我们从 $\mathbb{E}$\footnote{原文即如此；应为"从 $\mathbb{L}$ 到 $\mathbb{K}$".} 到 $\mathbb{K}$ 定义其\newterm{c5-norm-map}{范}{norm}为
\[
N_{\mathbb{L}|\mathbb{K}}(\alpha)=N_{\mathbb{K}}^{\mathbb{L}}(\alpha)=\prod_{n=1}^{r}\sigma_n(\alpha^{p^\mu})=\Bigl(\prod_{n=1}^{r}\sigma_n(\alpha)\Bigr)^{[\mathbb{L}:\mathbb{K}]_i}.
\]
\end{definition}

我们给出若干性质而不加证明.

\begin{proposition}\label{c5-prop-4-1}
若 $\mathbb{L}$ 在 $\mathbb{K}$ 上可分，则有
\[
N_{\mathbb{K}}^{\mathbb{L}}(\alpha)=\prod_{\sigma}\sigma(\alpha)
\]
其中乘积取遍 $\mathbb{L}$ 在 $\overline{\mathbb{K}}$ 中在 $\mathbb{K}$ 上的不同嵌入.
\end{proposition}

\begin{proposition}\label{c5-prop-4-2}
设 $\mathbb{L}$ 是 $\mathbb{K}$ 的有限扩张. 则范是从 $\mathbb{L}^*$ 到 $\mathbb{K}^*$ 的乘法同态.
\end{proposition}

\begin{theorem}\label{c5-thm-4-3}
设 $\mathbb{L}$ 是 $\mathbb{K}$ 的循环 Galois 扩张，其 Galois 群为 $G=\mathrm{Gal}(\mathbb{L},\mathbb{K})$，且 $[\mathbb{L}:\mathbb{K}]=n$. 设 $\sigma$ 是 $G$ 的一个生成元，即 $G=<\sigma>$. 设 $\beta\in\mathbb{L}$. 范 $N_{\mathbb{K}}^{\mathbb{L}}(\beta)=N(\beta)$ 等于 $1$ 当且仅当存在 $\alpha\in\mathbb{L}^*$ 使得
\[
\beta=\frac{\alpha}{\sigma(\alpha)}.
\]
\end{theorem}

\begin{proof}
先假设存在这样的 $\alpha$. 则 $N(\beta)=\frac{N(\alpha)}{N(\sigma(\alpha))}$. 但由于 $N(.)$ 是对 $G$ 中所有自同构取的乘积，作用 $\sigma$ 只是将这些自同构置换. 于是 $N(\sigma(\alpha))=N(\alpha)$，即 $N(\beta)=1$.

再假设 $N(\beta)=1$. 我们考虑映射
\[
id+\beta\sigma+\beta\sigma(\beta)\sigma^2+\ldots+\beta\sigma(\beta)\ldots\sigma^{n-2}(\beta)\sigma^{n-1}.
\]
各项互不相同，故可应用 Artin 定理（见 \cite[Theorem~4.1, \S4, Chapter~VI]{lang02a}）. 于是上述映射在 $\mathbb{L}$ 上不恒为零，即存在 $\theta\in\mathbb{L}$ 与 $\alpha\in\mathbb{L}^*$ 使得
\[
\alpha=\theta+\beta\sigma(\theta)+\beta\sigma(\beta)\sigma^2(\theta)+\ldots+\beta\sigma(\beta)\ldots\sigma^{n-2}(\beta)\sigma^{n-1}(\theta).
\]
于是
\[
\beta\sigma(\alpha)=\beta\sigma(\theta)+\beta\sigma(\beta)\sigma^2(\theta)+\ldots+\beta\sigma(\beta)\ldots\sigma^{n-1}(\beta)\sigma^n(\theta),
\]
但由于 $1=N(\beta)=\beta\sigma(\beta)\ldots\sigma^{n-1}(\beta)$，且 $\sigma^n(\theta)=\theta$，我们得到
\[
\beta\sigma(\alpha)=\alpha.
\]
这就完成了证明.
\end{proof}

这就是 Hilbert 的"第 $90$ 定理". 不过 Hilbert 原本有一些限制：域 $\mathbb{K}$ 被假定为数域，且扩张的次数为素数.

\subsection{另一种表述}
\label{c5-sub-4-2}

考虑 $\mathbb{L}$ 是域 $\mathbb{K}$ 的 Galois 扩张及其 Galois 群 $G$. 对 $\alpha\in\mathbb{L}$ 与 $\sigma\in G$ 我们使用记号 $\alpha^\sigma=\sigma(\alpha)$，并使用指数记号，即对 $\alpha\in\mathbb{L}$ 与 $\tau,\sigma\in G$ 有 $\alpha^{\tau+\sigma}=\tau(\alpha)\circ\sigma(\alpha)$.

设 $A$ 是 $G$-模 $\mathbb{L}^*$. 我们以 $A^{1-G}$ 记由形如 $\alpha^{1-\sigma}$（或 $\alpha/\sigma(\alpha)$，$\alpha\in A$、$\sigma\in G$）的元素生成的 $A$ 的子群. 我们记 ${}_N A$ 为范映射 $N:\mathbb{L}^*\to\mathbb{L}^*$ 的核. 于是我们有
\[
A^{1-G}\subseteq{}_N A,
\]
因为 $N(\alpha^\sigma)=N(\alpha)$（见范的定义）. 接下来我们可定义商
\[
H^{-1}(G,A)={}_N A/A^{1-G}.
\]
对任意 $\sigma,\tau\in A$ 与每个 $\alpha\in A$，我们有
\[
a^{1-\sigma\tau}=a^{1-\sigma+\sigma-\sigma\tau}=a^{1-\sigma}(a^\sigma)^{1-\tau}\in A^{1-\sigma}A^{1-\tau}.
\]
于是，当 $G$ 是以 $\sigma$ 为生成元的循环群时，我们有
\[
A^{1-G}=A^{1-\sigma}.
\]
从而 Hilbert 第 $90$ 定理可表述为

\begin{theorem}\label{c5-thm-4-4}
设 $\mathbb{L}$ 是 $\mathbb{K}$ 的有限循环扩张，其 Galois 群为 $G$. 则 $H^{-1}(G,\mathbb{K}^*)=1$.
\end{theorem}

我们现在寻找这一表述与前面给出的上同调群之间的相似之处. 回忆 $1$-闭链群 $Z^1(G,\mathbb{K}^*)$ 可明确描述为满足下式的映射 $f:G\to\mathbb{K}^*$ 之集
\[
f(\tau\sigma)=\tau(f(\sigma))f(\tau)=f(\sigma)^\tau f(\tau).
\]
若设 $G$ 是以 $\sigma$ 为生成元的循环群，则 $f$ 可由 $\sigma$ 的像来确定. 事实上，记 $\xi=f(\sigma)$，则
\begin{equation}\label{c5-eq-6}
f(\sigma^i)=f(\sigma)^{\sigma^{i-1}}f(\sigma^{i-1})=\xi^{1+\sigma+\sigma^2+\ldots+\sigma^{i-1}},
\end{equation}
对每个 $i\in\mathbb{N}$ 成立. 此外，$N(\xi)=\xi^{1+\sigma+\ldots\sigma^{n-1}}=f(\sigma^n)=1$. 注意满足 $f(\sigma)=\alpha^{1-\sigma}$ 的态射 $f:G\to A$ 恰是 $B^1(G,A)$ 的元素. 于是，若考虑映射
\[
\varphi:H^1(G,A)\to H^{-1}(G,A):f\mapsto\xi=f(\sigma),
\]
我们得到一个单群同态（群的第一同构定理）.

现在设 $\xi\in A$ 且 $N(\xi)=1$，我们利用\eqref{c5-eq-6} 构造映射 $f:G\to A$. 于是此映射定义良好且属于 $Z^1(G,A)$. 这意味着当 $G$ 为循环群时我们有同构
\[
H^1(G,A)\cong H^{-1}(G,A).
\]
故上述定理等价于

\begin{theorem}\label{c5-thm-4-5}
设 $\mathbb{L}$ 是 $\mathbb{K}$ 的有限循环扩张，其 Galois 群为 $G$. 则 $H^1(G,\mathbb{K}^*)=1$.
\end{theorem}

\subsection{Hilbert 第 \texorpdfstring{$90$}{90} 定理的推广}
\label{c5-sub-4-3}

我们在此给出 Hilbert 第 $90$ 定理最一般的形式之一. 不过，该定理还有其他一些形式，其中一个可在 \cite{sch02} 中找到. 我们不在本项目中介绍这一形式，因为它需要许多新的概念.

\begin{theorem}\label{c5-thm-4-6}
设 $\mathbb{L}$ 是域 $\mathbb{K}$ 的 Galois 扩张，则 $H^1(\mathrm{Gal}(\mathbb{L},\mathbb{K}),\mathbb{L}^*)=0$.
\end{theorem}

\begin{proof}
利用\autoref{c5-thm-3-6}，只需在 $\mathbb{L}$ 是 $\mathbb{K}$ 的有限扩张的情形证明本定理. 此时 $\mathrm{Gal}(\mathbb{L},\mathbb{K})$ 是离散群，故上同调群不存在歧义.

回忆 $H^1(\mathrm{Gal}(\mathbb{L},\mathbb{K}),\mathbb{L}^*)=Z^1(\mathrm{Gal}(\mathbb{L},\mathbb{K}),\mathbb{L}^*)/B^1(\mathrm{Gal}(\mathbb{L},\mathbb{K}),\mathbb{L}^*)$，这意味着我们只需证明
\[
Z^1(\mathrm{Gal}(\mathbb{L},\mathbb{K}),\mathbb{L}^*)\subseteq B^1(\mathrm{Gal}(\mathbb{L},\mathbb{K}),\mathbb{L}^*).
\]
注意对每个 $f\in B^1(\mathrm{Gal}(\mathbb{L},\mathbb{K}),\mathbb{L}^*)$，存在 $\alpha\in\mathbb{L}^*$ 使得 $f=d_0(\alpha)$. 于是 $f$ 是 $\mathrm{Gal}(\mathbb{L},\mathbb{K})$ 与 $\mathbb{L}^*$ 之间的一个态射，使得
\[
f(\sigma)=\sigma(\alpha)\alpha^{-1}
\]
对每个 $\sigma\in\mathrm{Gal}(\mathbb{L},\mathbb{K})$ 成立（注意这与\autoref{c5-sub-2-2} 中的表达式等价，只是采用乘法记号）.

设 $g\in Z^1(\mathrm{Gal}(\mathbb{L},\mathbb{K}),\mathbb{L}^*)$. 这意味着 $d_1(g)=0$，于是对每个 $\sigma,\tau\in\mathrm{Gal}(\mathbb{L},\mathbb{K})$
\[
1_{\mathbb{L}}=d_1(g)(\sigma,\tau)=\sigma g(\tau)\cdot g(\sigma\tau)^{-1}\cdot g(\sigma),
\]
从而
\begin{equation}\label{c5-eq-7}
g(\sigma\tau)=\sigma g(\tau)\cdot g(\sigma).
\end{equation}
对 $\alpha\in\mathbb{L}^*$ 定义
\[
A(\alpha)=\sum_{\sigma\in\mathrm{Gal}(\mathbb{L},\mathbb{K})}g(\sigma)\sigma(\alpha).
\]
由于 $\mathrm{Gal}(\mathbb{L},\mathbb{K})$ 的所有自同构互不相同，它们是线性无关的（Galois 理论的结果）. 于是存在 $\alpha\in\mathbb{L}^*$ 使得 $A(\alpha)\neq 0$，记 $A(\alpha)=\beta$.

对每个 $\sigma\in\mathrm{Gal}(\mathbb{L},\mathbb{K})$ 我们有
\[
\sigma(\beta)=\sum_{\tau\in\mathrm{Gal}(\mathbb{L},\mathbb{K})}\sigma(g(\tau))\sigma\tau(\alpha)
\]
两边乘以 $g(\sigma)\in\mathbb{L}^*$：
\[
\begin{aligned}
g(\sigma)\cdot\sigma(\beta)&=\sum_{\tau\in\mathrm{Gal}(\mathbb{L},\mathbb{K})}g(\sigma)\sigma(g(\tau))\sigma\tau(\alpha)\\
&\overset{?}{=}\sum_{\tau\in\mathrm{Gal}(\mathbb{L},\mathbb{K})}g(\sigma\tau)\sigma\tau(\alpha)
\end{aligned}
\]
现在这最后一个和等于 $\beta$，因为 $\sigma\tau$ 取遍 $\mathrm{Gal}(\mathbb{L},\mathbb{K})$. 将上式改写，我们得到
\[
g(\sigma)=\sigma(\beta^{-1})(\beta^{-1})^{-1}
\]
对每个 $\sigma\in\mathrm{Gal}(\mathbb{L},\mathbb{K})$ 成立. 于是 $g\in B^1(\mathrm{Gal}(\mathbb{L},\mathbb{K}),\mathbb{L}^*)$.
\end{proof}

接下来是一个与上述定理相关的定理.

\begin{theorem}\label{c5-thm-4-7}
设 $\mathbb{L}$ 是域 $\mathbb{K}$ 的 Galois 扩张，则对每个 $n\geq 1$，$H^n(\mathrm{Gal}(\mathbb{L},\mathbb{K}),\mathbb{L}^+)=0$.
\end{theorem}

\begin{proof}
利用\autoref{c5-thm-3-6}，只需在 $\mathbb{L}$ 是 $\mathbb{K}$ 的有限扩张的情形证明本定理. 接下来，我们回忆正规基定理（见 \cite[Theorem~13.1, \S13, Chapter~VI]{lang02a}），即存在 $\theta\in\mathbb{L}$ 使得
\[
\{\sigma(\theta)\mid\sigma\in\mathrm{Gal}(\mathbb{L},\mathbb{K})\}
\]
是 $\mathbb{L}^+$ 作为 $\mathbb{K}$-向量空间的一组基. 这等价于
\[
\mathbb{L}^+=\bigoplus_{\sigma\in G}\mathbb{K}^+\sigma(\theta).
\]
而此直和同构于 $K[G]\cong\operatorname{Hom}_{\mathbb{Z}}(\mathbb{Z}G,\mathbb{K}^+)$，即 $\mathbb{L}^+$ 是一个诱导 $\mathbb{Z}G$-模. 于是由\autoref{c4-prop-5-3}，对一切 $n\geq 0$ 有 $H^n(G,\mathbb{L}^+)=\{0\}$.
\end{proof}


\makebibliography

% 术语索引: \newterm 首次引入的术语自动汇入
\printterminology

\end{document}
